% =====================================================================================================================
%  Predicción de Llegadas de Autobuses con Posiciones en Tiempo Real y Horarios Interpolados
%  Memoria técnica del motor de Bizkaibus+, con el formato de las normas APA (7.ª edición)
%
%  Compilar:  latexmk -pdf memoria.tex    (o: pdflatex memoria && biber memoria && pdflatex memoria && pdflatex memoria)
%  Necesita biblatex con el estilo apa (biblatex-apa) y biber, incluidos en TeX Live completo.
% =====================================================================================================================
\documentclass[12pt,letterpaper,oneside,openany]{book}

% ---------------------------------------------------------------- idioma y tipografía
% APA: una sola fuente en todo el texto (Times New Roman de 12 puntos), sans serif en las figuras y monoespaciada de
% 10 puntos en el código.
\usepackage[T1]{fontenc}
\usepackage[spanish,es-tabla,es-noshorthands,es-nolists]{babel}
\usepackage{mathptmx}
\usepackage[scaled=0.92]{helvet}
\renewcommand{\ttdefault}{lmtt}
\usepackage{microtype}
\usepackage{csquotes}

% ---------------------------------------------------------------- página
% APA: carta, márgenes de 2,54 cm, doble espacio en todo (salvo notas al pie, cuerpo de tablas y código), sin espacio
% extra entre párrafos, alineado a la izquierda sin justificar, sin guiones y con sangría de 1,27 cm.
\usepackage[letterpaper,margin=2.54cm,headheight=16pt,headsep=0.75cm]{geometry}
\usepackage{setspace}
\doublespacing
\usepackage{ragged2e}
\AtBeginDocument{\RaggedRight}
\usepackage{footmisc}
\renewcommand{\footnotelayout}{\RaggedRight}
\setlength{\RaggedRightParindent}{1.27cm}
% Sin guiones, el margen derecho tiene que poder quedar tan irregular como haga falta (como \raggedright): con el
% estiramiento de 2em de ragged2e, una palabra larga dejaba la línea corta y TeX avisaba de un «Underfull \hbox».
\setlength{\RaggedRightRightskip}{0pt plus 1fil}
\setlength{\parindent}{1.27cm}
\setlength{\parskip}{0pt}
\hyphenpenalty=10000
\exhyphenpenalty=10000
\usepackage{etoolbox}

% Encabezado de un trabajo profesional: título corto en mayúsculas a la izquierda y número de página a la derecha, en
% todas las páginas; la portada es la página 1.
\usepackage{fancyhdr}
\newcommand{\titulocorto}{PREDICCIÓN DE LLEGADAS DE AUTOBUSES}
\fancypagestyle{apa}{\fancyhf{}\fancyhead[L]{\titulocorto}\fancyhead[R]{\thepage}\renewcommand{\headrulewidth}{0pt}}
\pagestyle{apa}
\fancypagestyle{plain}{\fancyhf{}\fancyhead[L]{\titulocorto}\fancyhead[R]{\thepage}\renewcommand{\headrulewidth}{0pt}}

% ---------------------------------------------------------------- títulos
% APA: títulos sin números, del mismo tamaño que el texto y en negrita. Nivel 1 centrado; nivel 2 a la izquierda;
% nivel 3 a la izquierda en cursiva; niveles 4 y 5 con sangría, terminados en punto y seguidos del texto.
\usepackage{titlesec}
\titleformat{\chapter}[block]{\normalfont\normalsize\bfseries\filcenter}{}{0pt}{}
\titlespacing*{\chapter}{0pt}{0pt}{0pt}
\titleformat{\section}[block]{\normalfont\normalsize\bfseries\filright}{}{0pt}{}
\titlespacing*{\section}{0pt}{0pt}{0pt}
\titleformat{\subsection}[block]{\normalfont\normalsize\bfseries\itshape\filright}{}{0pt}{}
\titlespacing*{\subsection}{0pt}{0pt}{0pt}
\titleformat{\subsubsection}[runin]{\normalfont\normalsize\bfseries}{}{0pt}{}[.]
\titlespacing*{\subsubsection}{1.27cm}{0pt}{0.5em}
\titleformat{\paragraph}[runin]{\normalfont\normalsize\bfseries\itshape}{}{0pt}{}[.]
\titlespacing*{\paragraph}{1.27cm}{0pt}{0.5em}
\setcounter{secnumdepth}{3}
\setcounter{tocdepth}{2}

% Índice: sin números de capítulo ni de sección; los apéndices, con su letra.
\usepackage{titletoc}
\newif\ifapendices
\newcommand{\puntos}{\titlerule*[0.6pc]{.}\contentspage}
\titlecontents{chapter}[0pt]{}{\ifapendices Apéndice~\thecontentslabel.\ \fi}{}{\puntos}
\titlecontents{section}[1.27cm]{}{}{}{\puntos}
\titlecontents{subsection}[2.54cm]{}{}{}{\puntos}
\titlecontents{figure}[0pt]{\addtolength{\rightskip}{0pt plus 1fil}}{Figura~\thecontentslabel.\ }{}{\puntos}
\titlecontents{table}[0pt]{\addtolength{\rightskip}{0pt plus 1fil}}{Tabla~\thecontentslabel.\ }{}{\puntos}

% ---------------------------------------------------------------- matemáticas
\usepackage{amsmath,amssymb,amsthm,mathtools}
\newtheorem{definicion}{Definición}
\newtheorem{proposicion}{Proposición}
\newtheorem{observacion}{Observación}
\theoremstyle{definition}
\newtheorem{ejemplo}{Ejemplo}
\DeclareMathOperator*{\argmin}{arg\,min}
\DeclareMathOperator{\clamp}{clamp}
\DeclareMathOperator{\redondeo}{rnd}
\DeclareMathOperator{\mediana}{med}
\DeclareMathOperator{\franja}{franja}
\newcommand{\R}{\mathbb{R}}
\newcommand{\N}{\mathbb{N}}
\newcommand{\Z}{\mathbb{Z}}
\newcommand{\tgps}{t_{\mathrm{gps}}}
\newcommand{\tref}{t_{\mathrm{ref}}}
\newcommand{\tnow}{t_{\mathrm{now}}}
\newcommand{\Tpos}{T_{\mathrm{pos}}}
\newcommand{\ETA}{\mathrm{ETA}}
\newcommand{\dest}{\mathrm{dest}}
\counterwithout{equation}{chapter}

% ---------------------------------------------------------------- tablas, figuras, pseudocódigo, código
\usepackage{booktabs,array,tabularx,longtable,multirow}
% Las columnas X de tabularx, sin justificar, como el resto del texto (sin guiones, justificadas quedaban huecos).
\renewcommand{\tabularxcolumn}[1]{>{\raggedright\arraybackslash}p{#1}}
\usepackage{graphicx}
\usepackage{xcolor}
\definecolor{ruta}{HTML}{01573C}
\definecolor{gps}{HTML}{D9480F}
\definecolor{gris}{HTML}{56635A}
\usepackage{tikz}
\usetikzlibrary{arrows.meta,positioning,calc,decorations.pathreplacing,shapes.geometric}
\usepackage{pgfplots}
\pgfplotsset{compat=1.18}
% APA: dentro de las figuras, fuente sans serif de 8 a 14 puntos.
\AtBeginEnvironment{tikzpicture}{\sffamily}
\usepackage[noend]{algpseudocode}
\algrenewcommand\algorithmicrequire{\textbf{Entrada:}}
\algrenewcommand\algorithmicensure{\textbf{Salida:}}
\algrenewcommand\algorithmicif{\textbf{si}}
\algrenewcommand\algorithmicthen{\textbf{entonces}}
\algrenewcommand\algorithmicelse{\textbf{si no}}
\algrenewcommand\algorithmicfor{\textbf{para}}
\algrenewcommand\algorithmicdo{\textbf{hacer}}
\algrenewcommand\algorithmicreturn{\textbf{devolver}}
\algrenewcommand\algorithmicwhile{\textbf{mientras}}
\algrenewcommand\algorithmicfunction{\textbf{función}}
\AtBeginEnvironment{algorithmic}{\setstretch{1.15}\raggedright}
\usepackage{listings}
% Sin guiones en el texto (APA), pero el código sí puede partir sus líneas largas: listings las corta con el mismo
% mecanismo que los guiones.
\makeatletter
\lst@AddToHook{Init}{\hyphenpenalty=50\relax\exhyphenpenalty=50\relax}
\makeatother
\usepackage{embedfile}
\lstset{
  basicstyle=\ttfamily\fontsize{10}{12}\selectfont\setstretch{1},
  columns=fullflexible,
  keepspaces=true,
  breaklines=true,
  frame=single,
  rulecolor=\color{gris!40},
  numbers=left,
  numberstyle=\tiny\color{gris},
  showstringspaces=false,
  keywordstyle=\color{ruta}\bfseries,
  commentstyle=\color{gris}\itshape,
  stringstyle=\color{gps},
  literate={á}{{\'a}}1 {é}{{\'e}}1 {í}{{\'i}}1 {ó}{{\'o}}1 {ú}{{\'u}}1 {ñ}{{\~n}}1 {ü}{{\"u}}1
           {Á}{{\'A}}1 {É}{{\'E}}1 {Í}{{\'I}}1 {Ó}{{\'O}}1 {Ú}{{\'U}}1 {Ñ}{{\~N}}1
           {«}{{\guillemotleft}}1 {»}{{\guillemotright}}1 {·}{{\textperiodcentered}}1 {×}{{$\times$}}1 {±}{{$\pm$}}1
           {-}{{-}}1
}

% APA: tablas y figuras numeradas seguidas en todo el texto, con el número en negrita y el título en cursiva en la línea
% siguiente, encima, alineados a la izquierda y a doble espacio. Las aclaraciones van debajo, en una nota.
\usepackage{caption}
\captionsetup{labelfont=bf,textfont=it,labelsep=newline,justification=raggedright,singlelinecheck=false,
  font=doublespacing,position=top,skip=0pt}
\counterwithout{figure}{chapter}
\counterwithout{table}{chapter}
\newcommand{\nota}[1]{\par\smallskip{\setstretch{2}\RaggedRight\noindent\textit{Nota.} #1\par}}
\AtBeginEnvironment{table}{\setstretch{1.15}}
\AtBeginEnvironment{figure}{\setstretch{1.15}}
\setlength{\LTcapwidth}{\textwidth}

% ---------------------------------------------------------------- referencias y enlaces
% APA 7 con biblatex-apa, en castellano: «y» entre los dos últimos autores, como en la edición en español.
\usepackage[style=apa,backend=biber,natbib=true]{biblatex}
\addbibresource{bibliografia.bib}
\DeclareDelimFormat[bib,biblist]{finalnamedelim}{%
  \ifthenelse{\value{listcount}>\maxprtauth}{}{\addspace y\space}}
\DeclareDelimFormat{finalnamedelim}{\addspace y\space}
\DefineBibliographyStrings{spanish}{retrieved={Recuperado},from={de}}
% Sin coma delante de la «y» (Brakewood, C. y Watkins, K.): la coma tras cada autor se pone solo hasta el penúltimo.
\makeatletter
\renewcommand*{\apablx@ifnotfinalname}{%
  \ifboolexpr{test {\ifnumless{\value{listcount}}{\numexpr\value{liststop}-1\relax}} or test \ifmorenames}}
\makeatother
\setlength{\bibitemsep}{0pt}
\AtBeginBibliography{\RaggedRight}
% Las direcciones se pueden partir, pero sin estirar el espacio entre sus caracteres.
\appto\biburlsetup{\Urlmuskip=0mu\relax}
% APA pide «(Vol. 1, pp. 361–379)» en los capítulos de obras en varios volúmenes; biblatex-apa pone las páginas antes.
\renewbibmacro*{addinfo1}{%
  \printtexte[parens]{%
    \usebibmacro{editor+trans+narr}%
    \setunit*{\addsemicolon\addspace}%
    \printfield{edition}%
    \setunit*{\addcomma\addspace}%
    \ifthenelse{\ifentrytype{report}\OR\ifentrytype{unpublished}\OR\ifentrytype{dataset}\OR
                \ifentrytype{software}\OR\ifentrytype{audio}\OR\ifentrytype{patent}\OR
                \ifentrytype{image}}{\usebibmacro{addinfo3}}{}%
    \setunit*{\addcomma\addspace}%
    \ifthenelse{\ifentrytype{book}\OR\ifentrytype{inbook}}{\printfield{volumes}}{}%
    \setunit*{\addcomma\addspace}%
    \ifthenelse{\ifentrytype{book}\OR\ifentrytype{inbook}}{\printfield{volume}}{}%
    \setunit*{\addcomma\addspace}%
    \ifentrytype{inbook}{\printfield{pages}}{}}}
\usepackage[hidelinks,unicode,pdfusetitle]{hyperref}
\urlstyle{same}
\hypersetup{pdfauthor={Yeray Garrido},
  pdftitle={Predicción de Llegadas de Autobuses con Posiciones en Tiempo Real y Horarios Interpolados}}
\usepackage{bookmark}

% ---------------------------------------------------------------- cortes de línea en nombres largos
% Los nombres del código (CONSTANTES_CON_GUIONES, franjas.con_punto) pueden partirse tras «_» y «.».
\let\guionbajo\_
\renewcommand{\_}{\guionbajo\penalty50\hspace{0pt}}
\setlength{\emergencystretch}{3em}
% Con el texto sin justificar quedan desbordes de décimas de punto que no se ven.
\hfuzz=0.5pt

% ---------------------------------------------------------------- apéndices
% APA: cada apéndice en página nueva, después de las referencias, con «Apéndice A» y su título centrados y en negrita;
% sus tablas y figuras se numeran A1, A2...
\newcommand{\apendices}{%
  \appendix
  \titleformat{\chapter}[display]{\normalfont\normalsize\bfseries\filcenter}{Apéndice~\thechapter}{0pt}{}
  \titlespacing*{\chapter}{0pt}{0pt}{0pt}
  \counterwithin{table}{chapter}\renewcommand{\thetable}{\thechapter\arabic{table}}
  \counterwithin{figure}{chapter}\renewcommand{\thefigure}{\thechapter\arabic{figure}}
  \counterwithin{equation}{chapter}\renewcommand{\theequation}{\thechapter\arabic{equation}}
  \addtocontents{toc}{\protect\apendicestrue}}

% Título de nivel 1 sin número que empieza página y va al índice (resumen, listas, referencias).
\newcommand{\seccionprincipal}[1]{\chapter*{#1}\addcontentsline{toc}{chapter}{#1}}

\begin{document}
\pagenumbering{arabic}

% =====================================================================================================================
% Portada (trabajo profesional): título, autor, afiliación y nota del autor.
\thispagestyle{apa}
\begingroup
  \centering
  \vspace*{3\baselineskip}
  {\bfseries Predicción de Llegadas de Autobuses con Posiciones en Tiempo Real y Horarios Interpolados\par}
  \vspace{\baselineskip}
  Yeray Garrido\par
  Bilbao, Bizkaia\par
  ORCID: \href{https://orcid.org/0009-0008-0459-0294}{0009-0008-0459-0294}\par
  \vfill
  {\bfseries Nota del Autor\par}
  \RaggedRight
  \setlength{\parindent}{1.27cm}
  Este es un trabajo personal e independiente: no se ha presentado como tesis doctoral, no ha sido revisado por
  pares y ninguna universidad, revista ni institución lo ha evaluado. Bizkaibus+ es una aplicación sin relación con
  Bizkaibus, con la Diputación Foral de Bizkaia ni con los operadores del servicio, y usa solo los datos abiertos que se
  publican: el horario en formato GTFS y el feed de posiciones en formato SIRI.\par
  El código del motor es privado. Este documento describe el motor tal como estaba el 7 de noviembre de 2026, con cada
  fórmula, constante y regla transcrita del código y comprobada con vectores de prueba; la implementación de
  referencia, los vectores y el programa que los comprueba van adjuntos al PDF. Es el informe técnico que acompaña al
  artículo en inglés que presenta el método. Ambos se publican, con los datos y el código de referencia, en
  \url{https://github.com/Garridoparrayeray/bide-eta}, donde también pueden dirigirse las consultas sobre este trabajo.\par
  Versión 2.1, 7 de noviembre de 2026. DOI: \href{https://doi.org/10.5281/zenodo.24041572}{10.5281/zenodo.24041572}.
  \copyright~2026 Yeray Garrido. El texto se publica con la licencia Creative
  Commons Atribución 4.0 Internacional (CC~BY~4.0) y el código de referencia, con la licencia MIT.\par
\endgroup
\clearpage

% =====================================================================================================================
% Resumen en castellano, inglés y euskera. APA: cada uno en página nueva, con el título centrado y en negrita, un solo
% párrafo sin sangría de 150 a 250 palabras y, debajo, las palabras clave en un párrafo con sangría.
\seccionprincipal{Resumen}
\noindent Bizkaibus, la red de autobuses interurbanos de Bizkaia, publica la posición GPS de cada vehículo y su
próxima parada, pero no su retraso. Las horas intermedias de su horario oficial están interpoladas a velocidad
constante: de 169\,308 pasos reales por parada, el 85\,\% de los que ocurren pasada la primera hora de viaje llega más
de dos minutos antes de la hora publicada. Esta memoria describe, para que pueda replicarse, el método con que la
aplicación Bizkaibus+ predice las llegadas. El método trabaja con diferencias de horario en lugar
de horas absolutas, aprende el tiempo real de cada tramo por franja horaria con medianas contraídas
y organiza la predicción en tres capas con reglas de abstención: la posición
proyectada sobre el recorrido, la última parada registrada y el horario corregido viaje a viaje. Se recalibra cada día y solo
adopta cambios que reducen el error. En 2,21 millones de predicciones a menos de veinte minutos, de 17 días,
contrastadas con el paso real, la llegada mostrada cae dentro de $\pm2$ minutos en el 96\,\% cuando el autobús está a
menos de tres minutos y en el 67\,\% entre diez y veinte minutos, frente al 24\,\% y el 25\,\% del horario publicado.
Evaluados por días, un filtro de Kalman sobre el ritmo del propio autobús y una red neuronal reducen el error medio de
las sumas de tramos un 4\,\% y entre un 6 y un 7\,\%, respectivamente. Se incluyen vectores de prueba y una
implementación de referencia que los reproduce.

\textit{Palabras clave:} predicción de llegadas, transporte público, posicionamiento GPS, GTFS, SIRI, estimación
robusta, contracción, validación temporal

\seccionprincipal{Abstract}
\noindent Bizkaibus, the interurban bus network of Biscay, publishes the GPS position and next stop of every
vehicle, but not its delay. The intermediate times of its official timetable are interpolated at constant speed:
of 169,308 actual stop passages, 85\% of those occurring after the first hour of travel arrive more than two minutes
before the published time. This report describes the method the Bizkaibus+ app uses to predict arrivals, so that it
can be reimplemented in any language. The method works with timetable differences instead of absolute times, learns the
actual travel time of every segment between stops by time band with medians shrunk towards the all-day median of the
segment, and organizes the prediction into three layers with explicit abstention rules: the position projected on the
route, the last stop registered by the feed and a timetable corrected trip by trip. It recalibrates itself daily and only adopts
changes that reduce the measured error. Over 17 days, in 2.21 million predictions within twenty minutes of arrival,
checked against the actual passage, the displayed arrival is within $\pm2$ minutes 96\% of the time when the bus is under three
minutes away and 67\% at ten to twenty minutes, against 24\% and 25\% for the published timetable. Under day-by-day
evaluation, a Kalman filter on the bus's own pace and a neural network reduce the mean error of sums of segment times by
4\% and by 6 to 7\%, respectively. Test vectors and a reference implementation that reproduces them are included.

\textit{Keywords:} arrival time prediction, public transport, GPS positioning, GTFS, SIRI, robust estimation,
shrinkage, temporal validation

\seccionprincipal{Laburpena}
\noindent Bizkaibusek, Bizkaiko hiriarteko autobus sareak, ibilgailu bakoitzaren GPS kokapena eta hurrengo geralekua
argitaratzen ditu, baina ez atzerapena, eta ordutegi ofizialeko tarteko orduak abiadura konstantean interpolatuta
daude: geralekuetako 169\,308 benetako igaroalditan, bidaiaren lehen orduaren ondoren, igaroaldien \%85 argitaratutako
ordua baino bi minutu baino lehenago igarotzen da. Memoria honek Bizkaibus+ aplikazioak iritsierak aurreikusteko
erabiltzen duen metodoa deskribatzen du, edozein lengoaiatan berriro egiteko behar den xehetasunarekin. Metodoak
ordutegiaren diferentziak erabiltzen ditu, ez ordu absolutuak; geralekuen arteko tarte bakoitzaren benetako denbora
ordu-tarteka ikasten du, tartearen egun osoko medianarantz hurbildutako medianekin; eta aurreikuspena hiru geruzatan
antolatzen du, abstentzio-arau esplizituekin: ibilbidean proiektatutako kokapena, jarioak erregistratutako azken
geralekua eta bidaiaz bidaia zuzendutako ordutegia. Egunero birkalibratzen da jarioaren harrapaketekin, eta neurtutako errorea
jaisten duten aldaketak bakarrik onartzen ditu. Benetako igaroaldiarekin alderatutako 17 eguneko 2,21 milioi
aurreikuspenetan (hogei minutu barrukoak), erakutsitako iritsiera $\pm2$ minuturen barruan dago kasuen \%96an autobusa hiru minutu baino
gutxiagora dagoenean, eta \%67an hamar eta hogei minutu artean; argitaratutako ordutegiak \%24 eta \%25 asmatzen ditu.
Egunez eguneko ebaluazio-protokoloarekin, autobusaren beraren erritmoaren gaineko Kalman iragazki batek eta sare
neuronal batek tarte-denboren baturen batez besteko errorea \%4 eta \%6-7 murrizten dute, hurrenez hurren. Memoriak formulak, konstanteak, sasikodea, proba-bektoreak
eta horiek erreproduzitzen dituen erreferentziazko inplementazio bat ditu.

\textit{Gako-hitzak:} iritsieren aurreikuspena, garraio publikoa, GPS kokapena, GTFS, SIRI, estimazio sendoa,
hurbilketa, denborazko balidazioa


\clearpage
\renewcommand{\contentsname}{Contenido}
\tableofcontents
\clearpage
\phantomsection
\renewcommand{\listfigurename}{Lista de Figuras}
\addcontentsline{toc}{chapter}{Lista de Figuras}
\listoffigures
\clearpage
\phantomsection
\renewcommand{\listtablename}{Lista de Tablas}
\addcontentsline{toc}{chapter}{Lista de Tablas}
\listoftables
\seccionprincipal{Acrónimos}

\begin{description}
  \item[APC] \emph{Automatic Passenger Counting}, conteo automático de viajeros.
  \item[API] \emph{Application Programming Interface}, interfaz de programación de aplicaciones.
  \item[AVL] \emph{Automatic Vehicle Location}, localización automática de vehículos.
  \item[CRC] \emph{Cyclic Redundancy Check}, código de redundancia cíclica.
  \item[ECM] Error cuadrático medio.
  \item[ETA] \emph{Estimated Time of Arrival}, hora estimada de llegada.
  \item[GPS] \emph{Global Positioning System}, sistema de posicionamiento global.
  \item[GTFS] \emph{General Transit Feed Specification}, formato abierto de horarios de transporte público.
  \item[ISO] \emph{International Organization for Standardization}, Organización Internacional de Normalización.
  \item[JSON] \emph{JavaScript Object Notation}, formato de intercambio de datos.
  \item[LSTM] \emph{Long Short-Term Memory}, tipo de red neuronal recurrente.
  \item[MD5] \emph{Message-Digest Algorithm 5}, función resumen de 128 bits.
  \item[PHP] Lenguaje de programación en el que está escrito el motor.
  \item[SIRI] \emph{Service Interface for Real Time Information}, norma europea EN~15531.
  \item[SIRI-VM] Servicio \emph{Vehicle Monitoring} de SIRI.
  \item[UTC] Tiempo universal coordinado.
  \item[WGS84] \emph{World Geodetic System 1984}, sistema de referencia de las coordenadas GPS.
  \item[XML] \emph{Extensible Markup Language}, lenguaje de marcado extensible.
\end{description}


% =====================================================================================================================
\chapter{Introducción}
\label{cap:introduccion}

\section{Motivación}
\label{sec:motivacion}

Quien espera un autobús en una parada no necesita saber a qué hora \emph{debería} pasar, sino cuánto falta para que
pase \emph{de verdad}. La diferencia entre ambas cosas es el objeto de esta memoria. En una red interurbana como
Bizkaibus, con líneas que combinan calles congestionadas, carreteras y autopistas, esa diferencia puede ser de muchos
minutos y cambia con la hora del día, el día de la semana y el tramo concreto del recorrido.

\subsection{Esperar un Autobús Interurbano}
\label{sec:esperar}

En una línea urbana que pasa cada cinco minutos, el viajero no necesita saber a qué hora llega el autobús: va a la
parada y espera, en promedio, la mitad del intervalo. En una red interurbana eso no funciona, porque los intervalos son
largos. En el horario de Bizkaibus de un miércoles laborable (28 de octubre de 2026), entre las 07:00 y las 21:00, la
mitad de las combinaciones de línea y parada tienen, como mucho, un autobús por hora: el intervalo mediano es de
60~minutos, el 68\,\% de las combinaciones tiene un intervalo mediano de una hora o más, el 87\,\% de media hora o más,
y solo el 3\,\% tiene un autobús cada cuarto de hora o con más frecuencia.\footnote{Calculado sobre las 4\,932
combinaciones de línea y parada con servicio ese día, tomando para cada una el intervalo mediano entre salidas
consecutivas de la misma línea en la misma parada.}

Con esos intervalos, el viajero planifica: decide a qué hora salir de casa a partir de la hora a la que cree que
pasará el autobús. Si llega tarde, aunque sea por un minuto, pierde el autobús y espera el siguiente, a menudo una
hora. Si llega con mucho margen, espera en una parada que con frecuencia está en el arcén de una carretera, sin
marquesina ni información. Cuanto menos fiable es la hora de que dispone, más margen deja y más tiempo pierde.

\subsection{Lo que Cuesta No Saber Cuándo Llega}
\label{sec:coste-incertidumbre}

La literatura sobre transporte público ha medido el efecto de esta incertidumbre. Los viajeros que esperan sin
información perciben su espera como más larga de lo que es \citep{mishalani2006}. En Seattle, los usuarios de una
aplicación con llegadas en tiempo real no percibían su espera como más larga que la medida, a diferencia de quienes no
la usaban, y además esperaban menos \citep{watkins2011}. Una revisión de la literatura sobre los beneficios de la
información en tiempo real recoge, entre otros, menos tiempo de espera, menos tiempo total de viaje, más satisfacción,
más sensación de seguridad y, en algunos estudios, más uso del transporte público \citep{brakewood2019}.

Estos beneficios dependen de que la información sea \emph{precisa}. Una hora equivocada no solo no ayuda: hace perder
el autobús o esperar más que sin ella.

\subsection{El Horario Publicado Se Desvía en los Dos Sentidos}
\label{sec:desvio-horario}

Las horas de paso por las paradas intermedias que publica el horario de Bizkaibus son, como se explica en la
sección «\nameref{sec:problema}», una interpolación a velocidad constante entre cabeceras. Son las que muestra cualquier
aplicación que use el horario oficial en formato GTFS sin otra corrección. Para saber cuánto se desvían de la realidad,
se compararon con el paso real de los autobuses por las paradas, reconstruido a partir de su GPS
(sección «\nameref{sec:paso-real}»): 169\,308 pasos por parada observados entre el 6 y el 31 de octubre de 2026,
frente al horario vigente el 30 de octubre.\footnote{Se descartan 521 pasos (el 0,3\,\%) con más de 45~minutos de
diferencia, que suelen indicar un cambio de horario entre la fecha de la observación y la del horario usado. El cálculo
está en \texttt{referencia/desvio-horario.py}.} La Tabla~\ref{tab:desvio-horario} y la Figura~\ref{fig:desvio-horario}
muestran el resultado según el tiempo transcurrido desde la salida de cabecera.

\begin{table}[ht]
\centering
\caption{Desvío del Paso Real Respecto a la Hora Publicada, Según el Tiempo Desde la Salida}\label{tab:desvio-horario}
\begin{tabular}{@{}lrrrrr@{}}
\toprule
\shortstack[l]{Desde la\\cabecera} & Pasos & \shortstack{Desvío\\mediano} & \shortstack{Más de 2 min\\antes} & \shortstack{Más de 5 min\\después} & \shortstack{Dentro de\\$\pm2$ min}\\
\midrule
0--10 min   & 29\,115 & $+1{,}3$ min  & 5,0\,\%  & 11,8\,\% & 57,0\,\%\\
10--20 min  & 27\,408 & $-0{,}2$ min  & 33,3\,\% & 15,0\,\% & 34,9\,\%\\
20--30 min  & 28\,449 & $-2{,}2$ min  & 51,5\,\% & 14,4\,\% & 23,6\,\%\\
30--45 min  & 35\,560 & $-5{,}6$ min  & 67,6\,\% & 10,5\,\% & 14,6\,\%\\
45--60 min  & 25\,931 & $-8{,}6$ min  & 74,8\,\% & 7,4\,\%  & 11,6\,\%\\
60 min o más & 22\,845 & $-13{,}6$ min & 85,4\,\% & 3,7\,\%  & 8,0\,\%\\
\midrule
Todos       & 169\,308 & $-2{,}4$ min & 52,1\,\% & 10,7\,\% & 25,3\,\%\\
\bottomrule
\end{tabular}
\nota{El tiempo se cuenta en minutos programados desde la salida de cabecera. Un desvío negativo significa que el autobús pasa \emph{antes} de la hora publicada.}
\end{table}

\begin{figure}[ht]
\centering
\caption{Desvío del Paso Real Según el Tiempo Programado Desde la Salida}\label{fig:desvio-horario}
\begin{tikzpicture}
\begin{axis}[
  width=0.92\textwidth, height=6cm,
  ybar, bar width=9pt,
  symbolic x coords={0--10,10--20,20--30,30--45,45--60,60+},
  xtick=data,
  xlabel={Minutos programados desde la salida de cabecera},
  ylabel={Pasos por parada (\%)},
  ymin=0, ymax=100,
  axis lines*=left, ymajorgrids, grid style={gris!20},
  tick label style={font=\footnotesize}, label style={font=\small},
  legend style={font=\footnotesize,draw=none,fill=none,at={(0.5,-0.28)},anchor=north,legend columns=2,/tikz/every even column/.append style={column sep=1em}},
  nodes near coords, every node near coord/.append style={font=\scriptsize,color=black},
  nodes near coords style={/pgf/number format/fixed,/pgf/number format/precision=0}]
\addplot[fill=gps,draw=gps] coordinates {(0--10,5.0) (10--20,33.3) (20--30,51.5) (30--45,67.6) (45--60,74.8) (60+,85.4)};
\addlegendentry{Pasa más de 2 min antes de la hora publicada}
\addplot[fill=ruta,draw=ruta] coordinates {(0--10,11.8) (10--20,15.0) (20--30,14.4) (30--45,10.5) (45--60,7.4) (60+,3.7)};
\addlegendentry{Pasa más de 5 min después}
\end{axis}
\end{tikzpicture}
\nota{Cerca de la cabecera el autobús va algo retrasado; a medida que avanza, se adelanta cada vez más a la hora publicada, porque el horario reparte el tiempo a velocidad constante y el autobús recorre los tramos de carretera más deprisa.}
\end{figure}

El resultado es el que predice la interpolación. Cerca de la cabecera el autobús va algo retrasado (desvío mediano de
$+1{,}3$~min en los diez primeros minutos). A partir de ahí se adelanta cada vez más al horario: a media hora de la
salida, la mitad de los pasos ocurre más de dos minutos antes de la hora publicada, y pasada la hora, el 85\,\%, con un
adelanto mediano de casi catorce minutos. En fin de semana el adelanto es mayor: se adelantan más de dos minutos el
60\,\% de los pasos, frente al 50\,\% en laborables. De las 69 líneas con al menos 500 pasos observados, en 39 la
mayoría de los pasos se adelanta más de dos minutos. Solo una cuarta parte de los pasos cae dentro de $\pm2$~minutos de la hora publicada.

Para el viajero, eso significa dos riesgos de signo contrario:

\begin{itemize}
  \item \textbf{Perder el autobús.} Quien se guiara por la hora publicada y llegara a la parada dos minutos antes
        encontraría que el autobús ya había pasado en el 52\,\% de los casos. Con un intervalo típico de una hora, cada
        autobús perdido es una hora de espera.
  \item \textbf{Esperar de más.} Cuando el autobús todavía no ha pasado, esa misma persona espera una mediana de
        3,8~minutos, pero una de cada diez veces, más de diez minutos, y el 10,7\,\% de los pasos llega más de cinco
        minutos tarde.
\end{itemize}

\noindent Estas cifras no son un juicio sobre la puntualidad del servicio: las horas intermedias del GTFS no son un
compromiso del operador (sección «\nameref{sec:problema}»), sino una aproximación. Son, sin embargo, las horas que ve quien
consulta el horario en una aplicación que no las corrige, y por eso el problema merece una solución.

\subsection{Este Trabajo}

Bizkaibus+ es una aplicación web independiente que muestra, para cada parada, las próximas salidas con su hora
estimada. Este documento describe el algoritmo que calcula esas horas: de dónde salen los datos, qué se hace con
ellos, por qué se hace así y cómo se ha comprobado que funciona. El objetivo es que el método pueda estudiarse,
criticarse y replicarse en cualquier otra ciudad y en cualquier lenguaje de programación.

\section{El Problema}
\label{sec:problema}

Los datos disponibles condicionan todo el diseño. Bizkaibus publica dos fuentes abiertas, que son las únicas que usa
este trabajo:

\begin{enumerate}
  \item Un \textbf{horario programado} en formato GTFS \citep{gtfs}, publicado en el Punto de Acceso Nacional de
        transporte multimodal \citep{nap}: líneas, viajes, paradas y la hora de paso de cada viaje por cada parada,
        junto con el trazado geográfico (\emph{shape}) de cada recorrido.
  \item Un \textbf{feed en tiempo real} en formato SIRI Vehicle Monitoring \citep{siri}, actualizado aproximadamente
        una vez por minuto, que para cada autobús en servicio indica el viaje que realiza, su posición GPS, la hora a la que se tomó
        esa posición y el número de orden de su próxima parada.
\end{enumerate}

Dos características de estas fuentes descartan el enfoque habitual de sumar un retraso al horario:

\begin{itemize}
  \item \textbf{El feed no publica retrasos.} El elemento \texttt{Delay} de SIRI no aparece, y el feed alternativo de
        actualizaciones de viaje declara siempre un retraso nulo (\texttt{PT0S}) con una marca de tiempo inválida.
  \item \textbf{El horario intermedio es una interpolación.} Las paradas intermedias del GTFS llevan
        \texttt{timepoint=0}: sus horas se obtienen repartiendo el tiempo entre cabeceras a velocidad constante. Un
        autobús que recorre 18~km de autopista en 15~minutos puede tener, para ese tramo, entre 34 y 49~minutos
        programados. En mitad del recorrido, el horario absoluto puede ir quince o veinte minutos por detrás del
        autobús sin que este se haya adelantado en ningún sentido operativo.
\end{itemize}

Así, ni existe un retraso que sumar ni el horario absoluto es una referencia fiable. Lo único vivo y fiable es la
posición.

\section{Objetivos}

\begin{enumerate}
  \item Estimar, para cada viaje en servicio y cada parada que le quede por recorrer, la hora de llegada con el menor
        error posible, a partir de la posición del vehículo.
  \item Seguir dando una estimación útil cuando el GPS falla o no existe, sin inventar datos.
  \item Aprender de forma automática el comportamiento real de la red, sin reglas escritas a mano sobre la velocidad
        de los autobuses.
  \item Medir la precisión contra el paso real de los autobuses por las paradas, y adoptar un cambio en el modelo solo
        si esa medida mejora.
  \item Documentarlo todo con la precisión suficiente para que otra persona pueda reconstruirlo sin leer el código.
\end{enumerate}

\section{Contribuciones}

\begin{itemize}
  \item Un método de predicción por \textbf{diferencias de horario} que neutraliza el error de forma de un horario
        interpolado (capítulo «\nameref{cap:modelo}», proposición~\ref{prop:diferencias}).
  \item Un procedimiento de \textbf{localización} que sitúa el GPS sobre el trazado oficial y es robusto ante los
        recorridos que vuelven sobre sí mismos (capítulo «\nameref{cap:localizacion}»).
  \item Un \textbf{modelo de tiempos por tramo} con estimadores robustos, contracción hacia la mediana del tramo y una
        cascada de respaldo para los tramos sin datos (capítulo «\nameref{cap:modelo}»).
  \item Una arquitectura de \textbf{tres capas con abstención}, que elige en cada predicción la mejor fuente fiable y
        nunca afirma un retraso que no pueda sostener, con una \textbf{corrección por viaje} para cuando no hay ningún
        dato del vehículo (capítulo «\nameref{cap:prediccion}»).
  \item Una \textbf{revisión de los datos de entrada} en la construcción del horario: los trazados dibujados al revés se
        invierten y, en las redes que no publican los metros de cada parada, estos se miden sobre el trazado
        (capítulo «\nameref{cap:datos}»).
  \item Un ciclo de \textbf{aprendizaje continuo} con validación cruzada y reglas de escritura conservadoras
        (capítulo «\nameref{cap:aprendizaje}»), un protocolo de \textbf{evaluación} por días que reproduce lo que ve el
        usuario, y una \textbf{comparación} con filtros de Kalman y redes neuronales hecha con ese mismo protocolo
        (capítulo «\nameref{cap:evaluacion}»).
  \item \textbf{Vectores de prueba} generados con el código de producción y una \textbf{implementación de referencia}
        independiente que los reproduce exactamente (Apéndices~\ref{ap:vectores} y~\ref{ap:referencia}).
\end{itemize}

\section{Visión General del Sistema}

La Figura~\ref{fig:tuberia} resume el recorrido de un dato desde el feed hasta la pantalla. Cada paso se describe en
el capítulo correspondiente.

\begin{figure}[ht]
\centering
\caption{Recorrido de un Dato Desde el Feed Hasta la Pantalla}\label{fig:tuberia}
\begin{tikzpicture}[
  caja/.style={draw=ruta,rounded corners=3pt,thick,align=center,minimum height=1.25cm,text width=2.25cm,font=\footnotesize},
  flecha/.style={-{Stealth[length=2.2mm]},thick,color=ruta}]
  \node[caja] (feed) {\textbf{Feed SIRI-VM}\\posición, próxima parada, hora};
  \node[caja,right=0.45cm of feed] (viaje) {\textbf{Emparejar}\\con el viaje del GTFS};
  \node[caja,right=0.45cm of viaje] (ruta) {\textbf{Proyectar}\\el GPS en el trazado ($\Tpos$)};
  \node[caja,right=0.45cm of ruta] (tramos) {\textbf{Sumar tramos}\\tiempo aprendido hasta la parada};
  \node[caja,below=0.9cm of tramos] (eta) {\textbf{Llegada}\\estado y retraso};
  \node[caja,below=0.9cm of feed] (captura) {\textbf{Captura}\\continua, por horas};
  \node[caja,right=0.45cm of captura] (muestras) {\textbf{Muestras}\\tramos y ritmo};
  \node[caja,right=0.45cm of muestras] (calibrar) {\textbf{Recalibrar}\\cada día, si mejora};
  \draw[flecha] (feed) -- (viaje);
  \draw[flecha] (viaje) -- (ruta);
  \draw[flecha] (ruta) -- (tramos);
  \draw[flecha] (tramos) -- (eta);
  \draw[flecha] (captura) -- (muestras);
  \draw[flecha] (muestras) -- (calibrar);
  \draw[flecha,dashed] (calibrar.north east) -- (tramos.south west);
\end{tikzpicture}
\nota{Arriba, la predicción en cada petición (capítulos «\nameref{cap:localizacion}» a «\nameref{cap:prediccion}»). Abajo, el aprendizaje fuera de línea que alimenta los tiempos por tramo (capítulo «\nameref{cap:aprendizaje}»).}
\end{figure}

\section{Convenciones}
\label{sec:convenciones}

\begin{itemize}
  \item \textbf{Tiempo.} Todas las horas son números enteros de segundos desde la medianoche de la zona
        \texttt{Europe/Madrid} del día de servicio. Siguiendo a GTFS, pueden superar $86\,400$: las 25:10 son
        $90\,600$~s. El feed y el reloj dan horas del día en curso; para un viaje del día anterior se les suman
        $86\,400$~s (sección «\nameref{sec:mas-de-24h}»). Las duraciones se expresan en segundos salvo que se diga otra
        cosa.
  \item \textbf{Distancia.} En metros. Las coordenadas geográficas, en grados decimales (WGS84).
  \item \textbf{Redondeo.} $\redondeo(x)$ redondea al entero más próximo con las mitades \emph{alejándose del cero}
        ($\redondeo(2{,}5)=3$, $\redondeo(-2{,}5)=-3$). Es el comportamiento de PHP y difiere del «redondeo del
        banquero» de otros lenguajes; replicarlo es necesario para obtener los mismos resultados.
  \item \textbf{Mediana.} Salvo que se indique otra cosa, $\mediana$ es la \emph{mediana baja}: con los valores
        ordenados $x_{(0)}\le\dots\le x_{(n-1)}$, $\mediana(x)=x_{(\lfloor (n-1)/2\rfloor)}$. En general, el percentil
        $p$ es $P_p(x)=x_{(\lfloor p\,(n-1)\rfloor)}$, sin interpolar.
  \item \textbf{Acotación.} $\clamp(x,a,b)=\max(a,\min(b,x))$.
  \item \textbf{Nombres en el código.} Cuando conviene, se indica entre paréntesis y en letra \texttt{monoespaciada} el
        nombre de la constante o función del código, para facilitar la correspondencia.
\end{itemize}

\section{Estructura del Documento}

El documento sigue la estructura habitual de un trabajo de investigación (introducción, método, resultados y
discusión) y el formato de las normas APA en su séptima edición. Este capítulo y el capítulo «\nameref{cap:marco}»
sitúan el trabajo respecto a la literatura de predicción de llegadas y a las herramientas estadísticas que usa. Los
cinco siguientes describen el método completo: las fuentes de datos (capítulo «\nameref{cap:datos}»), la localización del
vehículo (capítulo «\nameref{cap:localizacion}»), el modelo del tiempo de recorrido (capítulo «\nameref{cap:modelo}»), la
predicción en tres capas (capítulo «\nameref{cap:prediccion}») y el aprendizaje de los parámetros
(capítulo «\nameref{cap:aprendizaje}»). El capítulo «\nameref{cap:evaluacion}» presenta la precisión medida en una muestra
con las tres capas, a lo largo de 17 días y frente a otros métodos; el capítulo «\nameref{cap:discusion}» la interpreta y
discute las limitaciones del método y su aplicación a otras redes, y el capítulo «\nameref{cap:conclusiones}» cierra con las
conclusiones y las líneas futuras. Tras las referencias, los apéndices reúnen las constantes, el pseudocódigo
completo, los formatos de datos, los vectores de prueba, la implementación de referencia y un glosario.

\chapter{Marco Teórico y Estado del Arte}
\label{cap:marco}

Este capítulo sitúa el método en el contexto de la predicción de llegadas de transporte público y presenta las
herramientas estadísticas en que se apoya. Se centra en los trabajos y las técnicas que explican de dónde vienen las
ideas del método y por qué se eligieron unas y no otras.

\section{Predicción de la Hora de Llegada}

Predecir cuándo llegará un vehículo a un punto de su recorrido es un problema clásico de los sistemas inteligentes de
transporte. Los enfoques de la literatura pueden agruparse en cuatro familias.

\subsection{Medias Históricas} Se estima el tiempo de recorrido entre dos puntos con el valor típico observado en el
pasado para esa hora y ese día. Son sencillas, estables y fáciles de explicar, pero no reaccionan a lo que le ocurre
al vehículo concreto.

\subsection{Regresión y Modelos de Series Temporales} Relacionan el tiempo de recorrido con variables explicativas
(hora, distancia, número de paradas, demanda) o con su propia historia reciente. Exigen elegir y mantener esas
variables.

\subsection{Filtros de Kalman} Combinan una predicción \emph{a priori} con la observación más reciente, ponderando cada una
por su incertidumbre \citep{kalman1960}. Se han usado para predecir llegadas y salidas a partir de datos de
localización automática de vehículos (AVL) y de conteo automático de viajeros (APC) \citep{shalaby2004}.

\subsection{Redes Neuronales} Aprenden la relación entre la posición y el tiempo restante a partir de muchos datos
\citep{chien2002,jeong2005}. Suelen ser precisas cuando hay abundancia de datos homogéneos, pero son difíciles de
auditar y de explicar a quien usa el sistema.

\subsection{Trabajos Recientes} Dos revisiones recientes recogen el estado del campo: \citet{kumar2025review}, que cubre
el proceso completo (datos, modelos, discretización del recorrido y medidas de evaluación), y
\citet{singh2022review}, centrada en los modelos de inteligencia artificial. Partir el viaje en tramos entre paradas
consecutivas y predecir cada uno por separado es habitual \citep{gal2017journey}. La variabilidad del tiempo de
recorrido depende mucho de la hora del día \citep{mazloumi2010variability}, lo que justifica aprender por franjas.
\citet{lin1999experimental} compararon algoritmos que combinan el GPS con el horario en un servicio rural de baja
frecuencia, un caso parecido al de los autobuses interurbanos. Entre los filtros de Kalman, \citet{cathey2003prescription} propusieron el esquema de
seguimiento, filtrado y predicción para datos AVL, y \citet{vanajakshi2009kalman} lo aplicaron a autobuses en tráfico
heterogéneo. Entre las redes neuronales se han usado redes recurrentes \citep{pang2019rnn}, redes LSTM convolucionales
que predicen a la vez todos los enlaces de una línea \citep{petersen2019convlstm}, redes de grafos con atención para
datos escasos \citep{ma2022multiattention} y transformadores entrenados con datos abiertos
\citep{jeong2024battransformer}. En la predicción del tráfico, el aprendizaje profundo no siempre es la mejor técnica,
y los métodos de referencia sencillos y bien ajustados siguen siendo importantes \citep{manibardo2022deep}.

Otros trabajos parten, como este, de los datos que publican los operadores. \citet{wessel2017retrospective}
reconstruyeron el horario realmente cumplido a partir de un feed público de posiciones, igual que aquí se reconstruye
la verdad de referencia. \citet{cats2016realtime} evaluaron las predicciones que ya se mostraban a los viajeros en
Estocolmo y encontraron que subestimaban la espera. \citet{garciamaurino2024improving} corrigieron los errores
sistemáticos de las predicciones de los operadores de Madrid y París solo con datos de sus API públicas.
\citet{woodard2017predicting} modelaron, para viajes en coche, un efecto propio de cada viaje y recurrieron a
parámetros agrupados cuando un enlace tenía pocas observaciones, dos recursos análogos a la corrección por viaje y a la
contracción de este método.

El método de esta memoria pertenece a la primera familia en su forma de estimar tiempos (valores típicos por tramo y
franja) y combina esa estimación con la observación en tiempo real de la posición, que hace de «reinicio» de la predicción en cada
actualización del feed. La elección es deliberada: con un dato de posición por minuto y sin retrasos publicados, un
modelo explicable, robusto y medible tramo a tramo resultó más útil que uno de caja negra. Las ideas que se
probaron y no mejoraron la medida (entre ellas, corregir con la velocidad del último minuto del vehículo, la forma más
simple de filtrar) se recogen en la sección «\nameref{sec:descartadas}». Un filtro de Kalman sobre la desviación del
propio autobús en los tramos ya terminados, suavizada y olvidada con el tiempo, sí reduce un 4\,\% el error medio de las
sumas de tramos (sección «\nameref{sec:comparacion}»): lo que decide es qué se filtra y cuánto se recuerda
(sección «\nameref{sec:otros-enfoques}»).

\section{Datos de Transporte: GTFS y SIRI}

GTFS \citep{gtfs} es el formato \emph{de facto} para publicar horarios programados. Describe agencias, líneas
(\texttt{routes}), viajes (\texttt{trips}), paradas (\texttt{stops}), horas de paso (\texttt{stop\_times}),
calendarios (\texttt{calendar}, \texttt{calendar\_dates}) y trazados (\texttt{shapes}). Dos campos son centrales
aquí: \texttt{shape\_dist\_traveled}, la distancia recorrida a lo largo del trazado hasta cada parada y cada punto del
trazado, y \texttt{timepoint}, que indica si la hora de paso es exacta (\texttt{1}) o aproximada por interpolación
(\texttt{0}).

GTFS Realtime \citep{gtfsrt} y SIRI \citep{siri} son los formatos más usados para datos en tiempo real. SIRI Vehicle
Monitoring (SIRI-VM) publica, para cada vehículo, la actividad de su viaje: referencia del viaje
(\texttt{VehicleJourneyRef}), posición (\texttt{VehicleLocation}), hora de la posición
(\texttt{LocationRecordedAtTime}) y llamada a la próxima parada (\texttt{MonitoredCall}), con su orden
(\texttt{Order} o \texttt{VisitNumber}).

\section{Ajuste al Mapa}

Situar una posición GPS sobre una red o un recorrido se conoce como \emph{map matching} \citep{quddus2007}. Los
métodos van desde la proyección geométrica sobre el segmento más próximo hasta modelos probabilísticos que tienen en
cuenta la topología y la trayectoria. Aquí el recorrido se conoce de antemano (es el trazado del viaje
que el propio feed declara), lo que reduce el problema a una proyección sobre una polilínea. La dificultad que queda
es la ambigüedad cuando el recorrido pasa dos veces por la misma calle; el capítulo «\nameref{cap:localizacion}» la resuelve
con una ventana alrededor de la próxima parada y una regla de desempate.

Para distancias de unos pocos kilómetros basta una proyección equirrectangular local \citep{snyder1987}, que trata la
superficie terrestre como un plano alrededor del punto de interés.

\section{Estadística Robusta}

La media de una muestra es muy sensible a valores extremos: un único autobús parado veinte minutos por un accidente
desplaza la media de un tramo tanto como muchos autobuses normales. Los estimadores \emph{robustos}
\citep{huber1964,tukey1977} limitan esa influencia. La mediana es el más sencillo: su punto de ruptura es del 50\,\%,
es decir, la mitad de la muestra tendría que estar contaminada para llevarla a un valor arbitrario. El método usa
medianas en todos sus agregados, y además toma un solo valor por vehículo y viaje, para que un vehículo lento visto
muchas veces no pese más que uno visto una vez.

\section{Contracción y Modelos Jerárquicos}
\label{sec:marco-contraccion}

Cuando se estiman muchas medias a la vez (un tiempo para cada tramo y cada franja horaria), algunas se apoyan en muy
pocas observaciones y son ruidosas. \citet{jamesstein1961} demostraron que, al estimar simultáneamente tres o más medias de
distribuciones normales, un estimador que contrae las medias observadas hacia un punto fijo tiene un error cuadrático
total esperado menor que el de usar cada media por separado. \citet{efronmorris1975} mostraron su utilidad con datos
reales, contrayendo cada media hacia la media común del grupo. En términos bayesianos
\citep{gelman2013}, la contracción equivale a un modelo jerárquico con \emph{agrupamiento parcial}: cada grupo tiene su
propio parámetro, pero estos se parecen entre sí, y la estimación de cada grupo es una media ponderada de su dato
propio y del valor común, en la que el peso del dato propio crece con el número de observaciones del grupo.

El método aplica esta idea en dos sitios: el tiempo de un tramo en una franja se contrae hacia el tiempo del tramo a
cualquier hora (sección «\nameref{sec:contraccion-tramo}»), y el factor de ritmo se contrae hacia~1 (sección «\nameref{sec:k}»).

\section{Validación}
\label{sec:marco-validacion}

Medir el error con los mismos datos con que se aprendió sobreestima la precisión. La \emph{validación cruzada}
\citep{kohavi1995,hastie2009} separa los datos en particiones, aprende con unas y mide con otras. Cuando los datos
tienen orden temporal, lo honesto es aprender con el pasado y medir con el futuro (\emph{validación con origen
móvil}, \citealp{hyndman2021}). El método usa ambas: una partición por recorridos (un autobús en un viaje de un día) para
decidir si se escribe una nueva calibración y una partición por días para comparar reglas. La primera comparte los
días entre aprendizaje y prueba, y con ellos el tráfico y la meteorología de cada día, por lo que las reglas nuevas se
deciden con la segunda. \citet{hyndmankoehler2006} distinguen entre las medidas de
error que dependen de la escala de los datos, como el error absoluto medio, y las que no: las primeras sirven para
comparar métodos sobre unos mismos datos, pero no entre conjuntos de datos con escalas distintas. Aquí todas las
predicciones se miden en la misma unidad (segundos), de modo que se usan errores absolutos y porcentajes de aciertos
dentro de umbrales que tienen sentido para quien espera el autobús.

\chapter{Fuentes y Modelo de Datos}
\label{cap:datos}

\section{El Horario Programado}

El horario oficial se publica en GTFS y se transforma, una vez al día, en una base de datos SQLite de solo lectura
por red. La Tabla~\ref{tab:esquema} resume las tablas que usa el motor; el esquema completo está en el
Apéndice~\ref{ap:formatos}.

\begin{table}[ht]
\centering
\caption{Tablas de la Base de Datos del Horario que Usa el Motor}\label{tab:esquema}
\begin{tabularx}{\textwidth}{@{}lX@{}}
\toprule
Tabla & Contenido relevante para el motor \\
\midrule
\texttt{stops} & Paradas: identificador, nombre, latitud y longitud. \\
\texttt{service\_journeys} & Viajes: línea, recorrido (\texttt{journey\_pattern\_id}), número de viaje, calendario,
  hora de salida de cabecera y trazado (\texttt{shape\_id}). \\
\texttt{passing\_times} & Para cada viaje y cada parada en orden (\texttt{seq\_order}, desde 1): llegada y salida
  programadas en segundos y \texttt{dist\_m}, la distancia a lo largo del trazado. \\
\texttt{shapes} & Trazado de cada recorrido como lista JSON de puntos $[\varphi,\lambda,\sigma]$: latitud, longitud y
  distancia recorrida. \\
\texttt{trip\_aliases} & Correspondencia entre cada \texttt{trip\_id} original del GTFS y el viaje fusionado al que
  pertenece. \\
\texttt{journey\_patterns} & Recorridos: secuencias de paradas compartidas por varios viajes. \\
\bottomrule
\end{tabularx}
\end{table}

La columna \texttt{dist\_m} de \texttt{passing\_times} es el campo
\texttt{shape\_dist\_traveled} de \texttt{stop\_times.txt}, y la tercera componente de cada punto del trazado es el
\texttt{shape\_dist\_traveled} de \texttt{shapes.txt}. Ambas se miden sobre el mismo trazado y en la misma unidad (o
se corrigen al construir la base hasta que lo hacen; sección «\nameref{sec:trazados-gtfs}»), lo que permite pasar de un punto del trazado a la parada que tiene delante o detrás (sección «\nameref{sec:trazado}»). Los
puntos del trazado se guardan con 7 decimales en las coordenadas y 2 en la distancia. Si la red no tiene trazados, la
base no tiene la tabla \texttt{shapes} y el motor trata \texttt{dist\_m} como desconocida; esto tiene consecuencias en
el respaldo por longitud (sección «\nameref{sec:respaldo-largo}»).

\section{Trazados que No Casan con sus Paradas}
\label{sec:trazados-gtfs}

Lo anterior supone que el trazado y las distancias de las paradas casan, y en el GTFS de Bizkaibus no siempre es así.
Con el horario del 3 de noviembre de 2026, 55 de los 440 trazados que usan los viajes están dibujados del final al
principio: la primera parada del viaje queda junto al último punto del trazado y la última junto al primero. Afectan a
670 de los 6\,130 viajes fusionados, repartidos en 25 líneas. En otros trazados el sentido es el bueno, pero los metros
de \texttt{shape\_dist\_traveled} no caen sobre la parada: unos van desplazados 200 o 300~m, y en otros la diferencia
crece hasta varios kilómetros al final del viaje. En total, 1\,417 viajes tenían alguna parada a más de 200~m del punto
del trazado que marcan sus metros.

El efecto en el motor es grave. La ventana de búsqueda (sección «\nameref{sec:ventana}») se define con los metros de
las paradas, así que, con el trazado al revés, el autobús se busca en la otra punta del recorrido: un autobús a 9~m de su
parada queda «a kilómetros» y el pasillo de 50~m lo descarta. Esos autobuses no se situaban con el GPS y, como el
aprendizaje sale de las posiciones, tampoco aprendían: los viajes con el trazado al revés tenían el 47\,\% de sus
tramos con tiempo aprendido, frente al 95\,\% del resto.

La construcción de la base corrige las dos cosas antes de publicarla (el programa, \texttt{scripts/build/shapes.php},
es parte del código privado de la aplicación; el algoritmo completo es este):

\begin{enumerate}
  \item \emph{Sentido}. Un trazado está al revés si, en la mayoría de sus viajes, la suma de las distancias de la
        primera parada al último punto y de la última al primero es menor que la de cada parada a su extremo. Se le da
        la vuelta: los puntos en orden inverso y $\sigma'=\sigma_{\max}-\sigma$.
  \item \emph{Metros de las paradas}. Si el trazado se ha dado la vuelta, o si alguna parada cae a más de 100~m del
        punto del trazado que marcan sus metros, se miden de nuevo las de ese viaje: cada parada se proyecta sobre el
        trazado a partir de la anterior y, como mucho, hasta tres veces la distancia en línea recta entre las dos más
        500~m. Entre dos pasos del autobús por la misma zona gana el primero, salvo que el otro quede 30~m más cerca.
\end{enumerate}

\noindent Con el horario de ese día se dieron la vuelta los 55 trazados y se midieron de nuevo las paradas de 2\,905
viajes; los viajes con alguna parada a más de 200~m de sus metros bajaron de 1\,417 a 26. La corrección es general:
se aplica al construir cualquier red, porque nada en el formato GTFS impide estos errores.

Tampoco obliga a publicar los metros. El GTFS de Lurraldebus, Bilbobus, Tuvisa y La Unión no trae los de las paradas,
y el de Lurraldebus, La Unión y Bilbobus tampoco los del trazado (el de Bilbobus los deja a 0 en todos los puntos menos
el último). Sin metros, el motor sitúa al autobús con la recta entre paradas (sección «\nameref{sec:sin-trazado}»).
Por eso la construcción los mide: los del trazado, sumando la distancia de cada punto al siguiente, y los de las
paradas, como en el paso~2.

Los feeds de Lurraldebus y Bilbobus, las dos redes con motor, no traían vehículos el 3 de noviembre de 2026, así que
la medida se hizo con Bizkaibus: se le quitaron los metros del GTFS y se compararon tres versiones del mismo horario
con las mismas capturas, 36 bloques de una hora del 2 y el 3 de noviembre con 2,08 millones de llegadas de paso
conocido (tabla~\ref{tab:sin-metros}). Los metros medidos caen a una mediana de 21~m de los del GTFS, el 95\,\% a
menos de 103~m, y aciertan casi como ellos. La recta entre paradas sitúa con el GPS un 5,6\,\% menos de llegadas y
acierta 1,8 puntos menos dentro de $\pm$2~min. Sin tiempos aprendidos, como en una red recién activada, las tres
versiones quedan a menos de medio punto: entonces lo que falta es aprender, no situar.

\begin{table}[ht]
\centering
\caption{Llegadas de Bizkaibus Según los Metros de las Paradas}\label{tab:sin-metros}
\begin{tabular}{@{}lrrrrrr@{}}
\toprule
 & \multicolumn{4}{c}{Con tiempos aprendidos} & \multicolumn{2}{c}{Sin aprender} \\
\cmidrule(lr){2-5}\cmidrule(l){6-7}
Metros de las paradas & Con GPS & $\pm$1 min & $\pm$2 min & Error & $\pm$2 min & Error \\
\midrule
Del GTFS & 97,4\,\% & 55,6\,\% & 78,4\,\% & 92~s & 51,5\,\% & 192~s \\
Medidos sobre el trazado & 97,3\,\% & 55,3\,\% & 78,2\,\% & 93~s & 51,6\,\% & 192~s \\
Ninguno (recta entre paradas) & 91,9\,\% & 54,0\,\% & 76,6\,\% & 105~s & 51,2\,\% & 194~s \\
\bottomrule
\end{tabular}

\nota{Las mismas 2,08 millones de llegadas en las tres versiones. «Del GTFS»: los metros del GTFS con la corrección de
la sección anterior. «Con GPS»: las que el motor sitúa con la posición del autobús (capa~1). «Error»: error absoluto medio de la hora de llegada. «Sin aprender»: sin tiempos por tramo, factores de ritmo ni correcciones por viaje.}
\end{table}

En las cuatro redes la medida es coherente: los metros nunca bajan de una parada a la siguiente y tienen todas sus
paradas a menos de 100~m de su punto del trazado el 99,5\,\% de los viajes de Bilbobus, el 95,9\,\% de los de
Lurraldebus, el 99,6\,\% de los de Tuvisa y el 98,9\,\% de los de La Unión. Medirlos alarga la construcción de cada
red entre 0 y 9~s.

\section{Fusión de Variantes y Alias de Viaje}
\label{sec:alias}

Un mismo servicio aparece en el GTFS como varios \texttt{trip\_id} que solo se diferencian en el calendario (por
ejemplo, uno para los lunes y otro para el resto de laborables). Para no mostrar salidas duplicadas, la construcción
de la base agrupa esos \texttt{trip\_id} en un único viaje. Dos \texttt{trip\_id} se fusionan cuando coinciden en:

\begin{enumerate}
  \item la línea (\texttt{route\_id}) y el número de viaje;
  \item la secuencia de paradas (la clave del recorrido es un resumen MD5 de los identificadores de parada en orden);
  \item el grupo de calendario;
  \item la \emph{forma temporal}: para cada parada, la llegada relativa a la salida de cabecera y el tiempo de espera
        en la parada (un resumen MD5 de esa lista);
\end{enumerate}

\noindent y, además, sus salidas de cabecera son consecutivas con una separación no mayor de 90~s
(\texttt{iDepartureClusterGapSeconds}). El viaje resultante reúne los días de sus miembros (unión de las máscaras de
días de la semana y de las fechas incluidas y excluidas). Cada \texttt{trip\_id} miembro queda registrado en
\texttt{trip\_aliases} apuntando al viaje fusionado. Esa tabla es la que permite emparejar exactamente el feed en
tiempo real con el horario (sección «\nameref{sec:emparejamiento}»): el feed usa los \texttt{trip\_id} originales.

\section{Por Qué el Horario No Sirve Como Referencia Absoluta}
\label{sec:interpolado}

En el GTFS de Bizkaibus las paradas intermedias llevan \texttt{timepoint=0}. Sus horas no son compromisos del
operador, sino una interpolación lineal en la distancia entre las horas de las cabeceras. La
Figura~\ref{fig:distancia-tiempo} muestra la consecuencia. El horario es la recta $T(s)$: una velocidad constante. La
marcha real $R(s)$ es más lenta en la ciudad y mucho más rápida en la autopista. A mitad de recorrido, $R$ y $T$
pueden estar separadas muchos minutos sin que el autobús vaya «adelantado» en ningún sentido operativo.

\begin{figure}[ht]
\centering
\caption{Distancia y Tiempo de un Viaje Ilustrativo}\label{fig:distancia-tiempo}
\begin{tikzpicture}
\begin{axis}[
  width=0.88\textwidth, height=6.4cm,
  xlabel={Distancia recorrida $s$ (km)}, ylabel={Hora (min desde la salida)},
  xmin=0, xmax=30, ymin=0, ymax=62,
  axis lines=left, tick label style={font=\footnotesize}, label style={font=\small},
  legend style={font=\footnotesize,draw=none,at={(0.02,0.98)},anchor=north west}]
  \addplot[thick,color=gris,dashed,domain=0:30] {2*x};
  \addlegendentry{Horario interpolado $T(s)$}
  \addplot[very thick,color=ruta,smooth] coordinates {(0,0) (4,12) (6,17) (8,20) (14,26) (20,32) (24,40) (28,52) (30,57)};
  \addlegendentry{Marcha real $R(s)$}
  \addplot[only marks,mark=*,mark size=2.6pt,color=gps] coordinates {(14,26)};
  \draw[gps,thick,-{Stealth}] (axis cs:14,26) -- (axis cs:24,46);
  \node[font=\footnotesize,color=gps,anchor=west] at (axis cs:14.4,24.5) {GPS: $(s_p,\tgps)$};
  \draw[decorate,decoration={brace,amplitude=4pt}] (axis cs:14.3,28) -- (axis cs:14.3,26.2);
\end{axis}
\end{tikzpicture}
\nota{La recta discontinua es el horario interpolado; la curva, la marcha real. Desde el punto GPS, el método traza una paralela al horario (flecha) y la corrige tramo a tramo con los tiempos aprendidos.}
\end{figure}

\section{El Feed en Tiempo Real}
\label{sec:feed}

El feed SIRI-VM de Bizkaibus es un documento XML público que se regenera aproximadamente una vez por minuto y contiene
unos 250 vehículos en hora punta. El ritmo ha variado: en octubre de 2026 salía una versión cada 55~s, y el 5 de
noviembre, cada 64~s de mediana, con cada posición 66~s después de tomarse (sección «\nameref{sec:llegando}»). El
motor lo descarga con un tiempo máximo de 8~s y lo guarda en caché 25~s. Para cada \texttt{VehicleActivity} extrae
los campos de la Tabla~\ref{tab:feed}.

\begin{table}[ht]
\centering
\caption{Campos que el Motor Lee del Feed SIRI-VM}\label{tab:feed}
\begin{tabularx}{\textwidth}{@{}llX@{}}
\toprule
Símbolo & Elemento SIRI & Uso \\
\midrule
$\rho$ & \texttt{MonitoredVehicleJourney/VehicleJourneyRef} & Referencia del viaje; igual al \texttt{trip\_id} del GTFS. \\
$\ell,\ \nu,\ d_0^{\mathrm{feed}}$ & (extraídos de $\rho$) & Línea (sin ceros a la izquierda, como en el horario), número de viaje y salida de cabecera, por la expresión~\eqref{eq:regex}. \\
$v$ & \texttt{VehicleRef} & Identificador del vehículo. \\
$n$ & \texttt{MonitoredCall/Order} (o \texttt{VisitNumber}) & Orden de la próxima parada en el viaje. \\
$(\varphi,\lambda)$ & \texttt{VehicleLocation/Latitude}, \texttt{Longitude} & Posición GPS. \\
$\tgps$ & \texttt{LocationRecordedAtTime} (o \texttt{RecordedAtTime}) & Hora de la posición. \\
$t_{\mathrm{rec}}$ & \texttt{VehicleActivity/RecordedAtTime} & Hora del registro del vehículo. \\
$\delta_{\mathrm{decl}}$ & \texttt{Delay} (duración ISO 8601) & Retraso declarado; ausente en Bizkaibus ($0$). \\
\bottomrule
\end{tabularx}
\end{table}

\subsection{Referencia del Viaje} Solo se aceptan las actividades cuya referencia encaja con
\begin{equation}
\label{eq:regex}
\texttt{\^{}trp\_[A-Za-z]+(\textbackslash d+)\_(\textbackslash d+)\_[A-Za-z0-9]+\_(\textbackslash d+)}
\end{equation}
cuyos tres grupos son la línea $\ell$, el número de viaje $\nu$ y la salida de cabecera $d_0^{\mathrm{feed}}$ en
segundos. Por ejemplo, \texttt{trp\_A3414\_808\_OP9LSEPT\_54000\_\ldots} da $\ell=3414$, $\nu=808$ y
$d_0^{\mathrm{feed}}=54\,000$ (15:00:00). Las demás se ignoran. La línea se guarda sin ceros a la izquierda, como el
\texttt{route\_id} del horario: \texttt{trp\_A0651\_737\_\ldots} da $\ell=651$, no \texttt{0651} (el \texttt{0651} es el
relleno del nombre público \texttt{A0651}, que el horario guarda aparte en \texttt{route\_short\_name}).

\subsection{Horas} Las horas del feed son fechas ISO 8601 con zona. Se convierten a \texttt{Europe/Madrid} y a segundos
desde medianoche, $3600\,H+60\,M+S$. Si la fecha convertida no es la fecha de hoy, el valor se descarta: una posición
de ayer no dice nada del viaje de hoy. Si la hora de la posición no es válida, también se descartan la latitud y la
longitud. Una entrada sin ninguna hora de hoy, ni $\tgps$ ni $t_{\mathrm{rec}}$, no se usa: ni en la predicción, ni en
el registro de pasos, ni en el mapa (\texttt{RealtimeMatcher::hasTime}). Mientras el feed se actualiza, todas la traen: 1\,471
de 1\,471 lecturas del 6 de noviembre de 2026 entre las 23:00 y las 23:59, la última hora antes de que se detenga. De 23:58 a 05:00, en cambio, el feed de Bizkaibus no
se actualiza: repite la lectura de las 23:58 con sus horas de entonces, aunque haya autobuses nocturnos en servicio
(capturas de las noches del 31 de octubre y del 6 de noviembre de 2026). Pasada la medianoche, ninguna de esas entradas
tiene hora de hoy (sección «\nameref{sec:mas-de-24h}»).

\subsection{Retraso Declarado} El retraso, cuando existe, es una duración $\pm$\texttt{PT}$h$\texttt{H}$m$\texttt{M}$s$\texttt{S};
se convierte a $\pm(3600h+60m+s)$. En Bizkaibus no aparece y el motor de autobús no lo usa
(sección «\nameref{sec:otras-redes}»).

Las entradas se agrupan por la clave «$\ell|\nu$», porque puede haber varios vehículos con la misma línea y número de
viaje en distintos días de calendario.

\subsection{Lo que el Feed No Dice} El feed solo trae posiciones: no publica cancelaciones, ni qué vehículo hará cada
viaje, ni refuerzos fuera del horario. Con las capturas del 3 de noviembre de 2026 se midieron cuatro consecuencias:

\begin{itemize}
  \item \emph{Viajes sin autobús.} En cuatro momentos del día, 947 viajes debían estar en marcha según el horario y
        187 (el 20\,\%) no tenían autobús en el feed. Siguiendo esos viajes en las capturas, la mayoría no falta: el
        12\,\% del total ya había llegado al final antes de su hora (el horario deja margen al final del recorrido),
        un 2\,\% aparece después con unos 15~min de retraso (mediana), otro 2\,\% desaparece un rato y vuelve, y otro
        2\,\% desaparece a mitad y no vuelve. Solo el 2\,\% no aparece en todo el día; desde fuera no se distingue
        un autobús sin GPS de un viaje que no se hizo.
  \item \emph{El número de viaje no es el vehículo.} De los viajes seguidos de una misma clave «$\ell|\nu$», solo
        el 11\,\% los hizo el mismo autobús, y el GTFS no trae \texttt{block\_id}. No se puede saber de antemano qué
        autobús hará el próximo viaje, ni arrastrar a él el retraso del anterior (sección «\nameref{sec:descartadas}»).
  \item \emph{Orden y parada.} El orden $n$ y la parada \texttt{StopPointRef} coinciden en el 97{,}9\,\% de los
        casos. En unas pocas líneas (3129, 2610, 2611 y 652) el operador cuenta una o varias paradas más que el GTFS;
        el motor usa $n$, y fiarse de la parada no mejoró la medida.
  \item \emph{Dónde está el autobús.} El feed no dice si un autobús está en una parada. \texttt{StopPointRef} es
        siempre la próxima: mientras espera en la primera parada de su viaje, apunta a la segunda (359 de 360 casos en
        nueve capturas del 3 de noviembre). \texttt{ArrivalStatus} vale \texttt{arrived} en los 2\,142 autobuses
        mirados, parados o en marcha, y no hay \texttt{PreviousCalls}. Solo el GPS lo dice (sección
        «\nameref{sec:cabecera}»).
\end{itemize}

\section{Notación}
\label{sec:notacion}

\begin{definicion}[Viaje]
Un viaje $J$ es una secuencia de $N\ge2$ paradas $p_1,\dots,p_N$. Para cada $i$ se conocen la llegada y la salida
programadas $a_i\le d_i$, las coordenadas $(\varphi_i,\lambda_i)$, el identificador $u_i$ de la parada y,
opcionalmente, la distancia $s_i$ a lo largo del trazado, con $s_1\le s_2\le\dots\le s_N$. El orden de la parada
$p_i$ es $i$ (\texttt{seq\_order}). La salida de cabecera es $d_1$. El trazado, si existe, es una polilínea
$\Sigma=\big((\varphi_j,\lambda_j,\sigma_j)\big)_{j=1}^{M}$ con $\sigma_j$ no decreciente.
\end{definicion}

\begin{definicion}[Tramo]
El tramo $i$ de $J$ es el par de paradas consecutivas $(p_i,p_{i+1})$. Su identificador es «$u_i>u_{i+1}$», su
duración programada es $a_{i+1}-a_i$ y su longitud, si se conocen las distancias, es $s_{i+1}-s_i$. Dos viajes de
líneas distintas que pasan por las mismas dos paradas comparten el tramo.
\end{definicion}

\begin{definicion}[Entrada del feed]
Una entrada es la tupla $e=(\rho,\allowbreak\ell,\allowbreak\nu,\allowbreak d_0^{\mathrm{feed}},\allowbreak v,\allowbreak n,\allowbreak\varphi,\allowbreak\lambda,\allowbreak\tgps,\allowbreak t_{\mathrm{rec}})$ de la
Tabla~\ref{tab:feed}, donde cualquiera de $n,\varphi,\lambda,\tgps,t_{\mathrm{rec}}$ puede faltar.
\end{definicion}

\begin{definicion}[Franja]
\label{def:franja}
La franja de una fecha $D$ y un instante $t$ es el par
\begin{gather*}
\franja(D,t)=\big(\text{tipo}(D),\ h(t)\big),\\
\text{tipo}(D)=\begin{cases}\texttt{fin\_semana}&\text{si $D$ es sábado o domingo},\\ \texttt{laborable}&\text{en otro caso,}\end{cases}
\end{gather*}
con $h=\lfloor t/3600\rfloor \bmod 24$ y
\[
h(t)=\begin{cases}
\texttt{madrugada} & 0\le h<6\\
\texttt{punta\_manana} & 6\le h<9\\
\texttt{manana} & 9\le h<13\\
\texttt{mediodia} & 13\le h<16\\
\texttt{punta\_tarde} & 16\le h<20\\
\texttt{noche} & 20\le h<24.
\end{cases}
\]
Se escribe como texto «\texttt{tipo.nombre}», por ejemplo \texttt{laborable.punta\_manana}. Hay, por tanto, doce
franjas.
\end{definicion}

\begin{observacion}
La fecha que se usa es la del día de servicio, no la del reloj. Un instante de 25:00 ($90\,000$~s) de un miércoles
es la franja \texttt{laborable.madrugada} del miércoles, aunque en el calendario ya sea jueves. En los tiempos por
tramo, los festivos entre semana se tratan como laborables: se comprobó que su tráfico se parece más al de un laborable
que al de un fin de semana (sección «\nameref{sec:descartadas}»). La capa 3 es la excepción: para ella el festivo
cuenta como fin de semana en todo (franja de los tramos, factor de ritmo, parte aprendida y corrección por viaje),
porque así se midió mejor (sección «\nameref{sec:capa3}»).
\end{observacion}

\begin{definicion}[Hora actual]
$\tnow$ es la hora del servidor en el momento de responder, en segundos desde la medianoche de Madrid del día actual.
Para un viaje del día anterior se usa $\tnow+86\,400$, en las horas de su horario (sección «\nameref{sec:mas-de-24h}»).
\end{definicion}

\chapter{Emparejamiento y Localización}
\label{cap:localizacion}

Antes de predecir nada hay que responder dos preguntas: qué viaje del horario está haciendo cada autobús del feed, y
en qué punto de ese viaje se encuentra. La respuesta a la segunda es una hora programada, $\Tpos$: la hora a la que el
horario sitúa al autobús en el lugar donde el GPS dice que está.

\section{Emparejamiento con el Horario}
\label{sec:emparejamiento}

\subsection{Emparejamiento Exacto} La referencia $\rho$ del feed es un \texttt{trip\_id} del GTFS. Con la tabla
\texttt{trip\_aliases} (sección «\nameref{sec:alias}») se obtiene directamente el viaje fusionado $J=\mathrm{alias}(\rho)$.
Al construir el emparejador se consultan a la vez los alias de todas las referencias del feed; para cada viaje $J$ se
guarda la \emph{primera} entrada del feed cuyo alias es $J$. En las pruebas con el feed real, el 100\,\% de los
autobuses se emparejó así.

\subsection{Emparejamiento de Respaldo} Si la base de datos no tiene la tabla de alias, o el viaje no aparece en ella,
se buscan las entradas con la misma clave «$\ell|\nu$» que el viaje y se elige la de salida de cabecera más próxima:
\begin{equation}
e^\ast=\argmin_{e\in C}\ \big|d_0^{\mathrm{feed}}(e)-d_1(J)\big|
\quad\text{sujeto a}\quad \big|d_0^{\mathrm{feed}}(e)-d_1(J)\big|\le 180\ \mathrm{s},
\end{equation}
donde $C$ excluye las entradas que ya tienen un emparejamiento exacto con otro viaje. Si hay empate, gana la primera
encontrada; si ninguna cumple la tolerancia (\texttt{MATCH\_TOLERANCE\_SECONDS}), el viaje no tiene datos en vivo.

\section{Proyección de un Punto Sobre un Segmento}
\label{sec:proyeccion}

Toda la localización se basa en una operación: proyectar la posición del autobús sobre un segmento rectilíneo $AB$
(entre dos puntos del trazado, o entre dos paradas).

\subsection{Plano Local} Se usa una proyección equirrectangular centrada en el autobús, que queda en el origen. Un
punto $Q=(\varphi_Q,\lambda_Q)$ pasa a
\begin{equation}
\label{eq:plano}
x_Q=(\lambda_Q-\lambda)\cdot\cos\varphi\cdot 111\,320,\qquad y_Q=(\varphi_Q-\varphi)\cdot 111\,320,
\end{equation}
en metros, con $(\varphi,\lambda)$ la posición del autobús (\texttt{METERS\_PER\_DEGREE}).

\subsection{Proyección} Sean $\mathbf{a}=(x_A,y_A)$ y $\mathbf{b}=(x_B,y_B)$ y $\mathbf{u}=\mathbf{b}-\mathbf{a}$.
\begin{equation}
\label{eq:proyeccion}
f=\begin{cases}\clamp\!\left(-\dfrac{\mathbf{a}\cdot\mathbf{u}}{\|\mathbf{u}\|^2},\,0,\,1\right) & \text{si } \|\mathbf{u}\|^2>0\\[1ex]
0 & \text{si } \|\mathbf{u}\|^2=0\end{cases}
\qquad
d_\perp=\|\mathbf{a}+f\,\mathbf{u}\|,\qquad L=\|\mathbf{u}\|.
\end{equation}
$f\in[0,1]$ es la fracción del segmento hasta el punto más próximo al autobús, $d_\perp$ es la distancia del autobús a
ese punto y $L$ la longitud del segmento.

\begin{proposicion}
El punto $\mathbf{a}+f\,\mathbf{u}$ con $f$ dado por~\eqref{eq:proyeccion} es el punto del segmento más próximo al
origen.
\end{proposicion}
\begin{proof}
Si $\|\mathbf{u}\|=0$ el segmento es un punto. Si no, $g(t)=\|\mathbf{a}+t\mathbf{u}\|^2=\|\mathbf{a}\|^2+2t\,
\mathbf{a}\cdot\mathbf{u}+t^2\|\mathbf{u}\|^2$ es una parábola convexa con mínimo en
$t^\ast=-\mathbf{a}\cdot\mathbf{u}/\|\mathbf{u}\|^2$. Restringida a $[0,1]$, su mínimo es $\clamp(t^\ast,0,1)$, porque
$g$ es decreciente a la izquierda de $t^\ast$ y creciente a la derecha.
\end{proof}

\subsection{Error de la Aproximación} La proyección~\eqref{eq:plano} usa $\cos\varphi$ en la latitud del autobús. En
un punto a $\Delta\varphi$ radianes de distancia norte–sur, el factor correcto sería $\cos(\varphi+\Delta\varphi)
\approx\cos\varphi\,(1-\Delta\varphi\tan\varphi)$. A la latitud de Bilbao ($\varphi\approx43{,}26^\circ$,
$\tan\varphi\approx0{,}94$) y a 3~km ($\Delta\varphi\approx4{,}7\cdot10^{-4}$~rad), el error relativo es del orden
del 0,04\,\%, poco más de un metro. El factor de 111\,320~m por grado sobrestima en torno a un 0,2\,\% la longitud
de un grado de latitud a esa altura, unos 6~m en 3~km. Ambos errores son despreciables frente a los pasillos de
tolerancia (50 y 150~m). Además, la hora programada se obtiene de la \emph{fracción} $f$ y de las distancias del
propio GTFS, no de distancias medidas en el plano, así que un error de escala uniforme no altera $\Tpos$. El de este
plano no es del todo uniforme (frente al elipsoide WGS84, la longitud queda un 0,16\,\% corta y la latitud un
0,2\,\% larga), pero esa diferencia entre los ejes mueve el punto proyectado menos de un 0,4\,\% de la longitud del
segmento: también despreciable.

\begin{figure}[ht]
\centering
\caption{Proyección del Autobús Sobre un Segmento: Fracción $f$ y Distancia Perpendicular $d_\perp$}\label{fig:proyeccion}
\begin{tikzpicture}[scale=1.15]
  \coordinate (A) at (0,0);
  \coordinate (B) at (7,1.4);
  \coordinate (P) at (3.6,2.1);
  \coordinate (H) at ($(A)!(P)!(B)$);
  \draw[very thick,ruta] (A) -- (B);
  \fill[ruta] (A) circle (2.2pt) node[below left] {$A$};
  \fill[ruta] (B) circle (2.2pt) node[below right] {$B$};
  \fill[gps] (P) circle (2.6pt) node[above] {autobús $(0,0)$};
  \draw[gps,dashed,thick] (P) -- (H) node[midway,right] {$d_\perp$};
  \fill[black] (H) circle (1.6pt);
  \draw[decorate,decoration={brace,amplitude=5pt,mirror},thick,gris] ($(A)+(0.05,-0.25)$) -- ($(H)+(0.05,-0.25)$) node[midway,below=4pt,sloped,font=\small] {$f\cdot L$};
  \draw[decorate,decoration={brace,amplitude=5pt,mirror},thick,gris] ($(A)+(0.1,-0.95)$) -- ($(B)+(0.1,-0.95)$) node[midway,below=4pt,sloped,font=\small] {$L$};
\end{tikzpicture}
\end{figure}

\section{Ventana de Búsqueda}
\label{sec:ventana}

El feed indica el orden $n$ de la próxima parada. Esa información acota dónde puede estar el autobús: en algún punto
poco antes de $p_n$. La búsqueda se limita a los tramos
\begin{equation}
\label{eq:ventana}
i_{\min}=\max(1,\,n-3),\qquad i_{\max}=\min(N-1,\,n+2),
\end{equation}
es decir, del tramo $n-3$ al $n+2$: los tres que llevan hasta la próxima parada y los tres que salen de ella o de
las siguientes (\texttt{SEGMENTS\_BEHIND\_NEXT\_STOP} y \texttt{SEGMENTS\_AHEAD\_OF\_NEXT\_STOP}, 3 y 2, son los
desplazamientos respecto a $n$). La ventana cubre las paradas $p_{i_{\min}},\dots,p_{i_{\max}+1}$, de $p_{n-3}$ a
$p_{n+3}$ cuando el viaje las tiene. Mirar
hacia atrás tolera que el feed cambie de próxima parada con algo de retraso; mirar hacia delante, que lo haga con
antelación. Si $n$ no es el orden de ninguna parada del viaje, no hay posición.

\section{Con Trazado}
\label{sec:trazado}

Se usa este método cuando el viaje tiene trazado y se conocen las distancias $s_i$ de todas las paradas de la
ventana. Sean $S_{\mathrm{lo}}=s_{i_{\min}}$ y $S_{\mathrm{hi}}=s_{i_{\max}+1}$.

\begin{enumerate}
  \item \textbf{Candidatos.} Para cada par de puntos consecutivos del trazado $A=\Sigma_j$, $B=\Sigma_{j+1}$ que no
        quedan fuera de la ventana (es decir, salvo si $\sigma_B<S_{\mathrm{lo}}$ o $\sigma_A>S_{\mathrm{hi}}$), se
        proyecta el autobús con~\eqref{eq:proyeccion} y se calcula su posición a lo largo del trazado,
        $s^\ast=\sigma_A+f\,(\sigma_B-\sigma_A)$.
  \item \textbf{Dos grupos.} El candidato es \emph{esperado} si $s^\ast\le s_n+100$~m
        (\texttt{BEYOND\_NEXT\_STOP\_SLACK\_METERS}) y \emph{más allá} en otro caso. En cada grupo se queda el de menor
        $d_\perp$; con empate, el primero.
  \item \textbf{Desempate.} Se elige el mejor esperado salvo que el mejor de más allá esté al menos 20~m más cerca del
        recorrido (\texttt{AMBIGUITY\_MARGIN\_METERS}):
        \[
        \text{elegido}=\begin{cases}\text{más allá}&\text{si existe y } \big(\text{no hay esperado o }d_\perp^{\text{más allá}}+20<d_\perp^{\text{esperado}}\big)\\
        \text{esperado}&\text{en otro caso.}\end{cases}
        \]
  \item \textbf{Pasillo.} La posición es válida si $d_\perp\le50$~m (\texttt{MAX\_DISTANCE\_FROM\_SHAPE\_METERS}).
  \item \textbf{Hora programada.} Se acota $s^\ast\leftarrow\clamp(s^\ast,S_{\mathrm{lo}},S_{\mathrm{hi}})$ y se
        busca el primer tramo $i\in[i_{\min},i_{\max}]$ con $s^\ast\le s_{i+1}$ (o el último de la ventana). Entonces
        \begin{equation}
        \label{eq:tpos}
        f_i=\begin{cases}\clamp\!\Big(\dfrac{s^\ast-s_i}{s_{i+1}-s_i},0,1\Big)&\text{si }s_{i+1}>s_i\\ 0&\text{en otro caso}\end{cases}
        \qquad
        \Tpos=\redondeo\big(d_i+f_i\,(a_{i+1}-d_i)\big).
        \end{equation}
\end{enumerate}

La hora programada se interpola entre la \emph{salida} de la parada de inicio del tramo y la \emph{llegada} a la de
fin: dentro del tramo el autobús está en marcha, y la espera en la parada no forma parte de él. Como los segmentos del
trazado que tocan el borde de la ventana entran como candidatos y $s^\ast$ se acota a $[S_{\mathrm{lo}},S_{\mathrm{hi}}]$,
un GPS algo más allá de un extremo queda en ese extremo: con un trazado que solo tuviera puntos en las paradas, un
autobús 400~m después de $p_{i_{\max}+1}$ y dentro del pasillo toma $\Tpos=a_{i_{\max}+1}$.

Hacen falta dos grupos porque hay recorridos que vuelven sobre la misma calle: una ida y vuelta por el mismo eje, o
una rotonda. Un punto GPS en esa calle tiene dos proyecciones casi igual de cercanas. El grupo esperado prefiere la
que es coherente con la próxima parada que declara el feed. Solo se acepta una proyección más adelantada si es
claramente mejor (20~m), lo que cubre el caso en que el feed tarda en actualizar la próxima parada. Los 100~m de margen
tras la próxima parada absorben que el feed cambia de parada con retraso.

\begin{figure}[ht]
\centering
\caption{Ventana de Búsqueda Alrededor de la Próxima Parada}\label{fig:ventana}
\begin{tikzpicture}[scale=1.0]
  \draw[line width=6pt,gris!25,line cap=round] (0,0) -- (13,0);
  \draw[line width=6pt,ruta,line cap=round] (1.2,0) -- (11.4,0);
  \foreach \x/\l in {1.2/$p_{n-3}$,3.4/$p_{n-2}$,5.6/$p_{n-1}$,7.8/$p_{n}$,9.0/$p_{n+1}$,10.2/$p_{n+2}$,11.4/$p_{n+3}$} {
    \fill[white] (\x,0) circle (5pt); \draw[ruta,thick] (\x,0) circle (5pt);
    \node[below=7pt,font=\small] at (\x,0) {\l};
  }
  \fill[ruta] (7.8,0) circle (5pt);
  \draw[gps,dashed,thick] (7.8,0) -- (7.8,1.3) -- (8.6,1.3) -- (8.6,0);
  \node[gps,font=\footnotesize,above] at (8.2,1.3) {$+100$ m};
  \fill[gps] (6.6,0.55) circle (3pt);
  \draw[gps,thick] (6.6,0.55) -- (6.6,0);
  \node[gps,font=\footnotesize,above] at (6.6,0.6) {GPS};
  \draw[decorate,decoration={brace,amplitude=5pt},thick] (1.2,-0.95) -- (11.4,-0.95) node[midway,below=6pt,font=\small] {ventana: tramos $i_{\min}=n-3$ a $i_{\max}=n+2$, paradas $p_{n-3}$ a $p_{n+3}$};
\end{tikzpicture}
\nota{Se marca el límite de los candidatos «esperados» alrededor de la próxima parada $p_n$.}
\end{figure}

\section{Sin Trazado}
\label{sec:sin-trazado}

Si el viaje no tiene trazado, o falta la distancia de alguna parada de la ventana, se proyecta sobre las rectas entre
paradas. Para cada tramo $i$ de la ventana, en orden, se proyecta el autobús sobre el segmento $p_ip_{i+1}$ (con las
coordenadas de las paradas) y se conserva el de menor $d_\perp$ (con empate, el primero). Para ese tramo se calculan
\begin{equation}
c_i=\max\big(150,\ 0{,}15\,L_i\big),\qquad \Tpos=\redondeo\big(d_i+f\,(a_{i+1}-d_i)\big),
\end{equation}
y la posición es válida si $d_\perp\le c_i$ (\texttt{MAX\_DISTANCE\_FROM\_ROUTE\_METERS},
\texttt{MAX\_DISTANCE\_FROM\_ROUTE\_SEGMENT\_RATIO}). El pasillo es más ancho porque la calle real se aparta de la
recta entre dos paradas, tanto más cuanto más largo es el tramo. Sin trazado no hay dos grupos ni desempate: se queda
el tramo más cercano, y en un recorrido que vuelve sobre sí mismo la única protección es la ventana.

Desde que la construcción mide los metros de las redes que no los publican (sección «\nameref{sec:trazados-gtfs}»), este
caso solo queda para los viajes que no tienen ningún trazado: el resto se sitúa con el pasillo de 50~m.

\section{Condiciones Para Usar el GPS}
\label{sec:condiciones-gps}

La posición de una entrada $e$ en el viaje $J$ existe solo si se cumplen todas estas condiciones:

\begin{enumerate}
  \item $e$ tiene latitud, longitud, hora de la posición $\tgps$ y orden $n$ de la próxima parada;
  \item el dato es reciente: $\tnow-\tgps\le600$~s (\texttt{MAX\_LOCATION\_AGE\_SECONDS});
  \item el viaje tiene al menos dos paradas;
  \item el viaje ha empezado o está a punto: $\tgps\ge d_1-60$~s (\texttt{NOT\_STARTED\_GRACE\_SECONDS}); si no, solo
        cabe la posición en la cabecera de la sección «\nameref{sec:cabecera}»;
  \item $n$ es el orden de una parada del viaje;
  \item la proyección es válida (secciones «\nameref{sec:trazado}» o «\nameref{sec:sin-trazado}»);
  \item el GPS no contradice a la parada registrada (sección «\nameref{sec:desconfianza}»).
\end{enumerate}

Si se cumplen, la \emph{posición} es la cuádrupla $\big(\Tpos,\ \tgps,\ d_\perp,\ \text{método}\big)$ y el
\emph{desfase en la posición} es
\begin{equation}
\label{eq:desfase-pos}
\delta_{\mathrm{pos}}=\tgps-\Tpos.
\end{equation}
Un valor positivo indica que el autobús llega a ese punto más tarde de lo que dice el horario. Como el horario está
interpolado, este desfase \emph{no} se muestra al usuario tal cual (sección «\nameref{sec:retraso-mostrado}»).

\begin{ejemplo}
\label{ej:posicion}
Considérese el viaje sintético del Apéndice~\ref{ap:vectores}: cinco paradas en línea recta hacia el este a lo largo
del paralelo $43{,}26^\circ$, a 0, 400, 1\,000, 2\,500 y 3\,200~m, con horas 8:00:00, 8:01:30, 8:03:30, 8:08:30 y
8:10:30 ($28\,800$, $28\,890$, $29\,010$, $29\,310$ y $29\,430$~s). Un autobús declara como próxima parada $p_3$
($n=3$) y su GPS está a 700~m del origen y $0{,}0001^\circ$ al norte ($\approx11$~m), con $\tgps=28\,950$.

La ventana es $i_{\min}=\max(1,0)=1$, $i_{\max}=\min(4,5)=4$. La mejor proyección cae en el tramo $p_2p_3$, con
$d_\perp=11{,}132$~m $\le50$ y $s^\ast=700$~m. Entonces $f_2=(700-400)/(1000-400)=0{,}5$ y
$\Tpos=\redondeo(28\,890+0{,}5\cdot(29\,010-28\,890))=28\,950$. El desfase en la posición es
$\delta_{\mathrm{pos}}=0$: el autobús está donde el horario dice. Sin trazado se obtiene el mismo $\Tpos$, por
proyección sobre la recta $p_2p_3$.
\end{ejemplo}

\section{En la Cabecera, Antes de Salir}
\label{sec:cabecera}

Un autobús asignado a un viaje que aún no ha salido no cumple la condición~4. Muchos de ellos esperan en la primera
parada: en la captura del 3 de noviembre de 2026 a las 14:47, 36 de los 259 autobuses estaban antes de su hora y 34 de
ellos a menos de 50~m de su primera parada. Sin posición, la parada de cabecera y las siguientes salían solo con el
horario aunque el autobús estuviera allí. Por eso, cuando $\tgps<d_1-60$, la entrada tiene la posición
\begin{equation}
\label{eq:cabecera}
\Tpos=\min(a_1,d_1),\qquad \tgps'=d_1,\qquad\text{método «cabecera»},
\end{equation}
si el autobús está a 150~m o menos de la primera parada (\texttt{MAX\_DISTANCE\_FROM\_ROUTE\_METERS}) y su posición
se tomó como mucho 60~min antes de su salida, $d_1-\tgps\le3\,600$~s (\texttt{MAX\_WAIT\_AT\_ORIGIN\_SECONDS}). Más lejos puede estar acabando otro viaje, y más
de una hora antes la asignación es poco segura. La hora de la posición pasa a ser la de salida, $d_1$: el autobús no
ha recorrido nada y saldrá, como pronto, a su hora.

Con esta posición la fila sale «en directo», pero su llegada no se calcula desde $\tgps'$: es la de la capa 3
(sección «\nameref{sec:cabecera-eta}»). Tampoco cuenta la primera parada como pasada, aunque el feed ya apunte a la
segunda ($n=2$) mientras el autobús espera. La aplicación lo dice como «En la parada», pero solo en la fila de la
primera parada, decidida por su orden en el recorrido y no por su hora: dos paradas seguidas pueden tener la misma hora
publicada, y en un viaje circular la primera y la última son la misma. Las posiciones en la cabecera no entran en las
medidas del aprendizaje, porque no describen ningún recorrido. El 97\,\% de las veces que una entrada se sitúa así, ese
autobús hace después el viaje desde allí (sección «\nameref{sec:fiabilidad-estados}»).


\chapter{Modelo del Tiempo de Recorrido}
\label{cap:modelo}

El capítulo anterior sitúa al autobús en el horario: da la hora $\Tpos$ a la que el horario lo pone en el punto donde
está. Falta saber cuánto tardará en llegar desde ahí hasta una parada. Este capítulo explica por qué se trabaja con
diferencias de horario, qué unidad de tiempo se aprende (el tramo entre paradas, por franja), cómo se estima cada
tiempo, qué se hace cuando un tramo no tiene datos y cómo se suman los tramos para obtener el tiempo esperado entre dos
puntos del recorrido.

\section{Por Qué Diferencias de Horario}
\label{sec:diferencias}

Sea $J$ un viaje y $s\in[s_1,s_N]$ un punto de su recorrido, medido en metros desde la cabecera. Se definen tres
funciones:

\begin{itemize}
  \item $T(s)$, la hora a la que el \emph{horario} pone al autobús en $s$. En las paradas, $T(s_i)=a_i$; entre ellas,
        se interpola linealmente (ecuación~\eqref{eq:tpos}).
  \item $R(s)$, la hora \emph{real} a la que el autobús pasa por $s$. Es desconocida salvo en el punto donde se ha
        visto al autobús.
  \item $\Delta(s)=R(s)-T(s)$, el \emph{desfase} en $s$.
\end{itemize}

El autobús se ha visto en $s_p$ a la hora $\tgps=R(s_p)$, y la parada de interés está en $s_d>s_p$, con hora
publicada $a_\dest=T(s_d)$. La cantidad que se quiere predecir es $R(s_d)$.

\begin{proposicion}[El error de las diferencias]
\label{prop:diferencias}
Sean dos estimadores de $R(s_d)$:
\[
\widehat R_{\mathrm{hor}}=T(s_d)\qquad\text{(solo horario)},\qquad
\widehat R_{\mathrm{dif}}=\tgps+\big(T(s_d)-T(s_p)\big)\qquad\text{(diferencias)}.
\]
Sus errores son
\[
R(s_d)-\widehat R_{\mathrm{hor}}=\Delta(s_d),\qquad
R(s_d)-\widehat R_{\mathrm{dif}}=\Delta(s_d)-\Delta(s_p).
\]
Más en general, si $\widehat E$ es cualquier estimación del tiempo real entre $s_p$ y $s_d$ y
$\widehat R=\tgps+\widehat E$, el error es $\big(R(s_d)-R(s_p)\big)-\widehat E$, que no depende de $\Delta(s_p)$.
\end{proposicion}

\begin{proof}
La primera igualdad es la definición de $\Delta$. Para la segunda, como $\tgps=R(s_p)=T(s_p)+\Delta(s_p)$,
\begin{align*}
R(s_d)-\tgps-T(s_d)+T(s_p)&=\big(T(s_d)+\Delta(s_d)\big)-\big(T(s_p)+\Delta(s_p)\big)-T(s_d)+T(s_p)\\
&=\Delta(s_d)-\Delta(s_p).
\end{align*}
La tercera es inmediata: $R(s_d)-\tgps-\widehat E=\big(R(s_d)-R(s_p)\big)-\widehat E$.
\end{proof}

La proposición tiene tres consecuencias que guían todo el diseño.

\begin{enumerate}
  \item \textbf{El retraso acumulado no importa.} Un autobús que va veinte minutos tarde desde la cabecera tiene
        $\Delta(s_p)\approx\Delta(s_d)\approx1200$~s; el horario solo yerra en 20 minutos, las diferencias en casi
        nada. La predicción no necesita saber cuánto retraso lleva el autobús, solo dónde está.
  \item \textbf{El error de forma del horario sí importa.} Como el horario intermedio es una interpolación a velocidad
        constante (sección «\nameref{sec:interpolado}»), $T(s_d)-T(s_p)$ puede estar muy lejos del tiempo real entre esos dos
        puntos: en un tramo de autopista, el doble. Las diferencias de horario eliminan el desfase acumulado, pero no
        este error. Por eso se sustituye $T(s_d)-T(s_p)$ por una estimación aprendida $\widehat E$ del tiempo real
        (el resto de este capítulo).
  \item \textbf{El error crece con la distancia.} El error $\big(R(s_d)-R(s_p)\big)-\widehat E$ suma la variación de
        todos los tramos entre el autobús y la parada. Si los tiempos de los tramos fueran independientes, su varianza
        sería la suma de las varianzas de cada tramo y crecería aproximadamente de forma lineal con el número de
        tramos. Es lo que se observa: la precisión cae de forma regular con el horizonte (capítulo «\nameref{cap:evaluacion}»).
\end{enumerate}

\section{La Unidad de Aprendizaje: El Tramo}
\label{sec:unidad-tramo}

El error de forma del horario se produce tramo a tramo: en cada tramo de autopista el horario se queda corto y en
cada tramo urbano congestionado se queda largo. Si se aprende el tiempo real de cada tramo, el error de forma
desaparece al sumar. Esta es la razón de aprender \emph{por tramo} y no por línea:

\begin{itemize}
  \item un factor único por línea (por ejemplo, «la línea 3414 tarda un 4\,\% menos que su horario») no puede corregir
        un viaje que combina ciudad y autopista, porque la corrección tiene signos distintos en cada parte;
  \item el tramo se identifica por sus dos paradas (definición de tramo, sección «\nameref{sec:notacion}»), de modo que las
        medidas de todas las líneas que lo recorren se suman y el modelo aprende más deprisa;
  \item el tiempo de un tramo es una magnitud física estable (depende de la calle, los semáforos y la demanda de esa
        parada) y fácil de auditar: se puede comparar con lo que tarda uno mismo en recorrerlo.
\end{itemize}

El tiempo de un tramo cambia con la hora y el día. Se aprende por \emph{franja} (definición~\ref{def:franja}): doce
combinaciones de tipo de día (laborable, fin de semana) y periodo del día (madrugada, punta de la mañana, mañana,
mediodía, punta de la tarde, noche). Las fronteras de las franjas siguen los periodos de demanda habituales en el
transporte urbano. Con los datos de octubre de 2026, los tiempos de un mismo tramo difieren entre franjas un 8--13\,\%
de media, y la noche es un 12\,\% más rápida que la punta de la tarde.

\section{Estimador del Tiempo de un Tramo}
\label{sec:estimador-tramo}

\subsection{Las Muestras}

El capítulo «\nameref{cap:aprendizaje}» explica cómo se reconstruye, a partir de las capturas del feed, el instante real en
que un autobús pasa por cada parada. Con dos pasos consecutivos se obtiene una \emph{muestra de tramo}:

\begin{definicion}[Muestra de tramo]
Una muestra de tramo es una tupla $x=(\kappa,\ j,\ b,\ y)$. El primer elemento, $\kappa=u_i{>}u_{i+1}$, es el
identificador del tramo. $j$ es la \emph{pista}: el par vehículo y viaje en una fecha concreta, que identifica un
recorrido físico de un autobús. $b$ es la franja del instante en que el autobús pasó por $p_i$, e $y>0$ es el tiempo
real, en segundos, que tardó en ir de $p_i$ a $p_{i+1}$.
\end{definicion}

\subsection{Un Valor por Pista}

Para cada tramo $\kappa$ y cada \emph{grupo} $g$ (una franja concreta, o el grupo especial \texttt{all} que reúne
todas las franjas) se forma el conjunto de pistas que tienen alguna muestra de $\kappa$ en $g$. Cada pista aporta
\emph{un solo valor}: si una pista tuviera varias muestras del mismo tramo y grupo, se queda la \emph{última} en el
orden del fichero de muestras.\footnote{En la práctica una pista recorre cada tramo una sola vez. La regla de «la
última» importa solo para reproducir exactamente el resultado con datos duplicados; los vectores de prueba la
comprueban (Apéndice~\ref{ap:vectores}).} Sea $Y_{\kappa,g}$ la lista de esos valores y $n_{\kappa,g}=|Y_{\kappa,g}|$
el número de pistas distintas.

Un valor por pista evita que un autobús visto muchas veces pese más que otro visto una vez, y es el sentido en que el
estimador es \emph{robusto}: un autobús que se queda parado veinte minutos aporta un valor extremo, pero uno solo.

\subsection{Mediana y Contracción Hacia el Tramo a Cualquier Hora}
\label{sec:contraccion-tramo}

Para cada tramo $\kappa$ se calcula la mediana baja a cualquier hora,
\begin{equation}
m_{\kappa}=\mediana\big(Y_{\kappa,\texttt{all}}\big),
\end{equation}
y, para cada franja $b$ con $n_{\kappa,b}\ge2$ pistas, la mediana de la franja $m_{\kappa,b}=\mediana(Y_{\kappa,b})$
contraída hacia $m_\kappa$:
\begin{equation}
\label{eq:contraccion-tramo}
\boxed{\;\tau_{\kappa,b}=\redondeo\!\left(\frac{n_{\kappa,b}\,m_{\kappa,b}+n_0\,m_\kappa}{n_{\kappa,b}+n_0}\right),
\qquad n_0=5\;}
\end{equation}
(\texttt{SEGMENT\_BAND\_PRIOR\_TRACKS}). El valor a cualquier hora se guarda como
$\tau_{\kappa,\texttt{all}}=\redondeo(m_\kappa)$, pero solo si el tramo tiene muestras en al menos dos franjas
distintas. Las reglas completas de qué se guarda son:

\begin{enumerate}
  \item Un tramo se considera solo si $n_{\kappa,\texttt{all}}\ge2$ (\texttt{MIN\_TRACKS\_PER\_SEGMENT}).
  \item Una franja $b$ se guarda solo si $n_{\kappa,b}\ge2$.
  \item El grupo \texttt{all} se guarda solo si hay muestras de al menos dos franjas distintas, aunque alguna de ellas
        tenga una sola pista. Con una sola franja, el valor a cualquier hora sería el de esa franja con otro nombre, y
        se usaría en otras franjas para las que no hay ninguna medida: lo medido a mediodía acabaría usándose en hora
        punta.
  \item El tramo se guarda si se guardó al menos un grupo. Se guarda también $n_{\kappa,\texttt{all}}$ como
        información (clave \texttt{n}).
\end{enumerate}

Todos los valores se guardan como enteros (segundos, con $\redondeo$).

La ecuación~\eqref{eq:contraccion-tramo} es una media ponderada entre lo que dice la
franja y lo que dice el tramo a cualquier hora, con pesos $n_{\kappa,b}$ y $n_0$. Es exactamente la forma de la media
a posteriori de un modelo normal jerárquico \citep{gelman2013}: si las pistas de la franja $b$ se distribuyen
alrededor de su media $\mu_b$ con varianza $\sigma^2$, y las medias de las franjas se distribuyen alrededor de la media
del tramo $\mu$ con varianza $\sigma_0^2$, la media a posteriori de $\mu_b$ es
\[
\mathbb{E}[\mu_b\mid\text{datos}]=\frac{n\,\bar y_b+(\sigma^2/\sigma_0^2)\,\mu}{n+\sigma^2/\sigma_0^2}.
\]
La constante $n_0$ representa el cociente $\sigma^2/\sigma_0^2$: cuántas pistas de la franja hacen falta para que su
propia mediana pese lo mismo que la del tramo. Con pocas pistas manda el tramo; con muchas, la franja. El método
sustituye las medias por medianas para conservar la robustez, y fija $n_0$ por validación (no por una estimación de las
varianzas), lo que es más sencillo y se comprobó suficiente.

\begin{proposicion}[Por qué contraer reduce el error]
\label{prop:contraccion}
Sea $\hat\theta_b$ un estimador insesgado de $\theta_b$ con varianza $\sigma^2/n$, y $\mu$ un valor fijo con
$\theta_b-\mu=\beta$. Para el estimador contraído $\tilde\theta_b=w\hat\theta_b+(1-w)\mu$, con $w=n/(n+n_0)$, el error
cuadrático medio es
\[
\mathrm{ECM}(\tilde\theta_b)=w^2\frac{\sigma^2}{n}+(1-w)^2\beta^2 ,
\]
que es menor que el de $\hat\theta_b$ ($\sigma^2/n$) siempre que $\beta^2<\sigma^2\,(2n+n_0)/(n\,n_0)$. En particular,
la ganancia es mayor cuanto menor es $n$.
\end{proposicion}

\begin{proof}
$\tilde\theta_b-\theta_b=w(\hat\theta_b-\theta_b)-(1-w)\beta$; el primer término tiene media cero y varianza
$w^2\sigma^2/n$ y el segundo es constante, de donde sale la expresión. La desigualdad
$w^2\sigma^2/n+(1-w)^2\beta^2<\sigma^2/n$ equivale a $(1-w)^2\beta^2<(1-w^2)\sigma^2/n$, es decir,
$\beta^2<\frac{1+w}{1-w}\cdot\frac{\sigma^2}{n}$, y con $w=n/(n+n_0)$ se tiene $\frac{1+w}{1-w}=\frac{2n+n_0}{n_0}$.
\end{proof}

La proposición~\ref{prop:contraccion} explica el comportamiento medido (sección «\nameref{sec:eval-contraccion}»): la contracción mejora las franjas
con pocas pistas, donde la mediana es ruidosa, y apenas cambia las que tienen muchas. Un $n_0$ demasiado grande sesga
las franjas que de verdad son distintas (el término $\beta^2$), y por eso el valor 10 resultó peor que el 5 a 12--20
minutos. La proposición supone un estimador insesgado y un $\mu$ fijo; el método usa medianas bajas y estima $\mu$
($m_\kappa$) con datos que incluyen las pistas de la propia franja, así que es una justificación heurística, no una
garantía: el valor de $n_0$ lo decide la medida.

\begin{ejemplo}
\label{ej:contraccion}
Un tramo $A{>}B$ tiene dos pistas en \texttt{laborable.punta\_manana} (90 y 100~s) y seis en
\texttt{laborable.mediodia} (60, 58, 62, 64, 59 y 61~s). Las ocho juntas, ordenadas, son 58, 59, 60, 61, 62, 64, 90,
100, y su mediana baja es el elemento $\lfloor 7/2\rfloor=3$: $m_{AB}=61$. Entonces
\begin{align*}
\tau_{AB,\texttt{punta\_manana}}&=\redondeo\!\left(\tfrac{2\cdot90+5\cdot61}{7}\right)=\redondeo(69{,}29)=69,\\
\tau_{AB,\texttt{mediodia}}&=\redondeo\!\left(\tfrac{6\cdot60+5\cdot61}{11}\right)=\redondeo(60{,}45)=60,
\end{align*}
y $\tau_{AB,\texttt{all}}=61$ porque hay dos franjas. Sin contracción, la punta de la mañana valdría 90~s, apoyada solo
en dos autobuses. Un tramo $C{>}D$ con tres pistas, todas de noche (200, 210 y 190~s), guarda solo
\texttt{laborable.noche} $=200$ (la contracción hacia sí mismo no cambia nada) y ningún valor \texttt{all}.
\end{ejemplo}

\subsection{Consulta}

Dado un tramo $\kappa$ y una franja $b$, el tiempo aprendido es
\begin{equation}
\label{eq:consulta-tramo}
\tau(\kappa,b)=\begin{cases}
\tau_{\kappa,b} & \text{si existe},\\
\tau_{\kappa,\texttt{all}} & \text{si no, y si existe},\\
\text{indefinido} & \text{en otro caso.}
\end{cases}
\end{equation}

Se dice que el tramo está \emph{aprendido} en la franja $b$ si $\tau(\kappa,b)$ está definido. Obsérvese que un tramo
con datos solo de noche está aprendido de noche y no lo está a mediodía: entonces se usa el respaldo de la sección
siguiente, no el valor nocturno.

\section{Respaldo por Longitud del Tramo}
\label{sec:respaldo-largo}

Un tramo sin tiempo aprendido no se puede dejar sin estimar, porque la llegada suma todos los tramos del camino. La
primera idea fue usar el horario multiplicado por un factor de ritmo de la línea (sección «\nameref{sec:k}»), pero eso
reproduce el error de forma del horario interpolado. La observación que lleva al respaldo es que ese error depende
sobre todo de la \emph{longitud} del tramo: los tramos cortos son urbanos (el autobús tarda lo que dice el horario o
algo más), y los largos son interurbanos y de autopista (el autobús tarda bastante menos de lo programado, porque la
interpolación reparte el tiempo a velocidad constante).

\subsection{Cubos de Longitud}

La longitud de un tramo es $L=s_{i+1}-s_i$, a partir de \texttt{shape\_dist\_traveled}. Se clasifica en seis cubos
(\texttt{LENGTH\_CUTS\_METERS}):
\begin{equation}
\label{eq:cubo}
c(L)=\begin{cases}
0 & L<300\\
1 & 300\le L<600\\
2 & 600\le L<1000\\
3 & 1000\le L<1500\\
4 & 1500\le L<3000\\
5 & L\ge3000
\end{cases}
\end{equation}
Los cortes son cerrados por abajo: un tramo de exactamente 300~m está en el cubo 1.

\subsection{Proporción Real/Horario por Cubo y Franja}

Para cada muestra de tramo con duración programada $h=a_{i+1}-a_i>0$ y longitud conocida $L>0$, se calcula la
proporción $q=y/h$. Para cada cubo $c$ y grupo $g$ (franja o \texttt{all}) se toma la mediana baja de \emph{todas} las
proporciones del grupo (aquí sin reducir a un valor por pista, porque cada proporción ya es adimensional y compara
tramos distintos):
\begin{equation}
\label{eq:proporcion}
r_{c,g}=\redondeo_3\big(\mediana\{q\}\big),
\end{equation}
siempre que el grupo tenga al menos 20 proporciones (\texttt{MIN\_SAMPLES\_PER\_LENGTH\_RATIO}); si no, $r_{c,g}$ no
existe. Aquí $\redondeo_3(x)=\redondeo(1000x)/1000$ redondea a tres decimales con la misma regla de las mitades. La longitud
$L$ de un tramo, al calibrar, es la mediana baja de las longitudes $s_{i+1}-s_i$ de ese par de paradas en todos los
viajes del horario que lo recorren (puede variar algo entre viajes si los trazados difieren).

\subsection{Consulta}

Al predecir, para un tramo no aprendido con duración programada $h$ y longitud $L$ del viaje concreto,
\begin{equation}
\label{eq:respaldo}
\tau^{\mathrm{L}}(h,L,b)=\clamp\big(\hat r,\ 0{,}3,\ 1{,}5\big)\cdot h,\qquad
\hat r=\begin{cases}
r_{c(L),b} & \text{si existe},\\
r_{c(L),\texttt{all}} & \text{si no.}
\end{cases}
\end{equation}
Si no existe ninguna de las dos proporciones, si $L$ no se conoce o si $L\le0$, $\tau^{\mathrm{L}}$ está indefinido.
Las cotas $[0{,}3;\,1{,}5]$ impiden que un dato corrupto (una proporción de 0,1 o de 2) se traslade a la predicción.
Si el viaje no tiene distancias (sin \texttt{shape\_dist\_traveled}), $L$ no se conoce y este respaldo no se aplica.

La Tabla~\ref{tab:proporciones} muestra las proporciones calibradas el 31 de octubre de 2026. En la columna
\texttt{all} la tendencia es clara y monótona: los tramos de menos de 300~m se recorren en el tiempo programado y los
de más de 3~km en algo más de la mitad. En las franjas, los cubos 3 y 4 quedan casi iguales.

\begin{table}[ht]
\centering
\caption{Proporción Entre Tiempo Real y Tiempo Programado por Cubo de Longitud}\label{tab:proporciones}
\begin{tabular}{@{}clcccc@{}}
\toprule
Cubo & Longitud & $r_{c,\texttt{all}}$ & \shortstack{Laborable\\punta mañana} & \shortstack{Laborable\\mediodía} & \shortstack{Fin de semana\\mañana}\\
\midrule
0 & $<300$~m        & 0,987 & 1,016 & 1,052 & 0,894\\
1 & 300--600~m      & 0,907 & 0,937 & 0,949 & 0,823\\
2 & 600--1\,000~m   & 0,769 & 0,811 & 0,793 & 0,716\\
3 & 1\,000--1\,500~m & 0,647 & 0,665 & 0,659 & 0,606\\
4 & 1\,500--3\,000~m & 0,640 & 0,668 & 0,669 & 0,600\\
5 & $\ge3\,000$~m    & 0,534 & 0,565 & 0,524 & 0,523\\
\bottomrule
\end{tabular}
\nota{Columna $r_{c,\texttt{all}}$: grupo \texttt{all}; las demás, la franja de su título. Muestras acumuladas a 31
de octubre de 2026.}
\end{table}

\section{Factor de Ritmo}
\label{sec:k}

El factor de ritmo $k$ es el último recurso de la cascada y, además, el único modelo cuando no existe ninguna tabla de
tiempos por tramo (por ejemplo, al poner el método en marcha en una red nueva). Entonces las capas 1 y 2 usan
$E=k\,(t_1-t_0)$, la capa 3 no se aplica (las filas sin datos del autobús salen con el horario publicado) y el autobús
que espera en la cabecera se calcula con la ecuación~\eqref{eq:eta-capa1}. Mide cuánto tarda de verdad un autobús
frente a lo que dice el horario, de forma global por franja y, con datos suficientes, por línea y franja.

\subsection{Consulta}

\begin{equation}
\label{eq:k}
k(\ell,D,t)=\clamp\big(\hat k,\ 0{,}6,\ 1{,}2\big),\qquad
\hat k=\begin{cases}
k_{\ell,b} & \text{si existe},\\
k_b & \text{si no, y si existe},\\
1 & \text{en otro caso,}
\end{cases}\qquad b=\franja(D,t),
\end{equation}
donde $\ell$ es la línea del viaje (\texttt{MIN\_FACTOR}, \texttt{MAX\_FACTOR}). El número de línea se escribe
siempre como en el horario, sin ceros a la izquierda: el feed pone \texttt{A0651} y se lee como \texttt{651}, y una
tabla guardada con la otra forma vale igual (\texttt{PaceFactors::lineKey}).

\subsection{Estimación}

Las \emph{muestras de ritmo} (capítulo «\nameref{cap:aprendizaje}») son tuplas $(j,\ \ell,\ v,\ b,\ H,\ Y)$: la pista $j$,
la línea $\ell$, el vehículo $v$, la franja $b$ del instante de la observación, el tiempo que según el horario faltaba
para llegar a una parada, $H=a_\dest-\Tpos\ge120$~s, y el tiempo que de verdad faltaba, $Y=t_{\mathrm{real}}-\tgps>0$.

Para un conjunto de pistas $\mathcal J$ (las de una franja, o las de una línea en una franja), se calcula para cada
pista la proporción de sumas
\begin{equation}
\rho_j=\frac{\sum Y}{\sum H}\qquad\text{(sumas sobre las muestras de la pista $j$)},
\end{equation}
que pondera cada muestra por su duración (una predicción a 15 minutos pesa más que una a 2), y después
\begin{equation}
\label{eq:k-contraido}
\boxed{\;k=\redondeo_3\!\left(\clamp\!\Big(1+\big(\tilde\rho-1\big)\frac{m}{m+10},\ 0{,}6,\ 1{,}2\Big)\right),\qquad
\tilde\rho=\mediana_j\,\rho_j,\quad m=|\mathcal J|\;}
\end{equation}
(\texttt{SHRINKAGE\_TRACKS}). Es una contracción hacia $k=1$ (el horario tal cual) con peso equivalente a diez pistas:
la misma idea que la sección «\nameref{sec:contraccion-tramo}», con el horario como valor común.

Se estima $k_b$ si la franja tiene al menos $m\ge8$ pistas (\texttt{MIN\_TRACKS\_PER\_BAND}), y $k_{\ell,b}$ si la
línea en esa franja tiene al menos 25 pistas y 5 vehículos distintos (\texttt{MIN\_TRACKS\_PER\_LINE},
\texttt{MIN\_VEHICLES\_PER\_LINE}). La exigencia de vehículos distintos evita que el factor de una línea refleje a un
solo conductor o a un autobús averiado.

\begin{ejemplo}
Ocho pistas con proporciones 0,80, 0,82, \dots, 0,94 (en pasos de 0,02) tienen mediana baja $\tilde\rho=0{,}86$ (el
elemento $\lfloor3{,}5\rfloor=3$). Entonces $k=1-0{,}14\cdot8/18=0{,}93778$, que se redondea a $0{,}938$. Con 40
pistas de proporción 0,85, $k=1-0{,}15\cdot40/50=0{,}88$. Con diez pistas de proporciones 1,6, 1,7, \dots, 2,5, la mediana baja es $2{,}0$ y la fórmula
da $1+1{,}0\cdot10/20=1{,}5$, que queda en la cota: $k=1{,}2$.
\end{ejemplo}

Con los datos de octubre de 2026 los factores globales van de $0{,}756$ (fin de semana, punta de la mañana) a $0{,}928$
(laborable, punta de la tarde), y 54 líneas tienen factor propio en alguna franja.

\section{Tiempo Esperado Entre Dos Instantes del Horario}
\label{sec:tiempo-esperado}

Con los tres niveles anteriores se define la función central del método: el tiempo real que se espera que tarde el
autobús en ir desde el punto donde el horario lo pone a la hora $t_0$ hasta el punto donde lo pone a la hora $t_1$.
Los dos extremos se dan como \emph{horas del horario}, no como distancias, porque es lo que proporciona la
localización ($\Tpos$) y lo que define a cada parada ($a_\dest$).

\begin{definicion}[Tiempo de un tramo]
Para el tramo $i$ de $J$, con duración programada $h_i=a_{i+1}-a_i$, longitud $L_i=s_{i+1}-s_i$ (si se conoce),
identificador $\kappa_i$, franja $b$ y factor $k$:
\begin{equation}
\label{eq:cascada}
\tau_i=\begin{cases}
\tau(\kappa_i,b) & \text{si está definido (tiempo aprendido)},\\
\tau^{\mathrm{L}}(h_i,L_i,b) & \text{si no, y si está definido (respaldo por longitud)},\\
k\cdot h_i & \text{en otro caso (factor de ritmo).}
\end{cases}
\end{equation}
\end{definicion}

\begin{definicion}[Tiempo esperado]
\label{def:tiempo-esperado}
Sea $J$ un viaje con llegadas $a_1\le\dots\le a_N$, y sean $t_0<t_1$ dos horas del horario, una franja $b$ y un
factor $k$. Si no existe ninguna tabla de tiempos por tramo (la tabla \texttt{segments} está vacía),
\[
E(J,t_0,t_1,b,k)=k\,(t_1-t_0).
\]
En otro caso,
\begin{equation}
\label{eq:tiempo-esperado}
\boxed{\;E(J,t_0,t_1,b,k)=\sum_{i\in I(t_0,t_1)}\max(0,w_i)\,\tau_i\;}
\end{equation}
con
\begin{equation}
\label{eq:pesos}
w_i=\begin{cases}\dfrac{\min(a_{i+1},t_1)-\max(a_i,t_0)}{a_{i+1}-a_i} & \text{si } a_{i+1}>a_i,\\[2mm] 1 & \text{si } a_{i+1}=a_i,\end{cases}
\end{equation}
donde el conjunto de tramos $I(t_0,t_1)$ se obtiene recorriendo $i=1,\dots,N-1$ en orden y aplicando dos reglas:
se \emph{salta} el tramo si $a_{i+1}\le t_0$ (ya está detrás) y se \emph{termina} el recorrido en el primer tramo con
$a_i\ge t_1$ (ya está después del destino).
\end{definicion}

El peso $w_i$ es la fracción de la duración programada del tramo $i$ que queda entre $t_0$ y $t_1$: vale 1 en los
tramos completos y una fracción en el tramo donde está el autobús. Equivale a suponer que, \emph{dentro de un tramo},
el autobús avanza al ritmo del horario; entre tramos, en cambio, el ritmo es el aprendido. Se probó una curva no lineal
dentro del tramo y no mejoró (sección «\nameref{sec:descartadas}»).

Varias propiedades de la definición son importantes para replicarla exactamente:

\begin{observacion}[Se recorre por horas, no por índices]
Las reglas de salto y de fin se evalúan sobre las horas $a_i$, no sobre los índices de parada. Si dos paradas
consecutivas $p_j$ y $p_{j+1}$ tienen la misma hora programada, el destino $t_1=a_{j+1}=a_j$ termina el recorrido en el
tramo $j$ (porque $a_j\ge t_1$), de modo que ambas paradas reciben la misma llegada estimada.
\end{observacion}

\begin{observacion}[Tramos de duración nula]
Un tramo con $a_{i+1}=a_i$ (frecuente en el GTFS, donde dos paradas próximas comparten minuto) se incluye con peso 1,
y aporta su tiempo aprendido completo, solo si $t_0<a_i$ y $a_i<t_1$. Por eso el horario corregido de la capa 3
(sección «\nameref{sec:capa3}») empieza en $t_0=d_1-1$ y no en $d_1$: así se incluyen los tramos de duración nula del
principio del viaje. Por la misma razón, $E$ no es exactamente aditiva: si $t_0<t_m<t_1$ y hay un tramo de duración
nula con $a_i=t_m$, $E(t_0,t_1)$ lo cuenta y $E(t_0,t_m)+E(t_m,t_1)$ no.
\end{observacion}

\begin{observacion}[La espera en la parada]
La duración de un tramo, $h_i=a_{i+1}-a_i$, incluye la espera programada en $p_i$, mientras que $\Tpos$ se interpola
entre $d_i$ y $a_{i+1}$ (sección «\nameref{sec:trazado}»). Si $d_i>a_i$, un autobús que acaba de salir de $p_i$
($\Tpos=d_i$) recibe $w_i=(a_{i+1}-d_i)/h_i<1$ y no cuenta la parte del tiempo aprendido que corresponde a esa
espera. En la base de Bizkaibus todas las paradas tienen $d_i=a_i+2$~s, así que el efecto es de unos segundos; los
vectores de prueba, con $a_i=d_i$, no lo ejercitan.
\end{observacion}

\begin{observacion}[Franja única]
Todo el camino se evalúa con una sola franja $b$, la del instante de la observación, aunque el autobús vaya a cruzar a
otra franja antes de llegar. Calcular la franja de cada tramo a la hora prevista de paso por él se probó y no mejoró
(sección «\nameref{sec:descartadas}»).
\end{observacion}

\begin{ejemplo}
\label{ej:tiempo-esperado}
En el viaje sintético del ejemplo~\ref{ej:posicion}, con franja \texttt{laborable.punta\_manana} y $k=0{,}9$, los
tramos aprendidos son $S0{>}S1$ (110~s en esa franja) y $S1{>}S2$ (140~s), y las proporciones del cubo~2 y del cubo~4
en esa franja son $0{,}811$ y $0{,}668$. Desde $t_0=\Tpos=28\,950$ (a mitad del tramo $S1{>}S2$):
\begin{itemize}
  \item hasta $S2$ ($t_1=29\,010$): $w=(29\,010-28\,950)/120=0{,}5$ en $S1{>}S2$, luego $E=0{,}5\cdot140=70$~s;
  \item hasta $S3$ ($t_1=29\,310$): se añade $S2{>}S3$, no aprendido, de 1\,500~m (cubo 4) y 300~s programados:
        $0{,}668\cdot300=200{,}4$~s, luego $E=270{,}4$~s;
  \item hasta $S4$ ($t_1=29\,430$): se añade $S3{>}S4$, de 700~m (cubo 2) y 120~s: $0{,}811\cdot120=97{,}32$~s, luego
        $E=367{,}72$~s.
\end{itemize}
Sin distancias (viaje sin trazado), los dos últimos tramos caen al factor de ritmo: $0{,}9\cdot300=270$ y
$0{,}9\cdot120=108$~s, y $E$ vale $70$, $340$ y $448$~s.
\end{ejemplo}

\section{La Cascada Completa}
\label{sec:cascada}

La Tabla~\ref{tab:cascada} resume qué nivel del modelo se usa en cada caso. Todos los niveles devuelven un tiempo en
segundos para un tramo; el tiempo esperado de un camino es la suma ponderada de la
ecuación~\eqref{eq:tiempo-esperado}, y cada tramo del camino puede salir de un nivel distinto.

\begin{table}[ht]
\centering
\caption{Cascada de Estimación del Tiempo de un Tramo}\label{tab:cascada}
\begin{tabularx}{\textwidth}{@{}c>{\raggedright\arraybackslash}X>{\raggedright\arraybackslash}X>{\raggedright\arraybackslash}X@{}}
\toprule
Nivel & Se usa si & Valor & Aprende de\\
\midrule
1 & El tramo tiene valor en la franja & $\tau_{\kappa,b}$, contraído hacia el tramo & Pasos reales de $\ge2$ pistas en la franja\\
2 & El tramo tiene valor a cualquier hora (y muestras de $\ge2$ franjas) & $\tau_{\kappa,\texttt{all}}$ & Pasos reales de $\ge2$ pistas\\
3 & Se conoce la longitud y hay proporción para su cubo & $\clamp(\hat r,0{,}3,1{,}5)\cdot h$, con $\hat r=r_{c,b}$ o, si no, $r_{c,\texttt{all}}$ & $\ge20$ proporciones del cubo en el grupo\\
4 & Siempre & $k\cdot h$, con $k\in[0{,}6;\,1{,}2]$ & $\ge8$ pistas (franja) o $\ge25$ y 5 vehículos (línea)\\
\bottomrule
\end{tabularx}
\end{table}

El diseño en cascada responde a un compromiso entre sesgo y varianza: cada nivel es menos específico que el anterior
(más sesgo) pero se apoya en más datos (menos varianza). El nivel que se usa es siempre el más específico que tiene
datos suficientes, y los umbrales de «suficiente» se han fijado por validación cruzada (capítulos «\nameref{cap:aprendizaje}»
y «\nameref{cap:evaluacion}»). Con los datos de octubre de 2026, la tabla tiene 2\,454 tramos aprendidos, el 86\,\% de las
predicciones tiene todos los tramos de su camino en los niveles 1 o 2, y, como los seis cubos tienen proporción a cualquier hora, el nivel 4 solo actúa en la práctica
en viajes sin distancias a lo largo del trazado.

\chapter{Predicción en Tres Capas}
\label{cap:prediccion}

Este capítulo combina la localización (capítulo «\nameref{cap:localizacion}») y el modelo de tiempos
(capítulo «\nameref{cap:modelo}») en la predicción que ve el usuario. La idea organizadora es una jerarquía de fuentes,
cada una con una regla explícita para abstenerse:

\begin{description}
  \item[Capa 1 · GPS.] La posición proyectada sobre el recorrido. Es la fuente más precisa y la que permite afirmar un
        adelanto o un retraso.
  \item[Capa 2 · Parada registrada.] Cuando el GPS no es utilizable pero el feed dice cuál es la próxima parada y a qué
        hora se tomó el dato, se supone al autobús a mitad del tramo que lleva a esa parada.
  \item[Capa 3 · Horario corregido.] Cuando no hay ningún dato del vehículo, se recalcula el horario desde la cabecera
        con los tiempos aprendidos.
\end{description}

\noindent Si ninguna capa puede sostener una estimación, se muestra el horario publicado, marcado como tal. Ninguna
capa inventa: cada una se usa solo en las condiciones en que se ha medido que mejora a la siguiente.

\section{Entradas, Salidas y Estados}
\label{sec:estados}

La unidad de trabajo es una \emph{fila}: una salida concreta de una parada. Una fila identifica un viaje $J$ (su
línea, número de viaje y salida de cabecera $d_1$) y una parada de ese viaje, la parada de \emph{destino} de la
predicción, con su orden $q$ y su hora publicada $a_\dest$. Para cada fila el motor devuelve:

\begin{itemize}
  \item un \textbf{estado}, de la Tabla~\ref{tab:estados};
  \item una \textbf{llegada estimada} $\ETA$, en segundos desde medianoche;
  \item un \textbf{retraso} $\delta$ en segundos, que puede ser \emph{nulo} (no se afirma retraso alguno);
  \item el vehículo y su parada actual, si se conocen.
\end{itemize}

\begin{table}[ht]
\centering
\caption{Estados de una Fila y lo que Muestra la Aplicación}\label{tab:estados}
\begin{tabularx}{\textwidth}{@{}l>{\raggedright\arraybackslash}X>{\raggedright\arraybackslash}X@{}}
\toprule
Estado & Significado & Texto en la aplicación\\
\midrule
\texttt{live} & Capa 1: el GPS sitúa al autobús en el recorrido. & «En hora», «Retraso $+N$m» o «Adelantado $-N$m»; distintivo «En ruta», «Llegando» o «En la parada» (también mientras espera en la cabecera). \\
\texttt{estimated} & Capa 2 o capa 3: estimación sin seguimiento en directo; retraso $0$. & «Estimado con lo que tardan de verdad los buses»; distintivo «Estimado». \\
\texttt{departed} & El autobús ya pasó por la parada. & «Ya salió»; distintivo «Saliendo». \\
\texttt{likely\_departed} & Sin datos del autobús, el horario corregido indica que ya debió de pasar; o, con GPS, la hora estimada de llegada pasó hace más de 3~min sin que el feed ni el GPS confirmen el paso. & «Sin datos del bus: puede que ya haya pasado»; distintivo «Sin datos del bus»; se ordena al final. \\
\texttt{scheduled} & Ninguna capa puede sostener una estimación. & «Horario, sin seguimiento en vivo»; distintivo «Solo horario». Si el feed trae autobuses con hora de hoy y todavía no es la hora del viaje, la ficha dice que aún no tiene autobús asignado. \\
\bottomrule
\end{tabularx}
\end{table}

La fila se empareja con una entrada del feed como se explicó en la sección «\nameref{sec:emparejamiento}»: primero de forma
exacta por los alias de viaje y, si no, por línea, número de viaje y salida de cabecera con una tolerancia de 180~s.
Si no hay entrada, la fila se resuelve sin datos del vehículo: con el registro de pasos (sección
«\nameref{sec:registro-pasos}») o, si este no la da por pasada, con la capa~3.

En todo el capítulo, $D$ es la fecha de servicio del viaje (de madrugada, la de ayer para los viajes que pasan de las 24:00; sección «\nameref{sec:mas-de-24h}»), $\ell$ la línea, $n$ el orden de la próxima parada que declara el
feed y $E(\cdot)$ el tiempo esperado de la definición~\ref{def:tiempo-esperado} aplicada al viaje $J$.

\section{Capa 1: Llegada Desde el GPS}
\label{sec:capa1}

Si existe la posición $(\Tpos,\tgps)$ (sección «\nameref{sec:condiciones-gps}»), la franja y el factor se toman en el
instante de la posición,
\[
b=\franja(D,\tgps),\qquad k=k(\ell,D,\tgps),
\]
y la llegada a la parada de destino es
\begin{equation}
\label{eq:eta-capa1}
\boxed{\;\ETA=\redondeo\big(\tgps+E(J,\Tpos,a_\dest,b,k)\big),\qquad \delta=\ETA-a_\dest\qquad(a_\dest\ge\Tpos).\;}
\end{equation}
Si $a_\dest<\Tpos$, la parada ya ha quedado atrás según el horario y
\begin{equation}
\label{eq:paso-capa1}
\ETA=\tgps-(\Tpos-a_\dest),\qquad\delta=\delta_{\mathrm{pos}}=\tgps-\Tpos,
\end{equation}
que es la hora a la que se estima que pasó: la distancia programada entre la parada y el autobús, contada hacia atrás
desde la hora de la posición.

La llegada parte de $\tgps$, no de $\tnow$. Cada posición llega al motor un minuto o más después de tomarse (66~s
de mediana y 92~s en el percentil 90 el 5 de noviembre de 2026; sección «\nameref{sec:llegando}»); sumar el tiempo
restante a $\tnow$ retrasaría todas las llegadas estimadas en lo que tarda en recibirse el dato. Partiendo de $\tgps$, la antigüedad
del dato no sesga la predicción: solo la hace algo menos precisa. Una consecuencia es que $\ETA$ puede ser anterior a $\tnow$ si el autobús estaba muy cerca
de la parada cuando se tomó la posición.

\section{Capa 1 en la Cabecera}
\label{sec:cabecera-eta}

Con la posición en la cabecera~\eqref{eq:cabecera}, la ecuación~\eqref{eq:eta-capa1} daría $d_1+E(J,a_1,a_\dest,b,k)$,
lo que supone que el autobús sale exactamente a su hora. Los autobuses no salen así, y la medida lo confirmó: con esa
suposición, las llegadas a 10--20~min vista se adelantaban 1~min de media y los aciertos dentro de $\pm2$~min bajaban
del 67{,}0 al 64{,}0\,\% (sección «\nameref{sec:eval-trazados}»). La llegada es, por tanto, la de la capa 3,
\begin{equation}
\label{eq:eta-cabecera}
\ETA=\max\big(C(J,a_\dest),\ \tnow\big),\qquad \delta=\ETA-a_\dest\qquad(a_\dest>d_1),
\end{equation}
con $C$ el horario corregido~\eqref{eq:capa3}, cuyos 60~s iniciales y cuya corrección por viaje
$\Delta_v$~\eqref{eq:correccion-viaje} recogen lo tarde que salen de verdad los autobuses. La fila es «en directo»
porque se sabe que el autobús está allí; si no hay tiempos aprendidos, se usa la ecuación~\eqref{eq:eta-capa1}. La
parada de cabecera ($a_\dest\le d_1$) tiene como llegada la hora de salida, $d_1$.

\section{Capa 2: Llegada Desde la Parada Registrada}
\label{sec:capa2}

Muchas entradas del feed traen la próxima parada y la hora del registro, pero no una posición utilizable: sin
coordenadas, con un GPS fuera del pasillo o con una hora de posición que no es de hoy. La capa 2 aprovecha lo que sí
es un dato real: el autobús ha dejado atrás la parada $p_{n-1}$ y va hacia $p_n$, y eso se sabía a la hora $\tref$.

\subsection{El Registro}

El \emph{registro} de una entrada existe solo si se cumplen todas estas condiciones:

\begin{enumerate}
  \item la entrada tiene viaje y orden $n$, y el par (viaje, vehículo) no está marcado como no fiable (sección
        «\nameref{sec:desconfianza}»);
  \item $n\ge2$: hay una parada anterior;
  \item hay hora de referencia $\tref$ ($\tgps$ si existe y, si no, $t_{\mathrm{rec}}$) y es reciente:
        $\tnow-\tref\le600$~s;
  \item el viaje ha empezado: $\tref\ge d_1-60$~s;
  \item el viaje tiene paradas de orden $n-1$ y $n$;
  \item el tramo es corto: $a_n-a_{n-1}\le1\,200$~s (\texttt{REGISTRATION\_MAX\_SEGMENT\_SECONDS}).
\end{enumerate}

\noindent Entonces, con $b=\franja(D,\tref)$ y $k=k(\ell,D,\tref)$,
\begin{align}
\label{eq:registro}
\theta&=\redondeo\Big(\tfrac12\,E(J,a_{n-1},a_n,b,k)\Big) &&\text{(lo que le falta para llegar a $p_n$)},\\
\delta_{\mathrm{reg}}&=\tref+\theta-a_n &&\text{(retraso estimado en $p_n$)},\\
g&=a_n-a_{n-1} &&\text{(duración programada del tramo)}.
\end{align}

\noindent Si el tramo de $p_{n-1}$ a $p_n$ tiene duración programada nula ($a_{n-1}=a_n$, frecuente cuando dos paradas
comparten minuto), $E(J,a_{n-1},a_n,b,k)=0$ por la regla de los tramos de duración nula (sección
«\nameref{sec:tiempo-esperado}»), aunque ese tramo tenga tiempo aprendido. Entonces $\theta=0$ y $g=0$: el autobús se
supone ya en $p_n$, y el contraste de la sección «\nameref{sec:desconfianza}» se hace con
$\delta_{\mathrm{reg}}=\tref-a_n$.

\subsubsection{Por Qué la Mitad} El feed no dice en qué punto del tramo está el autobús. Si el instante de la observación
es independiente de la posición, el autobús está en cualquier punto del tramo con la misma probabilidad, y el valor
esperado del tiempo que le falta es la mitad del tiempo del tramo. Se midieron repartos de 0,35 a 0,65 y el óptimo fue
0,5 (sección «\nameref{sec:descartadas}»).

\subsubsection{Por Qué Tramos de Hasta 20 Minutos} En tramos largos (salidas de cabecera, regulaciones) la incertidumbre
de «a mitad del tramo» es de muchos minutos. Contrastada con el GPS en 201 autobuses, la estimación de la parada
registrada se desvía 0,7~min de mediana y 1,7~min en el percentil 90 en tramos de hasta 5~min, pero 10~min de mediana
en tramos de más de 10~min. Con el umbral de 20~min, la capa 2 mejora al horario en todas las longitudes de tramo en que
se usa; por encima de ese umbral, no.

\subsection{La Llegada}

Para una parada de destino con orden $q\ge n$,
\begin{equation}
\label{eq:eta-capa2}
\boxed{\;\ETA=\redondeo\big(\tref+\theta+E(J,a_n,a_\dest,b,k)\big)\;}
\end{equation}
donde el último término es cero si $a_\dest\le a_n$. Para $q<n$, el registro no da llegada (la parada ya ha quedado
atrás; véase la sección «\nameref{sec:paso}»).

\section{Contraste Entre el GPS y la Parada Registrada}
\label{sec:desconfianza}

Cuando una entrada tiene a la vez posición y registro, las dos fuentes deberían coincidir. Si no coinciden, una de
ellas está equivocada (un GPS que pertenece a otro vehículo, un orden de parada que no se actualizó) y no
hay forma de saber cuál. La regla es que el par (viaje, vehículo) deja de ser fiable si

\begin{equation}
\label{eq:contraste}
g\le360\ \text{s}\qquad\text{y}\qquad\big|\delta_{\mathrm{pos}}-\delta_{\mathrm{reg}}\big|>360\ \text{s}.
\end{equation}
Los dos umbrales son las constantes \texttt{REGISTRATION\_CROSSCHECK\_MAX\_SEGMENT\_SECONDS} y
\texttt{MAX\_GPS\_VS\_STOP\_DISAGREEMENT\_SECONDS}. Un par no
fiable no tiene posición ni registro durante el resto de la petición, de modo que sus filas pasan a la capa 3. El
contraste solo se hace en tramos de hasta 6~minutos, porque en los largos el registro es demasiado impreciso para
contradecir al GPS.

Para que el resultado sea reproducible, la posición (con su contraste) debe evaluarse antes que cualquier uso del
registro. En el código, cada consulta del estado de una fila empieza por la posición; el registro que se usa dentro del
contraste se calcula antes de marcar el par y, una vez marcado, todas las consultas posteriores lo tratan como
inexistente.

\begin{ejemplo}
En el viaje sintético, un autobús declara $n=3$ (va hacia $S2$, a 1\,000~m), pero su GPS está a 2\,800~m, entre $S3$ y
$S4$, a las $29\,400$~s. La proyección elige el candidato «más allá» (a 11~m del trazado, frente a 1\,800~m del mejor
«esperado») y da $\Tpos=29\,361$, luego $\delta_{\mathrm{pos}}=39$~s. El registro da $\theta=70$ y
$\delta_{\mathrm{reg}}=29\,400+70-29\,010=460$~s. Como $g=120\le360$ y $|39-460|=421>360$, el par se marca como no
fiable y todas las filas de ese viaje se resuelven sin datos del vehículo.
\end{ejemplo}

\section{Capa 3: Horario Corregido}
\label{sec:capa3}

Sin ningún dato fiable del vehículo, la mejor estimación de su llegada es la de un autobús que sale de la cabecera poco
después de su hora y recorre cada tramo en el tiempo aprendido. Si la red tiene tabla de tiempos por tramo, el viaje
tiene paradas y $a_\dest>d_1$, se toman la franja y el factor a mitad de camino,
\[
t_{\mathrm{mid}}=\Big\lfloor\frac{d_1+a_\dest}{2}\Big\rfloor,\qquad b=\franja(D',t_{\mathrm{mid}}),\qquad
k=k(\ell,D',t_{\mathrm{mid}}),
\]
donde $D'$ es la fecha de servicio $D$, salvo en un festivo entre semana, en que es el domingo de esa semana. Las
franjas solo dependen del día de la semana, así que en un festivo se usan las de fin de semana. El \emph{horario
corregido} es
\begin{equation}
\label{eq:capa3}
\boxed{\;C=d_1+60+\redondeo\big(E(J,\,d_1-1,\,a_\dest,\,b,\,k)\big)+\redondeo(\Delta_v).\;}
\end{equation}
Los $60$~s (\texttt{SCHEDULE\_ORIGIN\_OFFSET\_SECONDS}) recogen que los autobuses salen de la cabecera algo después de
la hora publicada; el resumen de cada captura vigila que este sesgo siga siendo el medido. El inicio en $d_1-1$ incluye
los tramos de duración nula del principio del viaje (sección «\nameref{sec:tiempo-esperado}»). El último término,
$\Delta_v$, es la corrección por viaje de la sección «\nameref{sec:correccion-viaje}»; sin tabla de correcciones vale 0.

La decisión es:

\begin{enumerate}
  \item Si $C\ge\tnow$: estado \texttt{estimated}, $\ETA=C$, $\delta=0$. El autobús todavía no debería haber pasado y
        $C$ es una estimación mejor que el horario publicado.
  \item Si no, y además
        \begin{equation}
        \label{eq:probablemente}
        C<\tnow-600,\qquad a_\dest-C\ge600,\qquad \lambda\ge0{,}8
        \end{equation}
        (constantes \texttt{WRONG\_SCHEDULE\_PASSED\_SECONDS}, \texttt{WRONG\_SCHEDULE\_GAP\_SECONDS} y
        \texttt{WRONG\_SCHEDULE\_LEARNED\_SHARE}): estado \texttt{likely\_departed}, $\ETA=a_\dest$, $\delta=0$.
  \item En otro caso, estado \texttt{scheduled}, $\ETA=a_\dest$, $\delta=0$.
\end{enumerate}

\noindent La \emph{parte aprendida} $\lambda$ del camino es la fracción de su duración programada que recorre tramos
aprendidos en la franja $b$:
\begin{equation}
\label{eq:parte-aprendida}
\lambda=\frac{\sum_{i\in I(d_1-1,\,a_\dest)}\mathbf{1}\big[\tau(\kappa_i,b)\text{ definido}\big]\,h_i}{\sum_{i\in I(d_1-1,\,a_\dest)}h_i},
\qquad\lambda=0\ \text{si el denominador es }0,
\end{equation}
con el mismo conjunto de tramos $I$ de la definición~\ref{def:tiempo-esperado} y las duraciones programadas completas
$h_i$ (sin los pesos $w_i$).

Las tres condiciones de~\eqref{eq:probablemente}, que llevan al estado \texttt{likely\_departed}, identifican un caso
concreto y frecuente: un horario publicado mucho más lento que la realidad (típicamente, un viaje con un tramo largo de
autopista), en el que el horario corregido dice que el autobús pasó hace más de diez minutos, la hora publicada queda al
menos diez minutos después de esa hora corregida (a menudo, todavía no ha llegado), y la corrección se apoya en tramos medidos en al menos el 80\,\% del camino. Mostrar la
hora publicada como si el autobús fuera a venir haría que alguien esperase un autobús que ya se fue. No se oculta
porque no hay ningún dato del vehículo que lo confirme: se muestra al final de la lista, con su hora publicada y el
texto «Sin datos del bus: puede que ya haya pasado», y las alertas de llegada lo tratan como pasado. El texto dice lo que
se sabe: en las capturas del 3 de noviembre de 2026, el 97{,}8\,\% de los viajes con este estado no apareció en el feed
en todo el día (sección «\nameref{sec:fiabilidad-estados}»).

\subsection{Corrección por Viaje}
\label{sec:correccion-viaje}

Sin la corrección, el horario corregido llegaba de media 131~s antes que el autobús (sección
«\nameref{sec:eval-correccion}»). Hay dos razones. La suma de las medianas de los tramos queda por debajo del tiempo
medio del recorrido, porque el tiempo de cada tramo tiene una cola larga hacia arriba. Además, algunos viajes salen tarde
o se retrasan en el mismo sitio casi todos los días. La corrección $\Delta_v$ recoge cuánto se aparta habitualmente cada
viaje del horario corregido.

Un viaje se reconoce por la clave $v=(\ell,\,p_1,\,p_N,\,d_1,\,N)$: su línea, su primera y su última parada, su
salida de cabecera y su número de paradas. A diferencia del identificador del GTFS, esta clave no cambia cuando el
operador publica un horario nuevo con los mismos viajes. Como en el factor de ritmo (sección «\nameref{sec:k}»), la
línea, tanto en la clave como en la tabla por línea, se compara sin ceros a la izquierda. El tipo de día $\theta$ es
laborable o fin de semana según $D'$, de modo que un festivo entre semana cuenta como fin de semana. La parada de destino, en la posición $i$ del
recorrido (0 es la cabecera), cae en el tercio
\begin{equation}
\label{eq:tercio}
p=\min\big(2,\ \lfloor 3i/N\rfloor\big).
\end{equation}
La corrección se busca de lo más fino a lo más general:
\begin{equation}
\label{eq:correccion-viaje}
\Delta_v=\begin{cases}
\Delta(v,\theta,p) & \text{si el viaje tiene valor para ese tipo de día},\\
l(\ell,b_1) & \text{si no, y su línea lo tiene en la franja } b_1=\franja(D',d_1),\\
g(b_1) & \text{si no, y la franja lo tiene},\\
0 & \text{en otro caso.}
\end{cases}
\end{equation}
Las tres tablas se aprenden cada día (sección «\nameref{sec:recalibracion}»). Un día es festivo entre semana si el
horario hace circular ese día, por excepción, algún calendario de solo domingo (máscara de días $64$), que es como
publica Bizkaibus sus festivos. La corrección solo afecta a la capa 3: las capas 1 y 2 parten de donde está el autobús
y ya recogen lo que le pasa ese día.

\subsection{Sin Conexión}
\label{sec:sin-conexion}

La aplicación puede guardar en el móvil el horario completo de cada red para usarlo sin cobertura. En Bizkaibus, cada
viaje lleva además su horario corregido~\eqref{eq:capa3}, ya con $\Delta_v$, calculado una vez para un laborable y otra
para un fin de semana. El móvil usa el de fin de semana en sábado, en domingo y en los festivos, que reconoce con la
misma regla de los calendarios de solo domingo. Aplica las reglas de estado del servidor, con la parte aprendida
$\lambda$ calculada sobre el recorrido entero: \texttt{estimated} si $C\ge\tnow$ y \texttt{likely\_departed} con las
condiciones~\eqref{eq:probablemente}. En 1\,200 paradas al azar de un laborable, un sábado y un festivo, las horas del
móvil y del servidor difieren como mucho 6~s, por el orden de los redondeos y por cómo se cuenta el primer tramo
cuando la salida de cabecera va unos segundos después de su llegada.

\section{Paradas Ya Pasadas}
\label{sec:paso}

Antes de calcular una llegada, se comprueba si el autobús ya pasó por la parada de destino. La hora de paso
$t_{\mathrm{paso}}$ se calcula así:

\begin{itemize}
  \item Con posición en la cabecera (sección «\nameref{sec:cabecera}»): el autobús aún no ha pasado por ninguna
        parada.
  \item Con otra posición: si $a_\dest<\Tpos$ y además el feed ya ha dejado atrás esa parada ($q<n$) o el GPS está a más
        de 150~m de ella (\texttt{PASSED\_DISTANCE\_METERS}), $t_{\mathrm{paso}}=\tgps-(\Tpos-a_\dest)$
        (ecuación~\eqref{eq:paso-capa1}); si no, pero $q<n$ (el feed ya ha dejado atrás esa parada aunque la proyección
        ponga al autobús antes de ella), $t_{\mathrm{paso}}=\tgps$; si no, no ha pasado.
  \item Sin posición, con registro y $q<n$: $t_{\mathrm{paso}}=\tref-\max(0,\ a_{n-1}-a_\dest)$.
\end{itemize}

\noindent Si hay hora de paso, la fila sale como \texttt{departed} con $\ETA=\min(t_{\mathrm{paso}},\tnow)$ y
$\delta=0$.

\section{Llegando sin Confirmar el Paso}
\label{sec:llegando}

El feed de Bizkaibus publica una versión cada 64~s de mediana (de 58 a 72~s) y cada posición llega 66~s después de
tomarse (mediana; 92~s en el percentil 90), según las capturas del 5 de noviembre de 2026 de 7:00 a 9:00. La posición
que se usa es, por tanto, de hace un minuto o más. Si el autobús va algo más lento de lo esperado, la llegada calculada
desde esa posición con la ecuación~\eqref{eq:eta-capa1} puede quedar ya en el pasado sin que el autobús haya llegado a
la parada. Antes, la fila salía en directo con esa hora ya pasada y la lista la retiraba al quedar atrás (sección
«\nameref{sec:retraso-mostrado}»): en esas dos horas, el autobús desaparecía de la lista antes de llegar en 170 de las
5\,041 llegadas a parada que el GPS ve a menos de 40~m (el 3{,}4\,\%).

Con posición y sin hora de paso, si $\ETA<\tnow$, la fila sale con $\ETA=\tnow$ (en la aplicación, «0~min») y un
retraso de al menos $\tnow-a_\dest$, hasta que el feed o el GPS confirmen el paso. Si $\tnow-\ETA$ pasa de 180~s
(\texttt{ARRIVING\_MAX\_SECONDS}) sin confirmación, la fila sale como \texttt{likely\_departed} («puede que ya haya
pasado») con $\ETA=\tnow$, en vez de desaparecer. Con las mismas capturas, el autobús desaparece antes de llegar en
22 de las 5\,041 llegadas (el 0{,}4\,\%; recorridos que pasan dos veces cerca de la parada o un feed que salta
paradas) y, a cambio, tarda algo más en dejar de mostrarse después de irse: 104~s de mediana y 152~s en el percentil
90, frente a 61 y 98~s. La ficha del viaje usa la misma regla de paso y la misma llegada. El estado \texttt{likely\_departed} tiene
así dos usos: en la capa 3 lleva la hora publicada y retraso 0; aquí, la hora actual y el retraso que da el GPS. Quien
use la API no debe tratar igual su llegada en los dos casos.

\section{Registro de Pasos}
\label{sec:registro-pasos}

Un viaje que ya terminó, o que ha dejado de aparecer en el feed, no tiene entrada, y sus filas irían a la capa 3 como
si el autobús no hubiera salido. Para evitarlo, un proceso externo guarda cada minuto, para cada referencia de viaje
del feed, el último orden de parada visto y su hora:
\[
\text{pasos}[D][\rho]=\big(n,\ t\big),\qquad t=t_{\mathrm{rec}}\ \text{si existe; si no, }\tgps,
\]
donde $t$ es la hora del registro en que el feed dio ese orden de parada (la capa 2 y la memoria del feed prefieren, en
cambio, $\tgps$) y $D$ es el día del viaje, no el del reloj: pasada la medianoche, un autobús de un viaje de ayer se
apunta en los pasos de ayer, con $t$ en las horas de su horario (sección «\nameref{sec:mas-de-24h}»). Hasta el 7 de
noviembre de 2026 se apuntaba en el día del reloj: el autobús de las 23:30 de ayer, visto a las 00:30, quedaba en los
pasos de hoy, y las paradas del viaje de las 23:30 de hoy, que tiene la misma referencia, salían esa noche como ya
pasadas hasta su salida. Los pasos de cada día caducan 36~horas después de terminar ese día. Al responder, y solo para
los viajes de las filas que no tienen entrada en el feed, se leen los pasos de todas sus referencias (un viaje
fusionado puede tener varias; sección «\nameref{sec:alias}») y se conserva el de mayor orden. Si la fila tiene $q<n$,
sale como \texttt{departed} con $\ETA=\min(t,\tnow)$. Esta comprobación se hace antes de recurrir a la capa 3.

\section{Autobuses que Faltan en una Lectura del Feed}
\label{sec:memoria-feed}

El feed deja fuera a veces, durante una lectura o unas pocas, a un autobús que sigue en su viaje. En las capturas del 2
al 4 de noviembre de 2026 pasa unas 700 veces al día con una sola lectura y otras 250 a 400 con dos o tres. Sin otra
medida, en esas lecturas las filas del viaje no tienen entrada y se resuelven con la capa 3, o salen como
\texttt{scheduled} («Solo horario») si el horario corregido ya ha pasado, aunque el autobús se viera un minuto antes.

Por eso, al leer el feed, se le añade la última entrada de cada autobús que falta en la lectura nueva si se cumplen
las tres condiciones siguientes:
\begin{itemize}
  \item el dato tiene como mucho $H=180$~s de antigüedad (\texttt{HOLD\_SECONDS}): $0\le\tnow-\tgps\le H$, con
        $t_{\mathrm{rec}}$ si no hay posición;
  \item el autobús no aparece en la lectura nueva con otro viaje;
  \item su viaje no lo lleva ya otro autobús en la lectura nueva.
\end{itemize}
La entrada se usa tal cual, y la capa~1 parte de su $\tgps$~\eqref{eq:eta-capa1}. Como ya se explicó, la antigüedad del
dato no sesga la llegada: solo la hace algo menos precisa. La memoria es un fichero temporal por red en el servidor;
en un servidor recién arrancado no hay nada guardado y se usa solo la lectura nueva.

\begin{table}[ht]
\centering
\caption{Llegadas de Autobuses que Faltaban en el Feed, con y sin Guardar su Último Dato}\label{tab:memoria-feed}
\begin{tabular}{@{}lrrrr@{}}
\toprule
Llegada en & $n$ & Sin guardar & Guardando 180~s & Mediana del error\\
\midrule
0--5 min & 9\,572 & 54{,}5\,\% & 77{,}5\,\% & de 1{,}8 a 1{,}0 min\\
5--10 min & 8\,945 & 51{,}6\,\% & 65{,}2\,\% & \\
10--20 min & 8\,676 & 37{,}3\,\% & 42{,}8\,\% & \\
\bottomrule
\end{tabular}
\nota{Aciertos dentro de $\pm2$~min contra el paso real, con las capturas del 2 al 4 de noviembre de 2026: cada
lectura en la que falta un autobús visto en los 180~s anteriores y que vuelve después al mismo viaje.}
\end{table}

La Tabla~\ref{tab:memoria-feed} recoge el resultado. Guardar el dato mejora los aciertos en los tres días y en los tres
horizontes; con 60, 120 y 300~s, la mejora a 0--5~min es del 55{,}7 al 80{,}9, del 54{,}8 al 79{,}6 y del 54{,}3 al
76{,}0\,\% (cada umbral mide sus propias lecturas, las de autobuses vistos en los $H$ segundos anteriores, y por eso
cambia también la cifra sin guardar). Cuanto más antiguo es el dato, algo menos acierta, y 180~s cubre la mayoría de
los cortes. El coste es el autobús que desaparece de verdad (acaba o se pierde el viaje): sus paradas siguen «en
directo» hasta que su último dato cumple $H$ segundos, cuando sin guardarlo habrían salido con la capa~3, que tampoco
se cumple.

\section{Viajes que Pasan de las 24:00}
\label{sec:mas-de-24h}

El horario cuenta las horas desde la medianoche del día del viaje: un autobús del viernes que pasa a las 02:10 del
sábado figura a las 26:10 ($94\,200$~s). En el horario de Bizkaibus del 5 de noviembre de 2026, 52 viajes pasan de las 24:00, y 43 de ellos salen ya pasada esa
hora: 6 circulan todas las noches, hasta las 26:05; 33, solo las noches del viernes o del sábado, y llegan hasta las
32:10, las 08:10 del día siguiente; y 4, solo en fechas concretas. El feed, en cambio, da $\tgps$ y $t_{\mathrm{rec}}$ en horas del día en curso, y la referencia
de un viaje que circula todos los días es la misma cada día. Con un solo emparejador, el autobús de las 23:30 de ayer,
visto a las 00:30, se emparejaría con el viaje de las 23:30 de hoy, y no con el suyo.

Por eso, cada entrada del feed se asigna al día de su viaje. Una entrada es de un viaje de ayer si la salida de cabecera
de su referencia queda más de 12~h por delante de la hora actual (\texttt{RealtimeMatcher::daysAgo}):
\begin{equation}
\label{eq:dia-viaje}
d_0^{\mathrm{feed}}-\tnow>43\,200\ \text{s}.
\end{equation}
A las 00:30 ($\tnow=1\,800$) son de ayer el viaje que salió a las 23:30 ($d_0^{\mathrm{feed}}=84\,600$) y el de las
24:10 ($87\,000$). La regla supone que el feed solo trae viajes en marcha o a punto de salir, y que la referencia da la
salida en horas del horario, como la de Bizkaibus~\eqref{eq:regex}.

Las entradas de ayer van a un emparejador propio en el que todo está en las horas del día del viaje: $\tgps$,
$t_{\mathrm{rec}}$ y $\tnow$ se cuentan con $86\,400$~s más (las 02:10 son las 26:10), la fecha de servicio $D$ es la de ayer y
el registro de pasos, el de ayer (\texttt{RealtimeMatcher::vehiclesOfDay}). El cálculo es el de las secciones
anteriores, sin cambios: la franja de las 26:10 de un viernes, por ejemplo, es la \texttt{laborable.madrugada} del
viernes (definición~\ref{def:franja}). El emparejador de hoy solo recibe las entradas de hoy, y una entrada sin ninguna
hora de hoy no va a ninguno (sección «\nameref{sec:feed}»).

La aplicación pasa al emparejador de ayer las salidas de ayer que pasan de las 24:00, con sus horas del horario, y
devuelve sus llegadas en horas de hoy, $\ETA-86\,400$ (sección «\nameref{sec:ultima-salida}»). De madrugada, la ficha de
un viaje muestra el viaje de ayer mientras este no haya terminado, es decir, hasta 10~min después de su última parada
(\texttt{LATE\_GRACE\_SECONDS}); en la red de autobús, solo si el feed no trae ningún autobús en el viaje de hoy. En el
mapa de una línea, cada autobús se sitúa con el emparejador de su día.

En la práctica, de 23:58 a 05:00 el feed de Bizkaibus no se actualiza (sección «\nameref{sec:feed}»), así que los
viajes nocturnos salen con la capa~3 o con el horario; el mapa no dice que siga a ningún autobús, y la ficha no dice
que a un viaje le falte autobús. Antes, el mapa de la línea 3136 mostraba toda la noche seis autobuses vistos a las
23:58. El emparejador de ayer sirve para los feeds que siguen actualizándose de noche y, en Bizkaibus, para los 11
viajes de los horarios del viernes y del sábado que acaban después de las 29:00 (las 05:00 del día siguiente), en las
líneas 3923 y 3932. De día, ninguna entrada es de ayer: con el código anterior y el nuevo sobre el mismo feed, 28 paradas, 256 fichas y 56 mapas de línea
dan las mismas respuestas.

\section{La Decisión Completa}
\label{sec:decision}

El algoritmo de la Figura~\ref{alg:fila} reúne las reglas anteriores en el orden exacto en que se aplican. Cualquier error
inesperado durante el cálculo de una fila (un dato corrupto, una base de datos inaccesible) se captura y la fila sale
con la capa 3 o con el horario publicado: una fila nunca impide responder a las demás.

\begin{figure}[ht]
\caption{Estado y Llegada de una Fila (Motor de Autobús)}\label{alg:fila}
\begin{algorithmic}[1]
\Require fila $(J,q,a_\dest)$, entrada del feed $e$ (o ninguna), $\tnow$
\Ensure (estado, $\ETA$, $\delta$)
\State $P\gets$ posición de $e$ en $J$ con su contraste (capítulo «\nameref{cap:localizacion}» y sección «\nameref{sec:desconfianza}»)
\State $G\gets$ registro de $e$ en $J$ (sección «\nameref{sec:capa2}»)
\If{$e$ no existe \textbf{o} ($P$ no existe \textbf{y} $G$ no existe)}
  \State \Return \Call{SinDatos}{$J,q,a_\dest$}
\EndIf
\If{hay $t_{\mathrm{paso}}$ (sección «\nameref{sec:paso}»)}
  \State \Return $(\texttt{departed},\ \min(t_{\mathrm{paso}},\tnow),\ 0)$
\EndIf
\If{$P$ no existe \textbf{y} ($q$ desconocido \textbf{o} la capa 2 no da llegada para $q$)}
  \State \Return \Call{SinDatos}{$J,q,a_\dest$}
\EndIf
\If{$P$ existe}
  \State $(\ETA,\delta)\gets$ capa 1, ecuaciones~\eqref{eq:eta-capa1}, \eqref{eq:paso-capa1} o, en la cabecera,~\eqref{eq:eta-cabecera}; $\pi\gets$ verdadero
\Else
  \State $\ETA\gets$ capa 2, ecuación~\eqref{eq:eta-capa2}; $\delta\gets\ETA-a_\dest$; $\pi\gets$ falso
\EndIf
\If{$|\delta|>2\,700$} $\delta\gets$ nulo \EndIf \Comment{retraso no creíble: se conserva la llegada}
\State $\alpha\gets(\ETA<\tnow-120)$ \Comment{ya debería haber pasado}
\If{$\alpha$ \textbf{y no} $\pi$} \State \Return \Call{SinDatos}{$J,q,a_\dest$} \EndIf
\If{\textbf{no} $\pi$} \State \Return $(\texttt{estimated},\ \ETA,\ 0)$ \EndIf
\If{$\ETA<\tnow$} \Comment{llegando, sin confirmar el paso}
  \If{$\delta$ no es nulo} $\delta\gets\max(\delta,\ \tnow-a_\dest)$ \EndIf
  \If{$\tnow-\ETA>180$} \State \Return $(\texttt{likely\_departed},\ \tnow,\ \delta)$ \EndIf
  \State \Return $(\texttt{live},\ \tnow,\ \delta)$
\EndIf
\State \Return $(\texttt{live},\ \ETA,\ \delta)$
\Statex
\Function{SinDatos}{$J,q,a_\dest$}
  \If{$e$ no existe \textbf{y} el registro de pasos tiene $(n,t)$ para $J$ con $q<n$} \State \Return $(\texttt{departed},\min(t,\tnow),0)$ \EndIf
  \State \Return capa 3 (sección «\nameref{sec:capa3}»): \texttt{estimated}, \texttt{likely\_departed} o \texttt{scheduled}
\EndFunction
\end{algorithmic}
\end{figure}

\subsection{El Tope de 45 Minutos} Un desfase de más de 45~minutos, de retraso o de adelanto ($|\delta|>2\,700$~s,
\texttt{MAX\_PLAUSIBLE\_DELAY\_SECONDS}), casi siempre indica un emparejamiento erróneo, no un retraso real: en los datos, los retrasos reales llegan a $+37$~min y
los viajes mal asignados dan valores como $+111$~min. Por encima del tope el retraso no se afirma ($\delta$ nulo), pero
la llegada se conserva, porque sigue saliendo de los datos del propio autobús y es útil. Se probó a ocultar también
la llegada, y la precisión empeoraba (del 85,2\,\% al 84,1\,\% de aciertos dentro de $\pm2$~min).

\subsection{El Margen de 120 Segundos} Una llegada de la capa~2 hasta dos minutos en el pasado
(\texttt{STALE\_GRACE\_SECONDS}) se sigue mostrando como estimada: el autobús puede estar en la parada o saliendo de
ella. Más allá, la estimación ya no es creíble y la fila pasa a la capa~3. Con GPS no hay margen: una llegada ya
pasada sale con $\ETA=\tnow$ hasta que se confirme el paso (sección «\nameref{sec:llegando}»).

\section{Retraso y Minutos que Se Muestran}
\label{sec:retraso-mostrado}

El desfase en la posición $\delta_{\mathrm{pos}}$ no es el retraso que entiende el usuario: compara al autobús con un
horario interpolado en un punto donde no hay parada. Lo que el usuario compara es la hora publicada de \emph{su} parada
con la llegada que se le muestra. Por eso, el retraso que se muestra es siempre una diferencia entre una llegada
estimada y una hora publicada:

\begin{itemize}
  \item \textbf{En la lista de una parada}, para cada fila,
        \begin{equation}
        \label{eq:retraso-lista}
        \text{retraso}=\redondeo\!\big(\delta/60\big)\ \text{min},\qquad \text{minutos}=\redondeo\!\big((\ETA-\tnow)/60\big),
        \end{equation}
        con el retraso nulo si $\delta$ es nulo. Las filas con menos de $-2$ minutos se retiran de la lista. Las demás
        se ordenan por $(\mathbf{1}[\text{estado}=\texttt{likely\_departed}],\ \ETA)$: por llegada estimada, no por la
        hora publicada (un autobús adelantado puede llegar antes que uno anterior en el horario), y con las filas
        \texttt{likely\_departed} al final.
  \item \textbf{En la ficha de un autobús}, el retraso es el de su próxima parada: con posición, $\ETA(p_n)-a_n$
        calculado con la capa 1. Se muestra $\redondeo(\cdot/60)$ minutos, o nada si $|\delta|$ pasa de 45~minutos. Sin
        posición se calcula $\delta_{\mathrm{reg}}$ (o 0 si tampoco hay registro), pero no se afirma: la ficha se marca
        como estimada y sin retraso, igual que la lista.
\end{itemize}

\noindent La aplicación traduce el retraso $r$ en minutos a un texto: «Retraso $+r$m» si $r>0$, «Adelantado $r$m» si
$r\le-2$, con su signo («Adelantado $-2$m»), y «En hora» en otro caso. La asimetría es deliberada: un autobús que llega un minuto antes no se anuncia como
adelantado, porque la resolución del dato es de un minuto. Nótese el efecto del redondeo de PHP:
$\redondeo(-90/60)=\redondeo(-1{,}5)=-2$, de modo que un adelanto de 90~s ya se anuncia.

Con esta definición, la lista y la ficha nunca se contradicen: se comprobó con el feed real que ninguna de 232 salidas
mostraba un retraso incoherente con su llegada.

\section{Después de la Última Salida del Día}
\label{sec:ultima-salida}

La lista se construye con el calendario del día y, de madrugada, con las salidas del día anterior que pasan de las
24:00, con el calendario de ayer y sus horas pasadas a las de hoy ($a-86\,400$: las 26:10 son las 02:10), intercaladas
con las de hoy por hora de salida. Las de la red de autobús pasan por el emparejador de ayer (sección
«\nameref{sec:mas-de-24h}») y las de las demás redes, por el emparejamiento por viaje de la sección
«\nameref{sec:otras-redes}». Cuando ya no queda ninguna salida por delante (las que acaban de salir siguen en la lista
por si alguna viene con retraso), se añaden detrás las primeras del día siguiente, con su calendario y sus horas
desplazadas 24~h ($a+86\,400$), en estado \texttt{scheduled} y sin pasar por el motor: el feed de hoy no dice nada de
esos viajes. La aplicación las muestra
con la hora («mañana 07:00») en lugar de los minutos.

\section{Ejemplo Completo}
\label{sec:ejemplo-completo}

La Tabla~\ref{tab:escenarios} recoge los dieciséis escenarios de los vectores de prueba (Apéndice~\ref{ap:vectores})
sobre el viaje sintético del ejemplo~\ref{ej:posicion}, con los tramos aprendidos, las proporciones y el factor del
ejemplo~\ref{ej:tiempo-esperado}. Las horas publicadas son $28\,800$, $28\,890$, $29\,010$, $29\,310$ y $29\,430$~s
(8:00:00 a 8:10:30), salvo en el escenario~8. Todas las cifras se han obtenido con el código de producción.

\begin{table}[ht]
\centering
\caption{Estado y Llegada Estimada en Cada Parada del Viaje Sintético}\label{tab:escenarios}
\setlength{\tabcolsep}{3.5pt}
\begin{tabular}{@{}lrccccc@{}}
\toprule
Esc. & $\tnow$ & $S0$ & $S1$ & $S2$ & $S3$ & $S4$\\
\midrule
1 & 28\,980 & S 28\,800 & S 28\,890 & V 29\,020 (+10) & V 29\,220 ($-90$) & V 29\,318 ($-112$)\\
2 & 28\,980 & S 28\,800 & S 28\,890 & V 29\,020 (+10) & V 29\,290 ($-20$) & V 29\,398 ($-32$)\\
3 & 28\,980 & S 28\,860 & S 28\,950 & E 29\,020 & E 29\,220 & E 29\,318\\
4 & 28\,980 & S 28\,860 & S 28\,950 & E 29\,020 & E 29\,220 & E 29\,318\\
5 & 29\,420 & P 28\,800 & P 28\,890 & P 29\,010 & P 29\,310 & P 29\,430\\
6 & 28\,700 & P 28\,800 & E 28\,970 & E 29\,110 & E 29\,310 & E 29\,408\\
7 & 29\,600 & P 28\,800 & P 28\,890 & P 29\,010 & P 29\,310 & P 29\,430\\
8 & 31\,200 & P 28\,800 & P 28\,890 & P 29\,010 & Pp 30\,810 & Pp 30\,930\\
9 & 28\,980 & P 28\,800 & P 28\,890 & E 29\,110 & E 29\,310 & E 29\,408\\
10 & 28\,700 & P 28\,800 & E 28\,990 & E 29\,150 & E 29\,350 & E 29\,468\\
11 & 28\,700 & P 28\,800 & E 28\,930 & E 29\,110 & E 29\,302 & E 29\,424\\
12 & 28\,210 & V 28\,800 (0) & V 28\,970 (+80) & V 29\,110 (+100) & V 29\,310 (0) & V 29\,408 ($-22$)\\
13 & 28\,210 & P 28\,800 & E 28\,970 & E 29\,110 & E 29\,310 & E 29\,408\\
14 & 29\,100 & S 28\,800 & S 28\,890 & V 29\,100 (+90) & V 29\,220 ($-90$) & V 29\,318 ($-112$)\\
15 & 29\,040 & S 28\,800 & S 28\,890 & V 29\,040 (+30) & V 29\,214 ($-96$) & V 29\,311 ($-119$)\\
16 & 29\,320 & S 28\,800 & S 28\,890 & Pp 29\,320 (+310) & V 29\,320 (+10) & V 29\,320 ($-110$)\\
\bottomrule
\end{tabular}
\nota{Esc.: escenario. Llegadas en segundos. V: \texttt{live}; E: \texttt{estimated}; S: \texttt{departed}; P: \texttt{scheduled}; Pp: \texttt{likely\_departed}. Entre paréntesis, el retraso en segundos cuando se afirma. Escenarios: 1, GPS a 700~m del origen y 11~m del trazado; 2, el mismo sin trazado; 3, GPS a 60~m del trazado; 4, sin GPS, con $n=3$ y $\tref=28\,950$; 5, GPS y registro discrepan; 6 y 7, sin datos del autobús, antes de salir y con el horario corregido ya pasado, respectivamente; 8, horario publicado muy lento ($S3$ y $S4$ a $30\,810$ y $30\,930$); 9, GPS de hace 680~s; 10, como el 6, con la corrección por viaje de los vectores de prueba; 11, el 10 en un festivo; 12, el autobús espera en la primera parada 10~min antes de salir; 13, el mismo a 600~m de ella; 14, como el 1 con $\tnow=29\,100$: la llegada a $S2$ ya pasó sin confirmar el paso; 15, GPS 50~m pasada $S2$ con el feed aún en ella (no basta para darla por pasada); 16, como el 14 con $\tnow=29\,320$: más de 3~min llegando sin confirmar el paso.}
\end{table}

\subsection{Escenario 1} Es el del ejemplo~\ref{ej:posicion}: $\Tpos=\tgps=28\,950$. Las paradas $S0$ y $S1$ tienen
$a_\dest<\Tpos$ y salen como pasadas a las $28\,950-(28\,950-28\,800)=28\,800$ y $28\,890$. Para $S2$, $S3$ y $S4$ la
llegada es $\tgps$ más los tiempos esperados del ejemplo~\ref{ej:tiempo-esperado}: $28\,950+70=29\,020$,
$\redondeo(28\,950+270{,}4)=29\,220$ y $\redondeo(28\,950+367{,}72)=29\,318$. El horario publicado de $S3$ y $S4$ es
90 y 112~s más lento de lo esperado, y la aplicación anunciaría «Adelantado $-2$m» en ambas
($\redondeo(-1{,}5)=-2$ y $\redondeo(-1{,}87)=-2$).

\subsection{Escenario 2} Sin trazado, la proyección sobre la recta $S1S2$ da el mismo $\Tpos$, pero los tramos sin
aprender ya no tienen longitud y caen al factor de ritmo: las llegadas a $S3$ y $S4$ pasan a $29\,290$ y $29\,398$.

\subsection{Escenarios 3 y 4} Con el GPS a 60~m (fuera del pasillo de 50~m) o sin GPS, se usa el registro: $n=3$,
$\tref=28\,950$, $\theta=\redondeo(0{,}5\cdot140)=70$. La llegada a $S2$ es $28\,950+70=29\,020$, y a $S3$,
$\redondeo(28\,950+70+200{,}4)=29\,220$: las mismas cifras que con GPS, porque en este ejemplo el autobús estaba
justo a mitad del tramo. Las paradas pasadas salen a $\tref-\max(0,a_2-a_\dest)$: $28\,950-90=28\,860$ para $S0$ y
$28\,950$ para $S1$. Las demás salen como estimadas y sin retraso.

\subsection{Escenario 5} El contraste de la sección «\nameref{sec:desconfianza}» invalida GPS y registro. La capa 3 da para
$S2$ un horario corregido de $28\,800+60+250=29\,110<\tnow$, que no cumple las condiciones~\eqref{eq:probablemente}
de \texttt{likely\_departed}, de modo que todas las paradas salen como \texttt{scheduled}, con su hora publicada.

\subsection{Escenario 6} Sin datos y antes de la salida, el horario corregido es $28\,860+E$: $28\,970$ para $S1$
(110~s aprendidos), $29\,110$ para $S2$ (110+140), $29\,310$ para $S3$ ($250+200{,}4=450{,}4$, redondeado a 450) y
$29\,408$ para $S4$ ($450{,}4+97{,}32=547{,}72$, redondeado a 548). La cabecera $S0$ no tiene $a_\dest>d_1$ y sale como
\texttt{scheduled}.

\subsection{Escenario 8} Con todos los tramos aprendidos (110, 140, 240 y 100~s) y un horario publicado que pone
$S3$ a las $30\,810$, el horario corregido es $C=28\,860+490=29\,350$. A las $31\,200$ se cumple $C<31\,200-600$,
$a_\dest-C=1\,460\ge600$ y $\lambda=1$: estado \texttt{likely\_departed}. Para $S2$, $C=29\,110$ es posterior a su hora
publicada y la regla no se aplica.

\subsection{Escenario 9} Un dato de hace 680~s supera la antigüedad máxima de 600~s: no hay posición ni registro, y
las filas van a la capa 3 como en el escenario 6, salvo $S1$, cuyo horario corregido ($28\,970$) ya es anterior a
$\tnow=28\,980$.

\subsection{Escenarios 10 y 11} Con la tabla de corrección de los vectores de prueba (Apéndice~\ref{ap:vectores}), el
viaje tiene en laborable 20, 40 y 60~s para sus tres tercios. Las paradas $S1$ a $S4$ ocupan las posiciones 1 a 4 de
cinco y caen en los tercios 0, 1, 1 y 2~\eqref{eq:tercio}: las llegadas del escenario 6 se desplazan 20, 40, 40 y
60~s. En el festivo, $D'$ es el domingo. La tabla de tramos no tiene valores en franjas de fin de semana, así que se
usan los de cualquier hora (\texttt{all}: 100 y 150~s); el factor vale 1, los tramos sin aprender usan las proporciones
\texttt{all} y la corrección de fin de semana es de $-30$, 0 y 30~s. Para $S4$, $C=28\,860+\redondeo(100+150+192+92{,}28)+30=29\,424$.

\subsection{Escenarios 12 y 13} Diez minutos antes de salir, el autobús está a 11~m de la primera parada con $n=2$: el feed
ya da la segunda como próxima parada. Queda situado en la cabecera (sección «\nameref{sec:cabecera}»), con
$\Tpos=a_1$ y la hora de la posición igual a la de salida, y todas las paradas salen en directo, $S0$ incluida, que no
cuenta como pasada aunque $n=2$.
Las llegadas son las de la capa 3 del escenario 6, porque el autobús aún no ha recorrido nada; el retraso que se
afirma es el de esas llegadas frente al horario publicado. A 600~m de la parada (escenario 13) puede estar acabando
otro viaje: el autobús no se sitúa en la cabecera, y las filas son las del escenario 6.

\chapter{Aprendizaje Continuo}
\label{cap:aprendizaje}

Los parámetros del modelo (los tiempos por tramo, las proporciones por longitud y los factores de ritmo) no se
escriben a mano: se aprenden del propio feed en tiempo real. Este capítulo describe el ciclo completo: cómo se captura
el feed, cómo se reconstruye el instante real en que cada autobús pasa por cada parada, qué muestras se extraen, cuánto
se conservan y con qué reglas se decide si una nueva calibración sustituye a la anterior.

\section{El Ciclo}
\label{sec:ciclo}

\begin{enumerate}
  \item \textbf{Captura}, continua y en bloques de una hora de reloj: se descarga el feed cada 20~s y se guarda cada
        versión distinta.
  \item \textbf{Muestras y medida}, al terminar cada bloque: se reconstruyen los pasos por parada, se extraen las
        muestras de tramo y de ritmo y se añaden a dos ficheros acumulados; además se mide la precisión de las
        llegadas (capítulo «\nameref{cap:evaluacion}») y se archivan las capturas en bruto del bloque.
  \item \textbf{Recalibración}, cada día a las 04:00 UTC: con todo lo acumulado se calculan nuevos parámetros, se
        validan y, solo si mejoran y cambian, se escriben.
  \item \textbf{Despliegue}, cada día a las 03:00 UTC: la aplicación se reconstruye con la última versión de los
        parámetros. Producción usa, por tanto, parámetros de entre 23 y 47~horas de antigüedad, salvo despliegue manual.
\end{enumerate}

Todo el ciclo usa el mismo código de localización y predicción que producción (la misma clase), de modo que lo que se
aprende y lo que se mide es lo que ve el usuario. La captura se ejecuta en una plataforma de integración continua que
limita cada trabajo a seis horas y que retrasa o se salta muchos trabajos programados: hasta el 31 de octubre de 2026
la captura era de 45~minutos cada hora y solo se ejecutaron 163 de unas 600 capturas previstas, algunas con horas de
retraso. Desde entonces cada trabajo captura unas cinco horas y media y, si termina bien, lanza el siguiente, que
espera su turno y empieza en cuanto acaba el anterior; un disparo programado cada hora queda como respaldo.

\section{Captura del Feed}
\label{sec:captura}

Cuando se fijó el intervalo de descarga, en octubre de 2026, el servidor del feed publicaba una versión nueva cada
55~s aproximadamente (el 5 de noviembre, cada 64~s de mediana). Cada 20~s (\texttt{intervalo\_s}) se
descarga, y la versión se guarda solo si su marca \texttt{RecordedAtTime} es distinta de la anterior. Con el intervalo
inicial de 45~s se perdía alguna versión cuando dos publicaciones caían entre dos descargas (46 versiones guardadas en
45~minutos, de unas 49 publicadas); con 20~s se guardan todas, y muestrear más deprisa no aporta nada, porque el feed
no cambia entre publicaciones.

La captura se organiza en bloques que terminan a la hora en punto, de modo que ninguno cruza la medianoche y cada uno
usa el horario de su día. Las capturas no usan el emparejador de ayer de la sección «\nameref{sec:mas-de-24h}»: la
parte de un viaje posterior a la medianoche no se sitúa con el GPS, y no aporta muestras ni entra en las medidas. En
Bizkaibus afecta sobre todo a los 11 viajes de fin de semana que acaban después de las 29:00, porque el resto de la
noche el feed no se actualiza. Mientras se captura un bloque se procesa el anterior, sin huecos entre ellos. Para medir la
precisión de las predicciones de un bloque se usan también las capturas del bloque siguiente, solo para ver cuándo
pasan los autobuses que estaban a punto de llegar al terminar el bloque. Las capturas en bruto de cada bloque,
comprimidas (unos 250~KB por hora), se archivan junto con las muestras nuevas que salieron de él, sin límite de
antigüedad: cualquier cálculo se puede repetir más adelante aunque los ficheros acumulados descarten lo antiguo.

Al procesar una captura, cada fichero se identifica por su \texttt{RecordedAtTime}
(los duplicados se ignoran) y se convierte a fecha y segundos de Madrid. Las instantáneas se ordenan por tiempo y se
agrupan en \emph{ventanas de observación}: una ventana termina cuando el hueco hasta la siguiente instantánea supera
300~s (\texttt{MAX\_GAP\_BETWEEN\_SNAPSHOTS\_SECONDS}). Cada instantánea guarda el instante final de su ventana,
$t_{\mathrm{fin}}$.

\section{Pistas}
\label{sec:pistas}

Para cada instantánea, en orden, se construye el emparejador con $\tnow$ igual a la hora de la instantánea y la fecha
de la instantánea como fecha de servicio, y para cada entrada del feed con viaje conocido se calcula su posición
(capa 1, con el contraste de la sección «\nameref{sec:desconfianza}»). Si existe, se añade el punto $(\tgps,\Tpos)$ a la
\emph{pista} de la clave (vehículo, viaje).

\begin{definicion}[Pista]
La pista de un vehículo $v$ en un viaje $J$ durante una captura es la sucesión, en el orden de las instantáneas, de
los puntos $(t_k,\theta_k)$, con $t_k=\tgps$ y $\theta_k=\Tpos$ de cada instantánea en que el vehículo tenía posición
en ese viaje.
\end{definicion}

Solo se aprende de posiciones GPS válidas: las capas 2 y 3, que son estimaciones, nunca alimentan el aprendizaje, y
tampoco las posiciones en la cabecera (sección «\nameref{sec:cabecera}»), que no describen ningún recorrido.

\section{El Instante Real de Paso por una Parada}
\label{sec:paso-real}

No existe en el feed ningún dato de «el autobús pasó por esta parada». Se reconstruye a partir de la pista, usando el
horario como escala de distancia: cada $\theta_k$ indica, en horas del horario, dónde estaba el autobús.

\begin{definicion}[Paso real]
\label{def:paso-real}
Dada una pista $\big((t_k,\theta_k)\big)_{k=1}^{K}$ y la hora publicada $a$ de una parada, el paso real $P(a)$ se
obtiene recorriendo los pares consecutivos $k=2,\dots,K$ y tomando el primero que cumpla
\[
\theta_k>\theta_{k-1},\qquad t_k-t_{k-1}\le300\ \text{s},\qquad \theta_{k-1}\le a\le\theta_k;
\]
entonces
\begin{equation}
\label{eq:paso-real}
P(a)=t_{k-1}+\frac{a-\theta_{k-1}}{\theta_k-\theta_{k-1}}\,\big(t_k-t_{k-1}\big).
\end{equation}
Si ningún par cumple las condiciones, el paso real no se conoce.
\end{definicion}

Es una interpolación lineal entre las dos observaciones que rodean a la parada. Los pares en que el autobús no avanza
($\theta_k\le\theta_{k-1}$, por ejemplo una posición repetida) o separados más de cinco minutos se descartan, porque no
permiten interpolar con fiabilidad.

\begin{ejemplo}
Una pista $(28\,900,\,28\,880)$, $(28\,960,\,28\,940)$, $(29\,020,\,29\,030)$ y una parada con $a=28\,950$: el par que
la rodea es el segundo, y $P=28\,960+\frac{28\,950-28\,940}{29\,030-28\,940}\cdot60=28\,966{,}67$~s. Para una parada con
$a=28\,800$, anterior a la primera observación, el paso no se conoce.
\end{ejemplo}

Con una posición por minuto, el error de la interpolación está acotado por el tiempo entre dos
observaciones: si el autobús se detiene o acelera entre ellas, el paso real puede estar en cualquier punto del
intervalo. En la práctica se trata como una resolución de $\pm30$~s, que es también la del método independiente que se
usa para evaluar (sección «\nameref{sec:verdad}»).

\section{Muestras de Tramo}
\label{sec:muestras-tramo}

Para cada pista y cada par de paradas consecutivas $(p_i,p_{i+1})$ de su viaje, si se conocen $P_i=P(a_i)$ y
$P_{i+1}=P(a_{i+1})$ y $P_{i+1}>P_i$, se produce una muestra de tramo con los campos
\begin{center}
\small
\begin{tabular}{@{}ll@{}}
\toprule
Campo & Valor\\
\midrule
\texttt{track} & «$D|v|J$»: fecha, vehículo y viaje\\
\texttt{line}, \texttt{date} & línea y fecha de servicio\\
\texttt{at} & $\lfloor P_i\rfloor$ (truncado: parte entera)\\
\texttt{band} & $\franja(D,\lfloor P_i\rfloor)$\\
\texttt{from}, \texttt{to} & identificadores de $p_i$ y $p_{i+1}$\\
\texttt{scheduled} & $a_{i+1}-a_i$\\
\texttt{actual} & $\redondeo_1(P_{i+1}-P_i)$, con un decimal\\
\bottomrule
\end{tabular}
\end{center}

\noindent La franja es la del momento en que el autobús entra en el tramo. La pista incluye la fecha para que el mismo
vehículo en el mismo viaje en dos días distintos cuente como dos pistas.

\section{Muestras de Ritmo}
\label{sec:muestras-ritmo}

Las muestras de ritmo sirven para estimar $k$ (sección «\nameref{sec:k}»). En cada instantánea, para cada entrada con
posición y cada parada de su viaje con $a_\dest>\Tpos$, se registra la predicción: la hora de la posición, $\Tpos$ y
$a_\dest$. Cuando se conocen las pistas completas, se calcula el paso real $P(a_\dest)$ y la predicción se convierte en
muestra si cumple:
\begin{enumerate}
  \item el paso real se conoce;
  \item la ventana de observación se extiende al menos 10~minutos después de la hora publicada:
        $t_{\mathrm{fin}}\ge a_\dest+600$ (\texttt{OBSERVED\_AFTER\_PREDICTION\_SECONDS});
  \item el tiempo programado que faltaba es de al menos dos minutos: $H=a_\dest-\Tpos\ge120$~s
        (\texttt{MIN\_REMAINING\_SECONDS\_FOR\_FACTOR});
  \item el tiempo real que faltaba es positivo: $Y=P(a_\dest)-\tgps>0$.
\end{enumerate}
Se guardan la pista, el vehículo, la línea, la fecha, $\texttt{at}=\tgps$, la franja de $\tgps$, $H$, $Y$ y el paso
real (estos dos con un decimal).

\subsection{Por Qué Exigir la Ventana} Solo se conoce el paso real de los autobuses que pasan antes de que termine la
captura. Si se aceptaran todas las predicciones con paso conocido, las de autobuses lentos (que pasan después del
final) quedarían fuera y $k$ saldría sesgado hacia autobuses rápidos: es un problema de \emph{censura}. Exigir que la
ventana cubra al menos diez minutos después de la hora publicada reduce ese sesgo a los autobuses que van más de diez
minutos tarde respecto a su horario en ese tramo. La misma regla, con la llegada estimada en lugar de la publicada,
se aplica al medir el error (capítulo «\nameref{cap:evaluacion}»).

\subsection{Aclarado} Un autobús visto durante 45~minutos con veinte paradas por delante produce cientos de
predicciones muy correlacionadas. Para que no pese más que otros, se conservan como mucho 30 por pista
(\texttt{MAX\_SAMPLES\_PER\_TRACK}): si una pista tiene $c>30$ muestras, en su orden original, se conservan las de
índices
\begin{equation}
\label{eq:aclarado}
\Big\lfloor\frac{i\,c}{30}\Big\rfloor,\qquad i=0,1,\dots,29 .
\end{equation}
Con $c=100$ son los índices 0, 3, 6, 10, 13, 16, 20, \dots, 96: un muestreo uniforme que conserva toda la duración de
la pista.

\section{Acumulación y Retención}
\label{sec:retencion}

Las muestras se añaden a dos ficheros JSON Lines acumulados (formatos en el Apéndice~\ref{ap:formatos}). Al añadir:

\begin{itemize}
  \item se descartan las muestras con fecha anterior a 120~días antes de hoy (\texttt{SAMPLES\_RETENTION\_DAYS});
  \item en las muestras de tramo, si un tramo tiene más de 300 de un mismo tipo de día, laborable o fin de semana
        (\texttt{MAX\_SAMPLES\_PER\_SEGMENT}), se conservan las 300 más recientes de ese tipo por (fecha, \texttt{at}),
        sin alterar el orden del resto del fichero;
  \item en las de ritmo, si el total pasa de $500\,000$ (\texttt{MAX\_PACE\_SAMPLES}), se ordenan por (fecha,
        \texttt{at}) y se conservan las $500\,000$ más recientes.
\end{itemize}

\noindent Estos topes mantienen los ficheros por debajo de unos 120~MB y la recalibración por debajo de 2~GB de memoria,
con un margen amplio respecto a los límites de la plataforma. El tope hace además de memoria: al quedarse solo lo más
reciente, un cambio en la calle (un desvío, unas obras) entra en lo aprendido en uno o dos días.

Hasta el 5 de noviembre de 2026 el tope era de 300 muestras por tramo, sin separar el tipo de día. Con las primeras
capturas, 300 muestras eran unos diez días de un tramo muy transitado, y diez días siempre incluyen un fin de semana: se
comprobó con 25~días que en los 109 tramos (de 2\,487) que llegaban al tope las franjas estaban igual de repartidas que
en el resto. Desde el 2 de noviembre las capturas recogen unas diez veces más pasos al día, y el 5 de noviembre 533 de
los 2\,757 tramos estaban en el tope: muchos guardaban solo los últimos 3 a 5~días, y 127 no tenían ninguna muestra de
fin de semana. En esos tramos el sábado y el domingo se usaba el tiempo de laborable, y el lunes, al revés, la lista
estaba llena de muestras del fin de semana.

Con un tope por tipo de día cada lista se renueva por su cuenta y siempre queda un fin de semana guardado. Para medirlo
se aprendió con las muestras anteriores a cada día y se probó en ese día (Tabla~\ref{tab:tope-tipo-dia}). Con el tope
de 300 no se puede comparar con los datos guardados, porque el fichero ya venía recortado por el tope anterior; por eso
se reprodujo la misma situación con topes menores en los tramos que conservan todo su historial. El tope por tipo de día
nunca empeora y mejora sobre todo el lunes que sigue al fin de semana.

\begin{table}[ht]
\centering
\caption{Error por Tramo con el Tope por Tramo y con el Tope por Tramo y Tipo de Día}\label{tab:tope-tipo-dia}
\begin{tabular}{@{}llrrr@{}}
\toprule
Tope & Día de prueba & Pasos & Error medio (s) & Dentro de $\pm30$~s (\%)\\
\midrule
20 & sábado 31 oct. & 22\,787 & 19,3 $\to$ 19,0 & 82,8 $\to$ 83,3\\
20 & domingo 1 nov. & 21\,441 & 19,1 $\to$ 18,9 & 83,3 $\to$ 83,6\\
20 & lunes 2 nov. & 52\,910 & 21,4 $\to$ 20,3 & 79,4 $\to$ 81,1\\
20 & martes 3 nov. & 58\,509 & 20,9 $\to$ 20,9 & 80,4 $\to$ 80,4\\
40 & sábado 31 oct. & 22\,787 & 19,0 $\to$ 19,0 & 83,3 $\to$ 83,5\\
40 & domingo 1 nov. & 21\,543 & 18,9 $\to$ 18,7 & 83,7 $\to$ 84,0\\
40 & lunes 2 nov. & 52\,910 & 20,7 $\to$ 19,9 & 80,5 $\to$ 81,7\\
40 & martes 3 nov. & 61\,081 & 20,3 $\to$ 20,3 & 81,0 $\to$ 81,1\\
80 & sábado 31 oct. & 22\,787 & 19,0 $\to$ 19,0 & 83,5 $\to$ 83,5\\
80 & domingo 1 nov. & 21\,543 & 18,8 $\to$ 18,7 & 83,8 $\to$ 84,0\\
80 & lunes 2 nov. & 52\,910 & 20,0 $\to$ 19,8 & 81,7 $\to$ 82,0\\
80 & martes 3 nov. & 61\,081 & 20,1 $\to$ 20,0 & 81,5 $\to$ 81,6\\
\bottomrule
\end{tabular}

\nota{Error de cada tramo entre dos paradas seguidas, no de la llegada completa: antes de la flecha, con un tope por
tramo; después, con un tope por tramo y tipo de día. Se aprendió con las muestras anteriores al día de prueba, con las
mismas medianas por franja de la sección «\nameref{sec:contraccion-tramo}». Solo cuentan los pasos con tiempo aprendido en los dos
casos.}
\end{table}

Para que una red recién puesta en marcha no empiece sin datos, el repositorio del motor incluye una semilla de muestras
de tramo que se usa si todavía no existe el fichero acumulado.

\section{Recalibración Diaria}
\label{sec:recalibracion}

Cada día se recalibran, en este orden, el factor de ritmo, los tiempos por tramo y la corrección por viaje. El orden
importa: la validación de los tiempos por tramo compara contra el factor de ritmo que se acaba de escribir, y la
corrección por viaje se mide contra el horario corregido que dan los tiempos recién escritos.

\subsection{Pliegues}

La validación es cruzada en dos pliegues por \emph{pistas}: cada muestra pertenece al pliegue
\begin{equation}
\label{eq:pliegue}
\mathrm{pliegue}(x)=\mathrm{crc32}\big(\texttt{track}(x)\big)\bmod2,
\end{equation}
donde $\mathrm{crc32}$ es el CRC de 32 bits de IEEE 802.3 (el de \texttt{zlib}, PNG o Ethernet) aplicado a los bytes
UTF-8 de la cadena, interpretado \emph{sin signo}. Todas las muestras de una misma pista caen en el mismo pliegue, de
modo que nunca se aprende y se mide con el mismo recorrido (un autobús en un viaje de un día). El mismo vehículo, en
otro viaje u otro día, puede caer en el otro pliegue, y los dos pliegues comparten los días. Se aprende con un pliegue
y se mide en el otro, y al revés.

\subsection{Factor de Ritmo}

\begin{enumerate}
  \item Se ajustan los factores con todas las muestras (ecuación~\eqref{eq:k-contraido} y los mínimos de la
        sección «\nameref{sec:k}»).
  \item Para cada pliegue $f$, se ajustan los factores con las muestras del otro pliegue. Después, para cada muestra
        de $f$ con factor aplicable $k\ne1$, se calculan los errores absolutos en minutos
        \[
        e_{\mathrm{antes}}=\frac{|\texttt{at}+H-P|}{60},\qquad e_{\text{después}}=\frac{|\texttt{at}+k\,H-P|}{60}.
        \]
  \item Se escribe el fichero de factores solo si hay alguna muestra de prueba, si la media de $e_{\text{después}}$
        es estrictamente menor que la de $e_{\mathrm{antes}}$, y si los factores nuevos no son iguales a los del fichero
        actual.
\end{enumerate}

\subsection{Tiempos por Tramo y Proporciones}

\begin{enumerate}
  \item Para cada pliegue $f$: se ajustan los tiempos por tramo con el otro pliegue y, para cada muestra de $f$ cuyo
        tramo tenga tiempo aprendido en su franja (ecuación~\eqref{eq:consulta-tramo}), se calculan
        \[
        e_{\mathrm{antes}}=\frac{|k\,h-y|}{60},\qquad e_{\text{después}}=\frac{|\tau-y|}{60},
        \]
        con $k$ el factor de ritmo vigente para la línea, la fecha y el instante de la muestra.
  \item Del mismo modo se validan las proporciones por longitud (con la longitud mediana de cada par de paradas en el
        horario): para cada muestra de prueba con $h>0$ y respaldo definido, $|k\,h-y|$ frente a
        $|\tau^{\mathrm{L}}-y|$. Las proporciones se incluyen en la tabla solo si la suma de errores baja.
  \item Se escribe el fichero solo si la suma de $e_{\text{después}}$ de los tiempos por tramo es estrictamente
        menor que la de $e_{\mathrm{antes}}$ y si los tiempos o las proporciones cambian respecto al fichero actual. Si
        los tiempos por tramo no mejoran, no se escribe nada, aunque mejoren las proporciones.
\end{enumerate}

\noindent La línea de base de la validación es siempre lo que haría el sistema sin lo aprendido (el horario por el
factor de ritmo). Con las 12\,755 medidas iniciales, el error medio por tramo bajó de 0,85~min a 0,39~min en las
medidas de prueba cuyo tramo estaba aprendido; con las proporciones por longitud, los tramos sin aprender pasaron de
0,88 a 0,69~min. La Tabla~\ref{tab:curva} cuenta todas las medidas de prueba, aprendidas o no, y por eso parte de
0,89~min.

\subsection{Corrección por Viaje}

\begin{enumerate}
  \item Con las muestras de tramo se reconstruye el paso real $R$ de cada viaje por cada parada y día: una muestra da
        el paso por su parada de salida (\texttt{at}) y, si falta el de la de llegada, también ese (\texttt{at} más el
        tiempo real del tramo). La pista dice el viaje, que se busca en el horario o en sus alias. Por eso el
        recalibrado reconstruye antes la base del horario del día.
  \item Para cada paso se calcula el desvío $e=R-C_0$, donde $C_0$ es el horario corregido~\eqref{eq:capa3} sin
        $\Delta_v$, con la $D'$ de ese día. Se saltan la cabecera, las paradas repetidas en un recorrido y los pasos a
        más de 45~minutos de su hora publicada. Cada desvío se recorta a $\pm900$~s para que un día de atasco no mande.
  \item Con $b$ la franja de la salida de cabecera según $D'$, se calculan
        \begin{equation}
        \label{eq:correccion-tablas}
        \begin{gathered}
        g(b)=Q_{0{,}4}(e_b),\qquad
        l(\ell,b)=\frac{n_{\ell,b}\,Q_{0{,}4}(e_{\ell,b})+20\,g(b)}{n_{\ell,b}+20},\\
        \Delta(v,\theta,p)=\frac{\sum_{\text{días}}Q_{0{,}4}(e_{\text{día}})+l(\ell,b)}{n_{\text{días}}+1},
        \end{gathered}
        \end{equation}
        donde $Q_{0{,}4}$ es el percentil 40 (con la misma regla que la mediana baja: el valor de la posición
        $\lfloor0{,}4(n-1)\rfloor$ de los ordenados), $e_b$ los desvíos de la franja, $e_{\ell,b}$ los de la línea en
        ella, $n_{\ell,b}$ su número y $e_{\text{día}}$ los del viaje en ese tercio cada día. Un tercio sin pasos toma el
        valor de su línea. Todo se guarda en segundos enteros. Hasta el 3 de noviembre de 2026 eran medias; el 40 es el
        percentil de partida, y la recalibración elige cada noche el que más acierta (sección
        «\nameref{sec:correccion-percentil}»).
  \item Se aprende con todos los días menos los 7 últimos y se mide en esos 7 (una sola partición; la evaluación de la
        sección «\nameref{sec:eval-correccion}» es más exigente y aprende cada día de prueba con todos los anteriores),
        con cada percentil entre el 30 y el 50;
        se queda el que más acierta a $\pm2$~minutos. La tabla, aprendida con todos los días y ese percentil,
        se escribe solo si el error medio absoluto de la capa 3 baja con la corrección y si cambia respecto a la
        actual. La herramienta informa además de cada festivo de las muestras, aprendiendo con los demás días.
\end{enumerate}

\noindent Aquí la validación es por días y no por pistas: la corrección es de cada viaje, y con pliegues por pistas
el mismo viaje de otro día estaría en las dos mitades.

\subsection{Qué Garantiza y Qué No}

Las reglas de escritura son conservadoras: un día con datos anómalos (una avería del feed, un festivo atípico) solo
cambia el modelo si el cambio mejora la predicción en recorridos que no se usaron para aprenderlo. Si no hay cambios,
no se escribe nada y el fichero conserva su fecha. La validación por pistas es una \emph{salvaguarda} diaria, no un
método para elegir reglas: para decidir si una regla nueva entra en el método se usa la evaluación por días del
capítulo «\nameref{cap:evaluacion}», que reproduce mejor lo que pasará con datos futuros.

\section{Curva de Aprendizaje}
\label{sec:curva}

Para saber cuánto mejora el modelo con más datos, se separó la mitad de los autobuses de la primera sesión de
captura (12\,755 medidas de tramo, un laborable a mediodía) para probar, y con la otra mitad se aprendieron los tiempos
por tramo usando cada vez más autobuses (Tabla~\ref{tab:curva}).

\begin{table}[ht]
\centering
\caption{Curva de Aprendizaje con la Primera Sesión de Captura}\label{tab:curva}
\begin{tabular}{@{}rrrrr@{}}
\toprule
Autobuses & Tramos vistos & Cobertura & Aprendido y $k$ & Solo $k\cdot h$\\
\midrule
16  & 278   & 1,2\,\%  & 0,89 & 0,89\\
45  & 610   & 13,0\,\% & 0,84 & 0,89\\
92  & 1\,001 & 32,2\,\% & 0,77 & 0,89\\
196 & 1\,458 & 62,8\,\% & 0,62 & 0,89\\
379 & 1\,916 & 77,5\,\% & 0,52 & 0,89\\
\bottomrule
\end{tabular}
\nota{Error medio por tramo, en minutos, en los autobuses de prueba. Tramos vistos: tramos distintos en las medidas
usadas para aprender (aprendidos o no). Cobertura: parte de las medidas de prueba cuyo tramo tiene tiempo aprendido en
su franja; las demás usan $k\cdot h$.}
\end{table}

\begin{figure}[ht]
\centering
\caption{Error por Tramo Según el Número de Autobuses}\label{fig:curva}
\begin{tikzpicture}
\begin{axis}[
  width=0.85\textwidth, height=5.5cm,
  xmode=log, log basis x=2,
  xlabel={Autobuses usados para aprender (escala logarítmica)},
  ylabel={Error medio por tramo (min)},
  xtick={16,45,92,196,379}, xticklabels={16,45,92,196,379},
  ymin=0.45, ymax=0.95,
  yticklabel style={/pgf/number format/fixed,/pgf/number format/precision=2,/pgf/number format/use comma},
  grid=major, grid style={gris!20},
  legend style={at={(0.02,0.05)},anchor=south west,font=\footnotesize,draw=none,fill=none},
  every axis plot/.append style={thick}]
\addplot[ruta,mark=*] coordinates {(16,0.89) (45,0.84) (92,0.77) (196,0.62) (379,0.52)};
\addlegendentry{Lo aprendido y, si no, $k\cdot h$}
\addplot[gris,dashed] coordinates {(16,0.89) (379,0.89)};
\addlegendentry{Solo $k\cdot h$}
\end{axis}
\end{tikzpicture}
\nota{Las medidas de prueba sin tramo aprendido (100\,\% menos la cobertura de la Tabla~\ref{tab:curva}) bajan del
98,8 al 22,5\,\%, un 40--45\,\% en cada una de las dos últimas duplicaciones, y el error sigue bajando.}
\end{figure}

El error no se había estabilizado con 379 autobuses. Desde entonces la captura pasó a ser horaria, y a finales de
octubre de 2026 las muestras acumuladas suman 195\,786 tramos de 18\,604 pistas en 26 días, con 2\,454 tramos aprendidos.

\chapter{Resultados}
\label{cap:evaluacion}

Un método de predicción solo vale lo que se puede medir. Este capítulo define la medida del error y explica cómo se
obtiene la verdad de referencia sin un dato de «el autobús pasó por aquí». Después presenta los resultados por
horizonte, por capa y por cobertura de tramos aprendidos, y los repite con 17 días de capturas. Con un protocolo por
días evalúa la contracción de la sección «\nameref{sec:contraccion-tramo}» y la corrección por viaje de la capa 3, y
compara el método con filtros de Kalman y redes neuronales. Mide también la corrección de los trazados, la espera en la
cabecera y la fiabilidad de los estados que muestra la aplicación, y termina con las ideas que se probaron y no
entraron. La interpretación de los resultados y sus limitaciones están en el capítulo «\nameref{cap:discusion}».

\section{Medida del Error}
\label{sec:medida}

Para cada predicción con paso real conocido, el error es
\begin{equation}
\label{eq:error}
e=\frac{t_{\mathrm{mostrada}}-t_{\mathrm{real}}}{60}\quad\text{minutos},
\end{equation}
donde $t_{\mathrm{mostrada}}$ es la llegada que vería el usuario (la $\ETA$ del capítulo «\nameref{cap:prediccion}») o, para
comparar, la hora publicada $a_\dest$. Un error positivo significa que el autobús llega \emph{antes} de lo que se
muestra.

Las predicciones se agrupan por \emph{horizonte}, el tiempo que de verdad faltaba para la llegada, porque la
precisión depende sobre todo de él. Para cada grupo se dan:

\begin{itemize}
  \item el \textbf{sesgo}, mediana de $e$: si es distinto de cero, el método se adelanta o se retrasa sistemáticamente;
  \item la \textbf{mediana} y el \textbf{percentil 90} de $|e|$;
  \item el \textbf{porcentaje de aciertos} dentro de $\pm1$, $\pm2$ y $\pm3$ minutos.
\end{itemize}

\noindent Los porcentajes dentro de un umbral son la medida principal porque responden a la pregunta del usuario
(«¿llegaré a tiempo si salgo cuando dice?») y, a diferencia del error cuadrático medio, no se dejan dominar por unos
pocos errores grandes. Como todas las predicciones se expresan en la misma unidad, las medidas en segundos y minutos son
comparables entre métodos (sección «\nameref{sec:marco-validacion}»). Los percentiles se calculan como en todo el método, sin interpolar
($P_p=x_{(\lfloor p(n-1)\rfloor)}$).

\section{La Verdad de Referencia}
\label{sec:verdad}

Se usan dos reconstrucciones del instante real de paso, independientes entre sí.

\subsection{Por Posición} Es la del aprendizaje (definición~\ref{def:paso-real}): interpolar la pista del autobús en
la hora publicada de la parada. Sirve para calibrar y para la validación cruzada, pero no es del todo independiente de
la predicción, porque ambas usan la misma proyección del GPS.

\subsection{Por Cambio de Parada} Para medir lo que ve el usuario se usa una reconstrucción que no depende del GPS: el
feed declara en cada instantánea la próxima parada $n$ de cada autobús. Si en la instantánea $m$ un autobús que sigue
en el mismo viaje desde la instantánea de la predicción declara por primera vez una próxima parada posterior a $q$,
el paso por $p_q$ ocurrió entre las instantáneas $m-1$ y $m$, y se toma
\begin{equation}
\label{eq:verdad}
t_{\mathrm{real}}=\frac{t_{m-1}+t_m}{2},
\end{equation}
con $t_m$ la hora más reciente de las posiciones de la instantánea $m$. La resolución es de medio intervalo entre
instantáneas, $\pm30$~s, y el error medido incluye ese margen. Se evalúan las paradas con $n\le q\le n+8$ y horizontes
de 0 a 20~minutos.

\section{Resultados de una Sesión de Captura}
\label{sec:resultados}

Las cifras de esta sección se midieron el 14 de octubre de 2026 con 15 instantáneas del feed real (14~minutos de un
laborable), que dan 7\,627 predicciones con paso real conocido. Es una muestra pequeña, pero en ella cada llegada se
puede calcular con las tres capas a la vez. La sección «\nameref{sec:varios-dias}» repite la medida principal con 17
días de capturas.

\subsection{Por Horizonte}

\begin{table}[ht]
\centering
\caption{Precisión de la Llegada Mostrada con GPS y del Horario Publicado, por Horizonte}\label{tab:resultados}
\begin{tabular}{@{}lccccc@{}}
\toprule
La parada queda a & Mediana de $|e|$ & $\pm1$ min & $\pm2$ min & $\pm3$ min & Horario, $\pm2$ min\\
\midrule
0--3 min   & 0,5 min & 80\,\% & 94\,\% & 96\,\% & 19\,\%\\
3--6 min   & 0,7 min & 63\,\% & 89\,\% & 95\,\% & 18\,\%\\
6--10 min  & 1,0 min & 49\,\% & 80\,\% & 90\,\% & 16\,\%\\
10--20 min & 1,2 min & 41\,\% & 73\,\% & 89\,\% & 13\,\%\\
\bottomrule
\end{tabular}
\nota{Capa 1 frente al paso real por la parada; el horario publicado se mide en las mismas llegadas.}
\end{table}

Con GPS, el método acierta dentro de $\pm2$~minutos unas cinco veces más que el horario publicado, y sin sesgo: la
mediana de $e$ es cero en todos los horizontes. La precisión decrece de forma regular con el horizonte, como predice la
tercera consecuencia de la proposición~\ref{prop:diferencias}.

\subsection{Por Capa}

Para medir las capas 2 y 3 se repite la evaluación sobre las mismas instantáneas, contra los mismos pasos reales,
quitando a todos los autobuses el GPS (capa 2) o cualquier dato del vehículo (capa 3).

\begin{table}[ht]
\centering
\caption{Aciertos Dentro de $\pm2$~Minutos por Capa, en las Mismas Llegadas}\label{tab:capas}
\begin{tabular}{@{}lcccc@{}}
\toprule
La parada queda a & \shortstack{Capa 1\\GPS} & \shortstack{Capa 2\\parada registrada} & \shortstack{Capa 3\\sin datos} & \shortstack{Horario\\publicado}\\
\midrule
0--3 min   & 94\,\% & 94\,\% & 42\,\% & 19\,\%\\
3--6 min   & 89\,\% & 84\,\% & 45\,\% & 18\,\%\\
6--10 min  & 80\,\% & 74\,\% & 43\,\% & 16\,\%\\
10--20 min & 73\,\% & 67\,\% & 39\,\% & 13\,\%\\
\bottomrule
\end{tabular}
\end{table}

En las mismas llegadas, la parada registrada, sin ningún GPS, queda a pocos puntos de la capa 1: el dato que más información aporta es saber
qué tramo recorre el autobús y desde cuándo. La capa 3, que no sabe nada del vehículo, duplica o triplica el acierto
del horario publicado, porque corrige el error de forma de la interpolación (esta muestra es anterior a la corrección
por viaje de la sección «\nameref{sec:eval-correccion}»). Según la longitud del tramo en que está el
autobús, la capa 2 acierta dentro de $\pm2$~min el 89\,\% en tramos de hasta 3~min, el 83\,\% de 3 a 6, el 86\,\% de 6
a 10 y el 52\,\% de 10 a 20 (frente al 18, 20, 15 y 0\,\% del horario en las mismas llegadas). En tramos de más de
20~min baja al 7\,\%, apenas por encima del horario (1\,\%); de ahí el umbral de la sección «\nameref{sec:capa2}».

\subsection{Por Cobertura de Tramos Aprendidos}
\label{sec:cobertura-tramos}

La variable que más explica el error, después del horizonte, es si todos los tramos del camino entre el autobús y la
parada tienen tiempo aprendido.

\begin{table}[ht]
\centering
\caption{Aciertos de la Capa 1 Según la Cobertura del Camino}\label{tab:cobertura}
\begin{tabular}{@{}lcc@{}}
\toprule
La parada queda a & Todo el camino aprendido & Algún tramo sin aprender\\
\midrule
0--5 min   & 96\,\% $\pm2$ · 99\,\% $\pm3$ (3\,923) & 80\,\% $\pm2$ · 88\,\% $\pm3$ (528)\\
5--10 min  & 87\,\% $\pm2$ · 96\,\% $\pm3$ (1\,811) & 59\,\% $\pm2$ · 75\,\% $\pm3$ (369)\\
10--20 min & 76\,\% $\pm2$ · 93\,\% $\pm3$ (439)   & 61\,\% $\pm2$ · 79\,\% $\pm3$ (77)\\
\bottomrule
\end{tabular}
\nota{Entre paréntesis, el número de llegadas. Solo cuentan las llegadas con datos del autobús, que son las que tienen cobertura
calculada: 7\,147 de las 7\,627.}
\end{table}

En llegadas a menos de 10~minutos, con todo el camino aprendido se acierta el 93\,\% dentro de $\pm2$~min; con un tramo
sin aprender, el 76\,\%; con dos o más, el 59\,\%. Con el 86\,\% de las llegadas con todo el camino aprendido, la
combinación da el 90\,\%. El 93\,\% es un techo para este feed: el 7\,\% restante se debe a la resolución de $\pm30$~s
del paso real, a una sola posición por minuto y a lo que nadie puede prever (semáforos, tráfico, esperas en parada).
Los métodos que además usan lo que ocurre ese mismo día lo mejoran poco (sección
«\nameref{sec:comparacion}»); superarlo con claridad exige otra fuente de datos, no más capturas de esta.

\section{Evaluación a lo Largo de Varios Días}
\label{sec:varios-dias}

Desde el 15 de octubre de 2026, cada captura horaria medía la precisión de la llegada mostrada con el método de la
sección «\nameref{sec:verdad}» y dejaba el resumen en el registro de la plataforma de ejecución. Se recuperaron esos
resúmenes: 88 capturas de 17 días (del 15 al 31 de octubre). De ellas, 61 tienen predicciones, porque en las de
madrugada el feed no trae posiciones del día (de 23:58 a 05:00 repite la lectura de las 23:58; sección
«\nameref{sec:feed}»). Suman 2\,369\,529 predicciones contrastadas con el paso real, de las que 2\,208\,468 quedan a
menos de 20~minutos (las de la Tabla~\ref{tab:varios-dias}). Hay dos diferencias con la sección anterior. La primera
es que estas cifras incluyen todas las llegadas que muestra la aplicación, de cualquier capa, y no solo las de la
capa 1. La segunda, que el método cambió durante esos días (por ejemplo, la contracción de la
sección «\nameref{sec:contraccion-tramo}» se adoptó al final del periodo), así que las cifras miden la aplicación tal
como estaba cada día.

Los intervalos de confianza del 95\,\% se obtienen por \emph{bootstrap} de capturas \citep{hastie2009}: se remuestrean
capturas enteras y no predicciones sueltas, porque las predicciones de una misma captura comparten autobuses, tráfico y
meteorología, y tratarlas como independientes daría intervalos demasiado estrechos.

\begin{table}[ht]
\centering
\caption{Precisión de Todas las Llegadas Mostradas en 17 Días, por Horizonte}\label{tab:varios-dias}
\begin{tabular}{@{}lrcccc@{}}
\toprule
La parada queda a & Predicciones & $\pm1$ min & $\pm2$ min (IC 95\,\%) & $\pm3$ min & Horario, $\pm2$ min\\
\midrule
0--3 min   & 534\,046 & 84\,\% & 96\,\% (96--96) & 98\,\% & 24\,\%\\
3--6 min   & 485\,601 & 67\,\% & 90\,\% (90--91) & 96\,\% & 24\,\%\\
6--10 min  & 548\,058 & 53\,\% & 82\,\% (81--82) & 92\,\% & 23\,\%\\
10--20 min & 640\,763 & 39\,\% & 67\,\% (65--68) & 82\,\% & 25\,\%\\
\bottomrule
\end{tabular}
\nota{Llegadas de cualquier capa en 61 capturas de unos 45~minutos, del 15 al 31 de octubre de 2026. IC: intervalo de
confianza por \emph{bootstrap} de capturas. El horario publicado se mide en las mismas llegadas. Los porcentajes de
cada captura se redondearon a enteros al imprimirse, y se promedian ponderados por su número de predicciones.}
\end{table}

La Tabla~\ref{tab:varios-dias} confirma la Tabla~\ref{tab:resultados} en el horizonte más corto (96\,\% frente a 94\,\%
dentro de $\pm2$~min a menos de 3~minutos) y la corrige en el más largo: entre 10 y 20~minutos se acierta el 67\,\% y no el 73\,\%,
porque la muestra de 14~minutos tenía pocas llegadas lejanas y la Tabla~\ref{tab:resultados} solo cuenta la capa 1,
mientras que aquí entran llegadas de cualquier capa. El horario publicado acierta en torno al 24\,\% en todos
los horizontes. La precisión apenas cambia con el tipo de día: dentro de $\pm2$~min, los laborables (40 capturas) dan
el 96, 90, 82 y 67\,\% por horizonte, y los fines de semana (21 capturas), el 96, 90, 80 y 64\,\%, con intervalos que se
solapan salvo a 6--10~minutos (82--83 frente a 79--81\,\%).

\begin{table}[ht]
\centering
\caption{Aciertos Dentro de $\pm2$~Minutos en 17 Días, por Origen del Dato y por Cobertura}\label{tab:varios-dias-origen}
\begin{tabular}{@{}lrc@{}}
\toprule
Llegadas a menos de 10 min & Predicciones & $\pm2$ min (IC 95\,\%)\\
\midrule
\multicolumn{3}{@{}l}{\emph{Origen del dato y tramo en que está el autobús}}\\
GPS, tramo corto                & 1\,299\,494 & 92\,\% (92--93)\\
GPS, tramo largo                & 58\,410     & 82\,\% (81--83)\\
Parada registrada, tramo corto  & 91\,533     & 77\,\% (76--78)\\
Parada registrada, tramo largo  & 1\,813      & 59\,\% (56--62)\\
Sin datos del autobús, tramo corto & 111\,557   & 71\,\% (69--73)\\
Sin datos del autobús, tramo largo & 4\,898     & 17\,\% (14--20)\\
\midrule
\multicolumn{3}{@{}l}{\emph{Tramos sin tiempo aprendido hasta la parada (llegadas con datos del autobús)}}\\
Ninguno     & 1\,251\,203 & 93\,\% (93--94)\\
Uno         & 120\,897    & 76\,\% (76--77)\\
Dos o más   & 49\,176     & 63\,\% (62--65)\\
\bottomrule
\end{tabular}
\nota{«Parada registrada»: llegadas con datos del autobús sin posición GPS utilizable (capa 2). «Sin datos del autobús»: llegadas
que no se muestran en directo (estimadas u horario corregido). Tramo corto: el autobús recorre un tramo de hasta
10~minutos programados; tramo largo: de más de 10~minutos.}
\end{table}

La Tabla~\ref{tab:varios-dias-origen} repite con 17 días los desgloses de las Tablas~\ref{tab:capas}
y~\ref{tab:cobertura}, y el orden se mantiene: en tramos cortos se acierta el 92\,\% dentro de $\pm2$~min con la
posición GPS, el 77\,\% con la parada registrada y el 71\,\% sin datos del autobús, y la cobertura del camino es, de
nuevo, la variable que más mueve la precisión
(del 93\,\% con todo el camino aprendido al 63\,\% con dos o más tramos sin aprender).

Desde el 31 de octubre la captura es continua y guarda, por bloque de una hora, los errores en histogramas de 15~s por capa, horizonte, franja y
tipo de día, con la llegada calculada también sin GPS y sin feed: la recalibración diaria publica con ellos las mismas
tablas con intervalos de confianza, ya separadas por capa.

\section{Evaluación por Días de la Contracción}
\label{sec:eval-contraccion}

La validación cruzada por pistas del capítulo «\nameref{cap:aprendizaje}» decide si se escribe una calibración, pero mezcla
en entrenamiento y prueba autobuses del mismo día, que comparten tráfico y meteorología. Para decidir si una
\emph{regla} nueva entra en el método se usa un protocolo más exigente, que reproduce cómo aprende el método en producción
\citep{hyndman2021}:

\begin{enumerate}
  \item Para cada uno de los últimos siete días $D$ de las muestras, se aprende con todas las muestras de días
        anteriores a $D$ y se prueba con las de $D$.
  \item Para cada pista de prueba, con sus tramos ordenados por \texttt{at}, y para cada tramo inicial, se encadenan los tramos
        siguientes mientras sean consecutivos (el destino de uno es el origen del siguiente) y tengan tiempo
        aprendido. Se van sumando el tiempo real y el predicho con la franja del tramo inicial.
  \item Cada vez que el tiempo real acumulado cae en uno de los horizontes $[4,8)$, $[8,12)$ o $[12,20)$~min, se anota
        el error absoluto de la suma.
  \item La regla se adopta si mejora en todos los horizontes y en casi todos los días, no solo en la media.
\end{enumerate}

\noindent Es una medida de punta a punta: un cambio puede mejorar cada tramo por separado y no la suma de varios, que
es de lo que se compone la llegada que ve el usuario. La Tabla~\ref{tab:contraccion} compara la mediana de la franja sola ($n_0=0$) con la
contraída, para dos valores de $n_0$, con las muestras acumuladas a 31 de octubre de 2026 (195\,786 muestras de tramo de 26 días;
días de prueba del 25 al 31 de octubre).\footnote{La herramienta de evaluación usa la mediana del tramo cuando la
franja tiene menos de dos pistas, aunque producción solo la guarda si hay muestras de dos franjas; esto afecta por igual
a las reglas comparadas.}

\begin{table}[ht]
\centering
\caption{Evaluación por Días de la Contracción Hacia la Mediana del Tramo}\label{tab:contraccion}
\begin{tabular}{@{}lrccc@{}}
\toprule
Horizonte & Recorridos & $n_0=0$ & $n_0=5$ & $n_0=10$\\
\midrule
\multicolumn{5}{@{}l}{\emph{Error medio (s)}}\\
4--8 min   & 107\,052 & 42,8 & \textbf{42,0} & 42,3\\
8--12 min  & 83\,336  & 62,7 & \textbf{61,7} & 62,4\\
12--20 min & 113\,581 & 89,3 & \textbf{88,2} & 89,4\\
\midrule
\multicolumn{5}{@{}l}{\emph{Dentro de $\pm60$~s}}\\
4--8 min   & & 75,6\,\% & \textbf{76,3\,\%} & 75,9\,\%\\
8--12 min  & & 58,2\,\% & \textbf{58,9\,\%} & 58,4\,\%\\
12--20 min & & 42,7\,\% & \textbf{43,0\,\%} & 42,8\,\%\\
\midrule
\multicolumn{5}{@{}l}{\emph{Dentro de $\pm120$~s}}\\
4--8 min   & & 96,0\,\% & \textbf{96,2\,\%} & 96,0\,\%\\
8--12 min  & & 87,4\,\% & \textbf{87,9\,\%} & 87,5\,\%\\
12--20 min & & 73,6\,\% & \textbf{74,1\,\%} & 73,4\,\%\\
\bottomrule
\end{tabular}
\nota{Error absoluto medio de la suma de tramos y aciertos dentro de $\pm60$ y $\pm120$~s.}
\end{table}

Con $n_0=5$ mejoran todas las medidas en todos los horizontes. Con $n_0=10$ la mejora tramo a tramo se mantiene en los
horizontes cortos, pero a 12--20~min empeora frente a no contraer: la contracción excesiva borra diferencias reales entre
franjas, que se acumulan al sumar muchos tramos (el término $\beta^2$ de la
proposición~\ref{prop:contraccion}). Cuando se adoptó la regla, con 193\,864 muestras de tramo de 25 días, la mejora se observó
en los siete días de prueba, y la validación por pistas bajó la suma de errores de 65\,854 a 63\,204~minutos.

\section{Corrección por Viaje de la Capa 3}
\label{sec:eval-correccion}

La corrección por viaje (sección «\nameref{sec:correccion-viaje}») se evaluó con el protocolo por días: de 238\,740
pasos reales con viaje reconocido en 27 días (del 6 de octubre al 1 de noviembre de 2026), se probaron los 7 últimos
días (87\,333 pasos de prueba), y para cada día de prueba se aprendió con todos los anteriores.

\subsection{Corrección y Partes del Recorrido}

La Tabla~\ref{tab:correccion} da el error de la hora que se mostraría sin ningún dato del autobús, sin corrección y
con cada nivel de la corrección.

\begin{table}[ht]
\centering
\caption{Error de la Hora sin Datos del Autobús, con y sin Corrección}\label{tab:correccion}
\begin{tabular}{@{}lrrrrr@{}}
\toprule
Hora mostrada & Error medio & $\pm60$~s & $\pm120$~s & $\pm300$~s & Desvío\\
\midrule
Horario publicado & 438~s & 13,1\,\% & 24,4\,\% & 49,8\,\% & $-263$~s\\
Capa 3 sin corrección & 210~s & 27,0\,\% & 46,6\,\% & 76,3\,\% & $+131$~s\\
Corrección de la franja, $g$ & 199~s & 20,9\,\% & 41,4\,\% & 80,2\,\% & $+14$~s\\
Corrección de la línea, $l$ & 177~s & 25,8\,\% & 47,7\,\% & 83,4\,\% & $+18$~s\\
Un valor por viaje & 164~s & 28,2\,\% & 51,1\,\% & 85,6\,\% & $+17$~s\\
\textbf{Por viaje y tercio}, $\Delta$ & \textbf{158~s} & \textbf{30,3\,\%} & \textbf{54,0\,\%} & \textbf{86,2\,\%} & $+17$~s\\
Por viaje y quinto & 159~s & 30,3\,\% & 53,9\,\% & 85,7\,\% & $+18$~s\\
Por viaje y parada & 162~s & 30,1\,\% & 53,3\,\% & 85,3\,\% & $+18$~s\\
\bottomrule
\end{tabular}
\nota{Desvío: media de la hora real menos la mostrada (el signo contrario al de $e$); positivo si el autobús llega después de lo mostrado. Pasos de
prueba: 87\,333, del 26 de octubre al 1 de noviembre de 2026. Cada fila usa las anteriores como respaldo: la del viaje
usa la de la línea cuando el viaje no tiene datos.}
\end{table}

Quitar el sesgo con una sola cifra por franja baja el error medio, pero empeora los aciertos ajustados ($\pm60$ y
$\pm120$~s): el sesgo no es igual en todas las líneas ni en todos los viajes. Con la corrección por viaje y tercio
mejoran todas las columnas. El error medio baja un 25\,\%, y el acierto a $\pm2$~minutos pasa del 47 al
54\,\%. Partir el recorrido en más trozos no mejora: con menos pasos por trozo, el ruido de cada día pesa más. Los
tiempos por tramo de la capa 3 se aprendieron con los 27 días, también los de prueba, y eso favorece a la fila sin
corrección: si acaso, la mejora medida se queda corta.

\subsection{La Media o un Percentil}
\label{sec:correccion-percentil}

Hasta el 3 de noviembre de 2026, cada valor de la corrección era una media de desvíos. Los desvíos tienen una cola
larga de retrasos (un día de atasco, un autobús que sale tarde de la cabecera), que sube la media por encima de lo
típico: con la media, la hora mostrada quedaba casi siempre algo después del paso del autobús, y quien se fiaba de ella
podía perderlo. Se probó a tomar en cada grupo (la franja; la línea y la franja; el viaje y el tercio de cada día) la
mediana o un percentil más bajo, con el resto del cálculo igual.

La prueba es la más exigente posible. Se reconstruyó lo que habría mostrado la aplicación el lunes 2 y el martes 3 de
noviembre de 2026, instantánea a instantánea (una de cada tres, de las 05:00 a las 23:00). Los tiempos por tramo y la corrección
se aprendieron solo con los días anteriores a cada uno, y cada fila «Estimado» sin autobús se comparó con el paso real
del autobús por esa parada (Tabla~\ref{tab:percentil}).

\begin{table}[ht]
\centering
\caption{Media o Percentil en la Corrección por Viaje}\label{tab:percentil}
\begin{tabular}{@{}lrrrrr@{}}
\toprule
Valor & $\pm2$~min & $\pm5$~min & Error & Antes & Después\\
\midrule
Media & 51,5 / 58,2\,\% & 83,5 / 87,8\,\% & 190 / 162~s & 26,6 / 22,4\,\% & 22,0 / 19,4\,\%\\
Mediana & 53,6 / 59,8\,\% & 83,5 / 87,9\,\% & 187 / 160~s & 22,4 / 19,1\,\% & 24,0 / 21,2\,\%\\
\textbf{Percentil 40} & \textbf{55,6 / 60,7\,\%} & 83,1 / 87,7\,\% & \textbf{187 / 159~s} & 17,1 / 15,2\,\% & 27,3 / 24,2\,\%\\
Percentil 35 & 56,0 / 60,5\,\% & 82,7 / 87,5\,\% & 188 / 160~s & 15,1 / 13,7\,\% & 28,9 / 25,8\,\%\\
\bottomrule
\end{tabular}
\nota{Lunes 2 / martes 3 de noviembre de 2026, sin verlos al aprender: 642\,427 y 488\,640 filas «Estimado» sin
autobús con paso real. Media: la que se usaba hasta el 3 de noviembre. Error: error absoluto medio. Antes: el autobús pasa más de 2~min antes de la hora mostrada (se
pierde si uno se fía). Después: pasa más de 2~min después (hay que esperarlo). El paso real es el punto medio de las
dos instantáneas entre las que el feed deja atrás esa parada.}
\end{table}

El percentil 40 acierta más que la media con $\pm2$~minutos los dos días (4,1 y 2,5 puntos), con casi el mismo acierto
a $\pm5$~minutos (0,4 y 0,1 puntos menos) y algo menos de error medio, y baja en un tercio las veces que el autobús pasa
más de 2~minutos antes de la hora mostrada. A cambio, hay que esperarlo más a menudo, que es el error menos costoso:
llegar un poco antes a la parada no hace perder el autobús. El percentil 35 acierta casi lo mismo que el 40 a $\pm2$~minutos (0,4 puntos más el lunes y 0,2 menos el
martes) y pierde a $\pm5$~minutos. Con la prueba de la recalibración (los 7 últimos días, aprendiendo a la vez con todos los anteriores:
216\,534 pasos de todos los viajes, con autobús o sin él, hasta el 3 de noviembre), la media da 162~s de error medio y
el 51,8\,\% a $\pm2$~min, y el percentil 40, 160~s y el 54,2\,\%.

La corrección también da la llegada del autobús que espera en la cabecera (sección «\nameref{sec:cabecera}»). En esas
filas, los mismos dos días, el acierto a $\pm2$~min pasa del 53,1 al 58,3\,\% y del 59,7 al 64,3\,\%, y las veces que
el autobús pasa más de 2~minutos antes de lo mostrado, del 31,8 al 21,5\,\% y del 26,1 al 17,3\,\%. «Sin datos del
bus» (\texttt{likely\_departed}) sale en un 13--14\,\% más de filas, con la misma fiabilidad: el 96,5 y el 94,8\,\% de
esos viajes no se ven en todo el día, como con la media. Estas cifras usan una de cada tres instantáneas de 05:00 a
23:00, sin ver el día al aprender; la Tabla~\ref{tab:fiabilidad}, con una de cada seis instantáneas de 07:00 a 16:00 y
lo aprendido hasta ese mediodía, da el 97{,}8\,\% el 3 de noviembre. Los estados con el autobús en ruta y «Ya salió» no
dependen de la corrección y no cambian. Desde el 3 de noviembre de 2026, la corrección usa un percentil y no la media.

El percentil no queda fijo: cada noche la recalibración prueba el 30, el 35, el 40, el 45 y el 50
(\texttt{TRIP\_OFFSET\_RANKS}) con la misma prueba de los 7 últimos días, se queda el que más acierta a $\pm2$~min
(en caso de empate, el de menor error medio) y lo anota en el campo \texttt{rank} de \texttt{trip-offsets.json}. Con las muestras del 3
de noviembre, el 35 y el 40 empatan a 54,2\,\% y gana el 40 por un segundo de error medio (161 y 160~s); el 30 y el 50
dan 53,8 y 53,3\,\%. Si el tráfico cambia la forma de los desvíos, el percentil cambia con él.

\subsection{Tipo de Día, Festivos y Temporada}

El horario ya distingue los días: cada fecha circula con los viajes de su calendario (laborable, sábado, domingo,
festivo, temporada de verano). Lo que se aprende es el tráfico, y la Tabla~\ref{tab:tipo-dia} compara varias formas de
agruparlo. El único festivo entre semana de las muestras, el lunes 12 de octubre, se probó aprendiendo con los otros
26 días.

\begin{table}[ht]
\centering
\caption{Error Medio de la Capa 3 Corregida Según el Tipo de Día}\label{tab:tipo-dia}
\begin{tabular}{@{}lrr@{}}
\toprule
Tipo de día de la corrección & \shortstack{7 días\\de prueba} & \shortstack{Festivo del\\12 de octubre}\\
\midrule
Sin corrección (el festivo, con tramos de laborable) & 210~s & 172~s\\
Laborable o fin de semana por el día de la semana & 158~s & 196~s\\
\textbf{Laborable o fin de semana; el festivo, fin de semana} & \textbf{158~s} & \textbf{162~s}\\
Laborable; sábado; domingo o festivo & 159~s & 156~s\\
Lunes a jueves, viernes, sábado y domingo & 160~s & ---\\
Solo con los 14 últimos días & 163~s & ---\\
\bottomrule
\end{tabular}
\end{table}

Tratar el festivo según su día de la semana le aplica el retraso de un laborable con el tráfico de un domingo y lo deja
peor que sin corrección. Tratado como fin de semana, con los tramos y la corrección aprendidos juntos, baja a
162~s. Separar el sábado del domingo acierta algo más en el festivo, pero algo menos en el resto, con la mitad de datos
por grupo; separar el viernes o usar solo los días recientes empeora. La temporada del año (verano, vacaciones escolares)
no se puede medir con muestras de un solo mes. Mientras tanto, lo aprendido se recalcula cada día con las muestras de
los últimos 120 días y se va adaptando a la estación.

\section{Comparación con Otros Métodos}
\label{sec:comparacion}

El capítulo «\nameref{cap:marco}» describe cuatro familias de métodos; aquí se comparan tres: medias históricas, filtros de
Kalman y redes neuronales. El método de este trabajo es de la primera. Para saber cuánto se pierde por no usar las otras, se comparó
el método con ellas mediante el protocolo de la sección «\nameref{sec:eval-contraccion}»: los mismos siete días de
prueba, los mismos recorridos encadenados y los mismos horizontes. A diferencia del método, los rivales pueden usar
también lo que está pasando ese día. Cada uno predice la suma de los tramos siguientes desde el instante $t_0$ en que el
autobús sale del tramo inicial, y solo puede usar los días anteriores y los tramos del propio día que terminaron al
menos 60~s antes de $t_0$, que es lo que tarda el feed en publicar una posición nueva. Los métodos son estos:

\begin{itemize}
  \item \textbf{Horario}: la suma de los tiempos programados.
  \item \textbf{Método actual}: las medianas por tramo y franja contraídas con $n_0=5$.
  \item \textbf{Escalado}: el método actual multiplicado por un factor único, para separar la corrección del sesgo de
        todo lo demás.
  \item \textbf{Buses previos}: el método actual multiplicado, en cada tramo, por la mediana del cociente real/aprendido
        de los últimos tres autobuses que lo terminaron, contraída hacia 1 como si valiera $n_0$ autobuses.
  \item \textbf{Kalman por tramo}: un filtro de Kalman \citep{kalman1960} por tramo cuyo estado es la desviación
        logarítmica de ese día frente a lo aprendido. Se actualiza con cada autobús que termina el tramo, vuelve a cero
        con una constante de tiempo $\tau$ (un proceso de Ornstein--Uhlenbeck) y se proyecta a la hora en que el autobús
        llegará al tramo. Es la idea de \citet{shalaby2004}, que usan los tiempos de los autobuses anteriores.
  \item \textbf{Kalman del propio autobús}: el mismo filtro, con la desviación del propio autobús en los tramos que ya
        ha recorrido.
  \item \textbf{Red neuronal}: un perceptrón con dos capas ocultas de 32 y 16 neuronas, como los de \citet{chien2002}
        y \citet{jeong2005}, que predice la desviación logarítmica de cada tramo frente al método actual. Usa 28
        variables: el tiempo aprendido y el programado, con cuántos autobuses se aprendió, la distancia y la posición en
        el recorrido, la hora, la franja y el tipo de día, y lo que se sabe del día (el ritmo del propio autobús, los
        autobuses previos en el tramo y la desviación media reciente en la red y en la línea). Se entrena para cada día
        de prueba con los anteriores (entre 424\,000 y 632\,000 filas), con pérdida de Huber.
  \item \textbf{Red sin datos del día}: la misma red, solo con las variables que no dependen de lo que pasa ese día.
\end{itemize}

\noindent Los parámetros de los rivales se eligieron por búsqueda en rejilla con los siete días anteriores a los de
prueba (del 18 al 24 de octubre), minimizando la suma de los errores medios de los tres horizontes: factor 1,02 para el escalado;
ventana de 30~minutos y $n_0=4$ para los buses previos; $\tau=120$~min para el filtro por tramo y $\tau=30$~min para
el del propio autobús, con una varianza estacionaria del estado de 0,5 y 0,1 veces la del ruido de observación, respectivamente. Ningún
parámetro quedó en el borde de la rejilla. El método actual reproduce exactamente la columna $n_0=5$ de la
Tabla~\ref{tab:contraccion}, lo que comprueba que todos los métodos se evalúan sobre los mismos recorridos.

\begin{table}[ht]
\centering
\caption{Comparación con Otros Métodos con el Protocolo por Días}\label{tab:comparacion}
\begin{tabular}{@{}lccccc@{}}
\toprule
& \multicolumn{3}{c}{Error medio (s)} & $\pm120$~s & Días mejor\\
\cmidrule(lr){2-4}
Método & 4--8 min & 8--12 min & 12--20 min & 12--20 min & que el actual\\
\midrule
Horario                    & 153,9 & 233,5 & 329,7 & 24,0\,\% & 0\\
Método actual              & 42,0  & 61,7  & 88,2  & 74,1\,\% & ---\\
Escalado                   & 42,6  & 61,9  & 86,9  & 74,9\,\% & 2 · 4 · 5\\
Buses previos              & 41,6  & 60,7  & 86,5  & 75,0\,\% & 5 · 7 · 7\\
Kalman por tramo           & 41,7  & 60,7  & 86,4  & 75,0\,\% & 5 · 7 · 7\\
Kalman del propio autobús  & 40,4  & 59,3  & 84,6  & 76,0\,\% & 6 · 7 · 7\\
Red sin datos del día      & 40,7  & 58,8  & 82,5  & 77,2\,\% & 6 · 6 · 6\\
Red neuronal               & \textbf{39,5} & \textbf{57,5} & \textbf{81,7} & \textbf{77,4\,\%} & 7 · 7 · 7\\
\bottomrule
\end{tabular}
\nota{Error absoluto medio de la suma de tramos en 107\,052, 83\,336 y 113\,581 recorridos (días de prueba del 25 al 31
de octubre de 2026). «Días mejor que el actual»: de los siete días de prueba, en cuántos el método tiene menos error
medio que el actual, en cada uno de los tres horizontes. Las redes se entrenaron con la semilla 20261031; con otra semilla, la red neuronal da
39,5, 57,3 y 81,4~s, y la red sin datos del día, 40,8, 59,7 y 85,3~s.}
\end{table}

La Tabla~\ref{tab:comparacion} muestra que el método actual se queda con casi toda la mejora sobre el horario: el
error medio a 4--8~minutos baja de 154 a 42~s. Los rivales que usan el propio día mejoran algo más, en casi todos
los días, aunque poco. Los autobuses que acaban de pasar por los mismos tramos, solos o con un filtro de Kalman, restan
entre 0,3 y 1,8~s, un 1 o 2\,\%. El ritmo del propio autobús resta entre 1,6 y 3,6~s, un 4\,\%. La red neuronal, que
combina todo, resta entre 2,5 y 6,5~s, entre un 6 y un 7\,\%, y es la única que mejora los siete días en los tres
horizontes. Las diferencias con el método actual tienen intervalos de confianza del 95\,\% que no incluyen el cero
(\emph{bootstrap} de pistas; a 12--20~minutos, de $-7{,}9$ a $-5{,}2$~s para la red y de $-4{,}5$ a $-2{,}8$~s para el
Kalman del propio autobús).

Corregir solo el sesgo no explica la mejora. Sumar medianas subestima el recorrido largo, porque los tiempos de tramo
tienen cola a la derecha: a 12--20~minutos la mediana del error del método actual es de $-29$~s. El factor único la
lleva a $-11$~s, pero apenas cambia el error medio (de 88,2 a 86,9~s). La red sin datos del día mejora con las dos
semillas, pero mucho menos con una que con otra (82,5 y 85,3~s a 12--20~minutos), así que esa parte de la mejora no es
fiable. La de la red completa sí lo es: las dos semillas dan prácticamente lo mismo (81,7 y 81,4~s). La
sección «\nameref{sec:otros-enfoques}» discute qué significan estas cifras para el método.

\section{Trazados Corregidos y Espera en la Cabecera}
\label{sec:eval-trazados}

Los dos cambios del 3 de noviembre de 2026 (secciones «\nameref{sec:trazados-gtfs}» y «\nameref{sec:cabecera}») se
midieron con las capturas de ese día: diez bloques de unos 55~minutos, de 07:00 a 16:00, usando también el bloque
siguiente para ver el paso real de las llegadas del final de cada uno. Son 689\,292 llegadas con paso real conocido,
de las que 586\,425 quedan a menos de 20~minutos. Todas las variantes usan la misma base del horario de ese día y las mismas capturas; solo cambian el motor y la corrección de los
trazados:

\begin{itemize}
  \item A: la aplicación antes de los cambios;
  \item B: con los trazados corregidos;
  \item C: B y, además, el autobús en la cabecera saliendo a su hora en punto~\eqref{eq:eta-capa1};
  \item D: B y el autobús en la cabecera con la llegada de la capa 3~\eqref{eq:eta-cabecera}, que es la adoptada.
\end{itemize}

\begin{table}[ht]
\centering
\caption{Aciertos Dentro de $\pm2$~Minutos con los Trazados Corregidos y la Cabecera, por Horizonte}\label{tab:trazados}
\begin{tabular}{@{}lrcccc@{}}
\toprule
La parada queda a & Llegadas & A & B & C & D\\
\midrule
0--3 min   & 124\,872 & 94,1\,\% & 95,5\,\% & 95,6\,\% & 95,5\,\%\\
3--6 min   & 114\,153 & 87,6\,\% & 88,7\,\% & 89,3\,\% & 88,7\,\%\\
6--10 min  & 143\,882 & 80,8\,\% & 80,9\,\% & 81,0\,\% & 80,9\,\%\\
10--20 min & 203\,518 & 67,0\,\% & 67,0\,\% & 64,0\,\% & 67,0\,\%\\
\bottomrule
\end{tabular}
\nota{Las mismas 586\,425 llegadas a menos de 20~minutos en las cuatro columnas, de cualquier capa. A: antes de los cambios; B: trazados
corregidos; C: B con el autobús en la cabecera saliendo a su hora; D: B con la llegada de la capa 3 en la cabecera.
Sesgo medio a 10--20~min (media de $e$, negativa si el autobús llega después de lo mostrado): $-0{,}51$ (A),
$-0{,}54$ (B), $-1{,}02$ (C) y $-0{,}54$~min (D).}
\end{table}

La corrección de los trazados (B) mejora o iguala en todos los horizontes. Sus efectos mayores no
están en la tabla, que mide la misma lista de llegadas en las cuatro columnas, sino en cuántas de ellas salen en
directo: las llegadas calculadas con datos del autobús pasan de 470\,596 a 507\,023, y las que solo tienen el horario, de
115\,829 a 79\,402. Contando cada autobús con GPS en una de cada tres instantáneas, el 73{,}8\,\% quedaba situado antes de
los cambios y el 97{,}4\,\% después: el 81{,}2\,\% en ruta y el 16{,}2\,\% en la cabecera. Los demás aún no han
salido y están lejos de la primera parada (1{,}7\,\%), van junto al recorrido sin situarse (0{,}7\,\%) o han
terminado o van a más de 1~km de su recorrido (0{,}3\,\%). Con las capturas de ese día, las líneas que aprenden tramos
pasan de 95 a 98 de 99, con 2\,624 tramos distintos frente a 2\,404; las líneas 2610, 3642 y 3932 no aprendían ninguno.

Si se supone que el autobús de la cabecera sale a su hora en punto (C), se gana algo a menos de 6~minutos y se
pierden 3 puntos a 10--20~minutos, con el doble de sesgo: los autobuses salen después de su hora. Con la llegada de la capa 3 (D), la precisión es la de B
en todos los horizontes y las llegadas con datos del autobús suben a 569\,769, con solo 16\,656 de horario. Una captura en
directo de 18~minutos de esa misma tarde, con 12\,825 llegadas, da el mismo resultado: D iguala o mejora a A en cada horizonte,
con 12\,421 llegadas con datos del autobús frente a 10\,941.

\section{Fiabilidad de lo que Dice la Aplicación}
\label{sec:fiabilidad-estados}

Cada estado de la Tabla~\ref{tab:estados} afirma algo: que el autobús ya pasó, que espera en la cabecera, que viene. Con
las capturas del 3 de noviembre de 2026 se comprobó cada afirmación con lo que ocurrió después: para una de cada seis
instantáneas se comparó el estado de cada parada de cada viaje cercano (de hora y media antes a una hora después),
calculado con el registro de pasos reconstruido desde las instantáneas anteriores, con el paso real por cambio de
parada. No se cuentan las afirmaciones de la última hora capturada, porque no queda tiempo para ver qué pasó.

\begin{table}[ht]
\centering
\caption{Lo que Afirma Cada Estado y Cuántas Veces Fue Cierto}\label{tab:fiabilidad}
\begin{tabularx}{\textwidth}{@{}>{\raggedright\arraybackslash}p{0.3\textwidth}rX@{}}
\toprule
Estado y lo que afirma & Veces & Resultado\\
\midrule
En directo en la cabecera: el autobús espera allí su viaje & 20\,924 & 97{,}0\,\% hizo después ese viaje desde allí\\
Ya salió (\texttt{departed}) & 158\,833 & 99{,}9\,\% había pasado; 0{,}1\,\% pasó después\\
Estimado sin autobús (capa 3) & 170\,649 & 90{,}6\,\% vino un autobús ($\pm2$~min el 65\,\%, $\pm5$~min el 90\,\%); 5{,}2\,\% el viaje no se vio en todo el día\\
Aún sin autobús asignado & 2\,372 viajes & 96{,}7\,\% recibió autobús después, con una mediana de 7~min antes de su salida (el 17\,\% después de su hora)\\
Sin datos del bus (\texttt{likely\_departed}) & 2\,601 & 97{,}8\,\% de esos viajes no se vio en todo el día; 2{,}0\,\% sí había pasado\\
\bottomrule
\end{tabularx}
\nota{Una de cada seis instantáneas de diez bloques del 3 de noviembre de 2026, de 07:00 a 16:00. El paso real es el primer
instante en que el feed da como próxima parada una posterior.}
\end{table}

La fila «Estimado sin autobús» se midió con lo aprendido hasta ese mediodía, que ya incluía la mañana medida. Sin ver
el día al aprender, el mismo martes da el 58,2\,\% a $\pm2$~min, y el 60,7\,\% con la corrección por el percentil
40 (sección «\nameref{sec:correccion-percentil}»).

Los estados que afirman algo del autobús aciertan casi siempre. El de \texttt{likely\_departed} no se puede comprobar:
solo se aplica a viajes de los que no hay ningún dato, y casi ninguno aparece después. Por eso su texto pasó de
«Probablemente ya ha pasado» a «Sin datos del bus: puede que ya haya pasado». También cambió el del autobús en la
cabecera, de «En ruta» a «En la parada».

\section{Ideas Probadas y Descartadas}
\label{sec:descartadas}

La regla del método es que una idea entra solo si la medida mejora. La Tabla~\ref{tab:descartadas} recoge las que se
probaron y no entraron, con lo medido, para no repetirlas.

\begin{longtable}{@{}>{\raggedright\arraybackslash}p{0.33\textwidth}>{\raggedright\arraybackslash}p{0.455\textwidth}>{\raggedright\arraybackslash}p{0.135\textwidth}@{}}
\caption{Ideas Probadas y No Adoptadas}\label{tab:descartadas}\\
\toprule
Idea & Resultado medido & Decisión\\
\midrule
\endfirsthead
\toprule
Idea & Resultado medido & Decisión\\
\midrule
\endhead
\bottomrule
\endfoot
Corregir con la velocidad del último minuto del autobús (la esencia de un filtro) & 92,5\,\% → 91,9\,\% dentro de $\pm2$~min, y hasta 78,3\,\% con más peso & Descartada\\
Mezclar lo aprendido con el horario cuando hay pocas muestras & 83,5\,\% → 83,2\,\% (peso 1) y hasta 80,0\,\% (peso 8) & Descartada\\
Rellenar los tramos sin aprender con la proporción real/horario del propio viaje & 70\,\% → 67\,\% & Descartada\\
Franja de cada tramo a la hora prevista de paso por él (no la de la observación) & Sin diferencia en conjunto (71,5\,\% → 71,6\,\% a 12--20~min); mejor en los recorridos que cruzan de franja a 12--20~min, peor a 4--12~min. Solo el 3--9\,\% de los recorridos cruza de franja & Descartada\\
Festivos como fin de semana en los tiempos por tramo (capas 1 y 2) & El lunes festivo 12 de octubre se parece más a un laborable: 20,6~s de error con tiempos de laborable frente a 24,2~s con los de fin de semana (387 tramos) & Descartada\\
Corrección por viaje con el sábado y el domingo por separado & Festivo del 12 de octubre: 162 → 156~s, pero en los 7 días de prueba 158 → 159~s, con la mitad de datos por grupo & No adoptada\\
Corrección por viaje con el viernes aparte & 158 → 160~s; los viernes, 145 → 151~s & Descartada\\
Corrección por viaje solo con los 14 últimos días & 158 → 163~s & Descartada\\
Corrección por viaje por quintos o por parada en lugar de tercios & 158 → 159 y 162~s & Descartada\\
Mínimo de 1 pista por tramo & Una sola medida supera al respaldo por longitud tramo a tramo (24,1 frente a 40,1~s), pero de punta a punta solo 72 → 74\,\% a 10--20~min y nada en el resto & No adoptada\\
Sin datos de la franja, la mediana del mismo tipo de día antes que la de todas las horas & 19,5 → 20,0~s de error medio en las 7\,493 llegadas afectadas & Descartada\\
Curva no lineal dentro del tramo (exponente sobre la fracción recorrida) & Mejora del 1,3\,\% con exponente 0,9; el error grande está al principio del tramo y es variabilidad, no sesgo & Descartada\\
Reparto del tramo en la parada registrada distinto de la mitad & 0,5 es el óptimo: 83,0\,\% frente a 82,1\,\% (0,35) y 81,4\,\% (0,65) & Sin cambio\\
Instante exacto del cambio de próxima parada en la capa 2 & 91\,\% → 91\,\% en 5\,086 llegadas; exigiría guardar estado entre peticiones & Descartada\\
Tope de $k$ más alto o mínimos de pistas menores para el $k$ por línea & 128,3~s frente a 128,0~s con la mejor variante; diferencias de $\pm0{,}7$~s con intervalos del 95\,\% que incluyen el cero & Descartada\\
Contracción de la franja con $n_0=10$ & Mejor tramo a tramo, peor a 12--20~min (Tabla~\ref{tab:contraccion}) & Descartada\\
Autobús en la cabecera saliendo a su hora en punto & 67,0\,\% → 64,0\,\% dentro de $\pm2$~min a 10--20~min y sesgo de $-1{,}02$~min (Tabla~\ref{tab:trazados}); se usa la capa 3 & Descartada\\
Usar la parada del feed (\texttt{StopPointRef}) cuando no coincide con el orden $n$ & 2\,329 llegadas más en directo, pero 88,7\,\% → 88,5\,\% a 3--6~min y 15~ms más por petición & Descartada\\
Arrastrar al viaje siguiente de la misma clave «$\ell|\nu$» el retraso del autobús que aún va en el anterior & Donde se aplica, 60,9\,\% → 5,8\,\% dentro de $\pm2$~min a 20--60~min y $+22$~min de sesgo: solo el 11\,\% de esos viajes los hace el mismo autobús & Descartada\\
Seguir el autobús que desaparece del feed desde donde se le vio por última vez (hasta 30~min) & 39,8\,\% → 23,2\,\% dentro de $\pm2$~min y $-25$~min de sesgo en 22\,690 llegadas medidas tras volver al feed: el autobús que falta no suele estar recorriendo el viaje & Descartada\\
Sin autobús pasada su hora de salida, no dar una llegada anterior a salir ahora mismo de la cabecera & En esas filas, 26,5\,\% → 27,8\,\% (2 de nov.) y 32,3\,\% → 32,7\,\% (3 de nov.) dentro de $\pm2$~min; en el conjunto, $+0{,}1$ y $+0{,}02$ puntos. Muchos de esos autobuses salen a su hora sin aparecer en el feed & Descartada\\
Corregir la capa 3 con lo que se apartan ahora los demás viajes de la misma línea, sentido o red & 51,5\,\% → 51,9\,\% en el mejor caso (2 de nov.) y peor en casi todos: el desvío de un viaje un día es suyo (lo explica el autobús que lo hace), no de la línea & Descartada\\
\end{longtable}

Por el mismo criterio, entraron estas otras:

\begin{itemize}
  \item el mínimo de 2 pistas por tramo en lugar de 3: un 17\,\% más de tramos, con el mismo error;
  \item conservar la llegada cuando el retraso no es creíble;
  \item la parada registrada hasta tramos de 20~min: sin GPS, del 79,8\,\% al 83,0\,\% dentro de $\pm2$~min;
  \item el horario corregido de la capa 3: sin datos del autobús, duplica o triplica el acierto del horario publicado
        dentro de $\pm2$~min (Tabla~\ref{tab:capas});
  \item la corrección por viaje de la capa 3: de 210 a 158~s de error medio y del 47 al 54\,\% dentro de $\pm2$~min;
  \item en la corrección por viaje, el percentil 40 de los desvíos en lugar de la media: en las filas «Estimado» sin
        autobús de dos días no vistos al aprender, del 51,5 al 55,6\,\% y del 58,2 al 60,7\,\% dentro de
        $\pm2$~min, y un tercio menos de autobuses que pasan más de 2~min antes de la hora mostrada;
  \item en la capa 3, el festivo entre semana como fin de semana: de 196 a 162~s el 12 de octubre;
  \item el respaldo por longitud: con un tramo sin aprender, del 69\,\% al 76\,\%; con dos o más, del 49\,\% al
        59\,\%;
  \item la captura horaria (después, continua) y la contracción con $n_0=5$;
  \item descargar el feed cada 20~s en lugar de 45~s: con 45~s se perdía alguna versión (46 guardadas en 45~minutos,
        cuando a una versión cada 55~s se publican unas 49).
\end{itemize}

\chapter{Discusión}
\label{cap:discusion}

Este capítulo interpreta las medidas del capítulo «\nameref{cap:evaluacion}», las sitúa respecto a otros enfoques de la
literatura, recoge las limitaciones del trabajo y discute qué hace falta para aplicar el método a otras redes.

\section{Interpretación de los Resultados}
\label{sec:interpretacion}

Los resultados permiten distinguir de dónde sale la precisión del método.

\subsection{El Tramo en que Está el Autobús} Saber en qué tramo está el autobús aporta la mayor parte de la información. La capa 2, que solo conoce la
próxima parada y la hora del dato, queda a pocos puntos de la capa 1 con GPS (Tabla~\ref{tab:capas}). La posición
exacta dentro del tramo mejora la predicción, pero lo decisivo es partir de dónde está el autobús y no de la hora
publicada, como anticipa la proposición~\ref{prop:diferencias}.

\subsection{La Corrección del Horario} Corregir el error de forma del horario ya lo mejora mucho. La capa 3 no sabe
nada del autobús y, aun así, duplica o triplica los aciertos del horario publicado. Esa mejora solo puede venir de
sustituir las duraciones interpoladas por los tiempos aprendidos de cada tramo. Le quedaba, sin embargo, un sesgo:
llegaba de media 131~s antes que el autobús. La corrección por viaje, que usa lo que cada viaje suele hacer y nada del
autobús ese día, lo deja en 17~s (sección «\nameref{sec:eval-correccion}»). Quitar el sesgo de la media no es lo mismo
que acertar lo típico: con el percentil 40 de los desvíos en lugar de la media, la hora sin autobús acierta más y el
autobús pasa antes de lo mostrado un tercio menos de veces (sección «\nameref{sec:correccion-percentil}»). Lo que
queda es el desvío de cada viaje ese día, que se explica por el autobús que lo hace. Si ese desvío se conociera, el
acierto a $\pm2$~min de las filas «Estimado» sin autobús del 2 y el 3 de noviembre pasaría del 52--58\,\% (con la
media) al 81--85\,\%; pero ni los demás viajes de la línea ni la hora que ya ha pasado sin autobús lo predicen.

\subsection{La Cobertura de Tramos Aprendidos} La cobertura de tramos aprendidos es lo que más mueve la precisión. En
las llegadas a menos de diez minutos, con todo el camino aprendido se acierta el 93\,\% dentro de $\pm2$~minutos,
frente al 76\,\% con un tramo sin aprender y el 59\,\% con dos o más (Tabla~\ref{tab:cobertura} y el texto que la sigue; con 17 días, el 93, el 76 y el 63\,\%,
Tabla~\ref{tab:varios-dias-origen}). La curva de aprendizaje (Tabla~\ref{tab:curva}) indica que esa
cobertura sigue creciendo con los datos.

\subsection{El Error y la Distancia} El error crece con la distancia, y el acierto tiene un techo. Cada tramo suma su
propia variación, de modo que el
acierto baja de forma regular con el horizonte. En las llegadas a menos de diez minutos con todo el camino aprendido, el techo medido
está en torno al 93\,\% dentro de $\pm2$~minutos, limitado por la resolución del dato (una posición por minuto y $\pm30$~s en el paso real) y por lo que
ningún modelo puede prever con estos datos: semáforos, tráfico y esperas en parada.

\section{Relación con Otros Enfoques}
\label{sec:otros-enfoques}

El método pertenece a la familia de las medias históricas (capítulo «\nameref{cap:marco}»), con dos
diferencias: los tiempos se aprenden por tramo y franja en lugar de por línea, y la posición en tiempo real reinicia la
predicción en cada actualización. Otros trabajos han usado filtros de Kalman con datos de localización y de conteo de
pasajeros \citep{shalaby2004} o redes neuronales \citep{chien2002,jeong2005}; en el caso de \citet{jeong2005}, las
redes neuronales superaron a un modelo de datos históricos y a modelos de regresión en una línea de Houston.

Las cifras de este trabajo no se pueden comparar directamente con las de esos estudios: las redes, los datos, los
horizontes y las medidas de error son distintos. Por eso se implementaron esos métodos sobre las mismas muestras y con
el mismo protocolo por días (sección «\nameref{sec:comparacion}»). El resultado tiene dos lecturas. La primera es que
el modelo sencillo se queda con casi toda la mejora sobre el horario: de 154 a 42~s de error medio a 4--8~minutos,
mientras que la mejor alternativa, una red neuronal que además usa lo que pasa ese día, baja a 39,5~s. La segunda es
que esa diferencia, aunque pequeña, es real: la red reduce el error medio entre un 6 y un 7\,\% y mejora al método
todos los días de prueba en todos los horizontes, y un filtro de Kalman sobre el ritmo del propio autobús consigue algo más de la mitad de esa mejora con dos
parámetros.

Esto no contradice que corregir con la velocidad reciente del autobús empeorara la medida
(sección «\nameref{sec:descartadas}»). Esa corrección usaba la velocidad del último minuto, muy ruidosa, sobre la llegada
completa; el filtro usa la desviación de los tramos ya terminados, la suaviza y la olvida con una constante de 30~minutos.
Que los autobuses previos en el mismo tramo aporten tan poco (un 1 o 2\,\%) indica que, en esta red, la variación de
un día a otro en un tramo pesa menos que la de cada conductor y cada viaje, que es la que capta el ritmo del propio
autobús. La comparación se hizo sobre sumas de tramos, no sobre la llegada completa que ve el usuario; antes de
adoptar el filtro habría que medirlo así, con la regla de la sección «\nameref{sec:eval-contraccion}».

\section{Limitaciones}
\label{sec:limitaciones}

\begin{itemize}
  \item \textbf{Una muestra pequeña con las tres capas.} Las Tablas~\ref{tab:resultados} a~\ref{tab:cobertura}, las
        únicas que calculan cada llegada con las tres capas a la vez, salen de 14~minutos de feed de un laborable; sus cifras a horizonte
        largo tienen intervalos de confianza anchos. La medida principal se repitió con 17 días
        (Tabla~\ref{tab:varios-dias}), y con ellos el acierto a 10--20~minutos quedó en el 67\,\% y no en el 73\,\%, ya
        con llegadas de cualquier capa.
  \item \textbf{Un método que cambia.} Los 17 días miden la aplicación tal como estaba cada día, con calibraciones y
        reglas que cambiaron en ese periodo; no son la precisión de una versión fija del método.
  \item \textbf{Resolución de la verdad de referencia.} El paso real se conoce con $\pm30$~s. Las medidas dentro de
        $\pm1$~minuto están especialmente afectadas.
  \item \textbf{Definición de paso.} La verdad de referencia por cambio de parada mide cuándo el feed cambia de
        próxima parada, lo que puede ir algo por detrás del paso físico. Afecta por igual al método y al horario publicado.
  \item \textbf{Censura.} Solo se evalúan los autobuses cuyo paso se observa dentro de la ventana de captura; los muy
        retrasados pueden quedar fuera.
  \item \textbf{Dependencia temporal.} Las predicciones de un mismo autobús en instantáneas sucesivas están
        correlacionadas; el número efectivo de observaciones independientes es menor que el número de predicciones.
  \item \textbf{Temporada y festivos.} Las muestras van del 6 de octubre a los primeros días de noviembre de 2026 y
        solo tienen un festivo entre semana. La regla de
        los festivos se apoya en ese día, y el efecto de la temporada (verano, vacaciones) está sin medir.
  \item \textbf{Lo que el feed no dice.} Un viaje sin autobús en el feed sale con la capa 3 aunque no se haga, y un
        refuerzo que no está en el horario no sale en las paradas. Según las capturas del 3 de noviembre de 2026, en
        torno al 4\,\% de los viajes que deberían estar en marcha no se pueden interpretar: los que salen tarde y los
        que no aparecen en todo el día se ven igual en el momento (sección «\nameref{sec:feed}»).
  \item \textbf{De madrugada.} De 23:58 a 05:00 el feed de Bizkaibus repite la lectura de las 23:58 y los autobuses
        nocturnos salen con la capa~3 o con el horario (sección «\nameref{sec:mas-de-24h}»). Esas llegadas no están
        medidas, y la parte de un viaje posterior a la medianoche no aporta muestras al aprendizaje.
  \item \textbf{Calidad del horario.} El método depende de que trazado y paradas casen. En el horario de Bizkaibus del
        3 de noviembre de 2026, 1\,417 de los 6\,130 viajes (el 23\,\%) tenían alguna parada a más de 200~m del
        punto del trazado que marcan sus metros, 670 de ellos (el 11\,\% del total) con el trazado al revés. La
        corrección de la sección «\nameref{sec:trazados-gtfs}» lo arregla en la construcción, pero otra red puede
        traer errores distintos, que habría que buscar antes de medir el método.
  \item \textbf{Una sola red.} Todos los resultados son de Bizkaibus. La generalización a otras redes es una hipótesis
        razonable (sección «\nameref{sec:aplicabilidad}»), no un resultado.
  \item \textbf{Evaluación hecha por el autor.} El método lo ha evaluado quien lo ha diseñado, con sus propias
        herramientas. Los vectores de prueba y la implementación de referencia permiten comprobar que el método descrito
        es el que se ha medido, pero no sustituyen a una replicación independiente con otros datos.
  \item \textbf{Sin revisión por pares.} Este documento no ha pasado por una revisión externa.
\end{itemize}

\section{Límites de la Infraestructura}
\label{sec:infraestructura}

Todo el sistema funciona sin un servidor propio y casi sin coste, y eso condiciona el método tanto como los datos.

\begin{itemize}
  \item \textbf{La aplicación, sin estado.} La aplicación y la API corren en una plataforma sin servidor (Vercel): cada
        petición es una función independiente que no recuerda nada de la anterior, salvo una caché corta en disco y un
        almacén clave-valor externo (Upstash Redis) que se paga por comando. Cada función admite 250~MB, de los que la
        aplicación usa 208, y cada red nueva suma su base de datos del horario.
  \item \textbf{El aprendizaje, en una plataforma de pruebas.} La captura y la recalibración corren en una plataforma
        de integración continua (GitHub Actions) pensada para compilar y probar, no para trabajar sin parar: cada trabajo
        dura como mucho seis horas y los disparos programados se retrasan o se saltan (sección «\nameref{sec:ciclo}»).
        Hasta el 31 de octubre de 2026 solo se ejecutaron 163 de unas 600 capturas previstas.
  \item \textbf{Los datos, en ficheros.} Las muestras acumuladas no están en una base de datos, sino en ficheros JSON
        Lines guardados como adjuntos del repositorio. Cada recalibración los descarga enteros, los recalcula en memoria
        y los vuelve a subir. Por eso existen los topes de la sección «\nameref{sec:retencion}»: 120~días, 300 muestras
        por tramo y tipo de día, y 500\,000 muestras de ritmo, con ficheros de unos 100--120~MB y menos de 2~GB de memoria.
  \item \textbf{Una vez al día.} Los parámetros se recalculan una vez al día y producción usa parámetros de entre 23 y
        47~horas de antigüedad: un desvío o unas obras de hoy no entran en lo aprendido hasta uno o dos días después.
\end{itemize}

\noindent Estos límites explican varias limitaciones del método:

\begin{itemize}
  \item \textbf{Memoria corta.} Con 120~días y un tope por tramo no hay un año de historia. No se puede aprender una
        fiesta con la misma fiesta del año anterior (Navidad, Semana Santa, Aste Nagusia), ni el verano o las
        vacaciones escolares, y los días especiales se mezclan con los normales al aprender.
  \item \textbf{Grupos gruesos.} Con pocas muestras por tramo, lo aprendido se agrupa en seis franjas y dos tipos de
        día. Con más historia se podría separar por día de la semana, por hora o por fecha especial sin quedarse sin
        datos en cada grupo.
  \item \textbf{Sin otras fuentes.} No se cruzan la meteorología, el calendario escolar, los eventos (un partido o un
        concierto) ni las incidencias publicadas, porque no hay dónde guardar su historia junto a la de los pasos.
  \item \textbf{Sin modelos más capaces.} La red neuronal de la sección «\nameref{sec:comparacion}», que usa además
        lo que está pasando ese día, reduce el error medio de la suma de tramos entre un 6 y un 7\,\% frente al método. Se descartó para producción porque lo
        sencillo se puede auditar, pero también porque entrenarla y servirla en directo exige un servidor que no existe.
\end{itemize}

\noindent Con inversión en infraestructura se podría quitar cada uno de estos límites:

\begin{itemize}
  \item un servidor que capture el feed sin pausa, sin huecos ni capturas saltadas;
  \item una base de datos de series temporales que guarde años de pasos sin tope y permita consultar solo lo necesario,
        en lugar de descargar y recalcular ficheros enteros;
  \item aprender por fecha especial (festivos, Navidad, fiestas de cada pueblo, vacaciones escolares) comparando con
        el mismo día del año anterior, y cruzar los pasos con la meteorología, los eventos y las incidencias;
  \item recalibrar cada hora o de forma continua, de modo que un desvío entre en lo aprendido en horas y no en días;
  \item servir en directo modelos más capaces, como la red neuronal ya medida, e intervalos de predicción;
  \item añadir redes sin el límite de tamaño de cada función.
\end{itemize}

\noindent Ninguna de estas mejoras está medida. Son la dirección que señalan las medidas de este trabajo: la cobertura
de tramos aprendidos es lo que más mueve la precisión (sección «\nameref{sec:cobertura-tramos}»), y más historia y más fuentes
dan más cobertura y grupos más finos. Cuánto mejoraría la precisión en festivos o en verano solo se sabrá con esos
datos y con el mismo protocolo de medida.

\section{Aplicabilidad a Otras Redes}
\label{sec:aplicabilidad}

El método se diseñó para Bizkaibus, pero casi nada de él es específico de esa red. Esta sección separa lo que es
general de lo que hay que adaptar y da una lista de pasos para poner el método en marcha en otra ciudad. Describe
además las otras dos situaciones que ya resuelve la aplicación: redes que publican su propio retraso, que se tratan en
el motor, y redes que publican su propia predicción, que se tratan fuera de él.

\subsection{Requisitos}
\label{sec:requisitos}

El método necesita, por orden de importancia:

\begin{enumerate}
  \item \textbf{Un horario GTFS} con viajes, paradas y horas de paso. Sin él no hay recorrido sobre el que situar al
        vehículo.
  \item \textbf{Un feed en tiempo real que identifique el viaje} de cada vehículo y la \textbf{próxima parada}, con la
        hora del dato. Con eso funcionan las capas 2 y 3 y el aprendizaje queda limitado a lo que se pueda reconstruir.
  \item \textbf{La posición GPS} de cada vehículo y la hora en que se tomó. Con ella funcionan la capa 1 y el
        aprendizaje.
  \item \textbf{Trazados} (\texttt{shapes.txt}). Si no traen los metros de cada parada
        (\texttt{shape\_dist\_traveled}), la construcción de la base los mide sobre el trazado; sin trazados, la
        localización usa rectas entre paradas, con un pasillo más ancho, y no hay respaldo por longitud.
  \item \textbf{Una forma de capturar} el feed periódicamente y de guardar las muestras: aquí, un trabajo
        programado de una plataforma de integración continua y un fichero adjunto a una publicación del repositorio.
\end{enumerate}

\noindent Un operador que no publique retrasos ni predicciones, como Bizkaibus, es precisamente el caso para el que el
método aporta más. Uno que sí los publique puede usarlos directamente (sección «\nameref{sec:otras-redes}»).

\subsection{Qué Hay que Adaptar}
\label{sec:adaptar}

\subsubsection{La Lectura del Feed y el Emparejamiento} Es la única parte que depende del formato concreto del operador:
cómo se extraen la línea, el número de viaje y la salida de cabecera de la referencia del viaje
(expresión~\eqref{eq:regex}) y cómo se encuentra el viaje del GTFS (sección «\nameref{sec:emparejamiento}»). Si la referencia
del feed coincide con el \texttt{trip\_id} del GTFS, basta la tabla de alias. El motor deduce además de qué día es cada
viaje por la salida de cabecera del feed~\eqref{eq:dia-viaje}, que en Bizkaibus va en horas del horario. Los lectores
que la toman de \texttt{OriginAimedDepartureTime}, en hora del reloj (Bilbobus, Lurraldebus), considerarían de hoy un
viaje que empezara pasadas las 24:00; en esas dos redes no hay ninguno.

\subsubsection{Las Franjas} Las seis franjas por tipo de día responden a la demanda de Bizkaia. Una ciudad con otra
distribución de la hora punta debería revisarlas, con el mismo protocolo de evaluación por días.

\subsubsection{Las Constantes Medidas} Las constantes del Apéndice~\ref{ap:constantes} son de dos clases. Unas son
geométricas o de diseño y no deberían cambiar (la ventana de tramos y la regla de desempate). Otras se midieron en Bizkaibus y conviene volver a medirlas: el pasillo de 50~m depende de la calidad de los
receptores GPS, los 20~minutos de la capa 2 dependen de la longitud típica de los tramos, el tope de 45~minutos de los
retrasos reales de la red y los 60~s de salida de cabecera de la práctica de los conductores. La mitad del tramo de la capa 2 también se midió:
fue el mejor reparto entre 0,35 y 0,65.

\subsubsection{Sin Muestras} En una red nueva, sin muestras propias (la semilla del repositorio solo sirve para las
paradas de Bizkaibus), la tabla de tiempos por tramo está vacía y el modelo se reduce al factor de
ritmo, que sin muestras vale 1: las diferencias de horario de la proposición~\ref{prop:diferencias}, sin corrección de
forma. Solo con eso ya mejora mucho al horario publicado, porque elimina el retraso acumulado. Con unos días de capturas aparecen los
primeros tramos aprendidos y el modelo mejora sin intervención.

\subsection{El Mismo Motor en Otras Redes de Autobús}
\label{sec:lurraldebus}

El motor se reutiliza por herencia en dos redes más. Lurraldebus (Gipuzkoa) recibe su información en tiempo real en
varios feeds SIRI-VM, uno por operador; Bilbobus (los autobuses urbanos de Bilbao) la recibe en un solo feed del
Consorcio de Transportes de Bizkaia. En las dos, una clase derivada del emparejador solo cambia la red a la que se
aplica y los ficheros con los datos aprendidos (factor de ritmo, tiempos por tramo y correcciones por viaje); el
cálculo es el mismo código. La lectura del feed sí es distinta: un lector genérico toma la referencia del viaje tal
cual la publica el operador y la empareja con los alias del horario, con el nombre del operador delante cuando hay
varios. Las dos redes están activadas, pero sus feeds llegan vacíos durante largos periodos: entonces la aplicación
muestra el horario oficial sin afirmar que sigue a ningún vehículo, y el aprendizaje de tramos no tiene con qué
empezar hasta que los feeds traigan vehículos de forma regular. Sus GTFS no traen los metros de cada parada; la
construcción los mide sobre el trazado (Tabla~\ref{tab:sin-metros}), así que, cuando lleguen vehículos, se situarán
con el mismo pasillo de 50~m que Bizkaibus y no con la recta entre paradas.

\subsection{Redes que Publican Retraso o Predicción}
\label{sec:otras-redes}

\subsubsection*{Retraso Declarado}

En el código, «red de autobús» es la red para la que se configuró el emparejador; las demás, sean o no de autobuses,
usan este camino. En ellas las capas 2 y 3 y el contraste no se aplican, y la entrada se considera en directo si no tiene
posición o si la tiene y es válida. La llegada se calcula a partir del retraso declarado por el operador, $\delta_{\mathrm{decl}}$:
\begin{equation}
\ETA_{\mathrm{plano}}=a_\dest+\delta_{\mathrm{decl}}.
\end{equation}
Si la entrada no tiene posición pero sí parada actual y orden $n$, y el horario da la llegada $a_n$ a esa parada, se
calcula también una llegada por diferencias desde ahora,
\[
\ETA_{\mathrm{pos}}=\tnow+(a_\dest-a_n),
\]
y se usa en lugar de la plana si $a_\dest-a_n\ge0$ y $|\ETA_{\mathrm{pos}}-\ETA_{\mathrm{plano}}|\le900$~s
(\texttt{POSITION\_SANITY\_SECONDS}). Con posición válida se usa la capa 1 tal cual. En estas redes no se aplica el
tope de 45~minutos ni el estado \texttt{estimated}.

\subsubsection*{Predicciones de Llegada (GTFS-Realtime Trip Updates)}

Metro Bilbao, Euskotren, los tranvías de Bilbao y Vitoria-Gasteiz y Álavabus publican, para cada vehículo, la hora
prevista en cada parada que le queda; Renfe, Tuvisa y Dbus, un retraso por viaje. Ahí no hace falta estimar: basta con
emparejar cada vehículo con su viaje del horario y comprobar que lo que dice el operador es un dato en directo. Ese
trabajo no aprende nada, así que está en la aplicación y no en el motor.

\paragraph{Emparejar por el Viaje} Todos esos feeds traen el identificador del viaje del horario (con el horario del
4 de noviembre de 2026: los 71 viajes de Metro en servicio, 188 de Euskotren, 46 y 73 de los tranvías), y se empareja
por él. La hora en una parada es la del feed; si el feed se salta esa parada, la de la siguiente que trae, menos lo que
se tarda hasta ella según el horario; con un retraso por viaje, la del horario más ese retraso. El retraso que se
muestra es la diferencia con la hora publicada, redondeado a cero por debajo de 60~s (\texttt{ON\_TIME\_SECONDS}).
Las horas del feed son del reloj y las del horario pueden pasar de las 24:00: cada viaje guarda los días enteros que
separan unas de otras (24~h para un tren anotado a las 26:00 que pasa a las 02:00; 0 el resto del día), y todas las
comparaciones los descuentan. Sin eso, ese tren salía con $\pm1\,440$~min de retraso.

La versión anterior emparejaba por paradas y hora: cada tren del feed con la fila del horario más cercana dentro de
10~minutos, asignando de menor a mayor diferencia. Se suponía que con trenes separados varios minutos la ambigüedad era
rara. El 4 de noviembre, con retrasos generalizados de 15 a 20~minutos en Metro, no lo era: cada tren se pegaba a la
salida siguiente, la aplicación mostraba «+2~min» o «4~min adelantado» a trenes que iban 18~minutos tarde y dejaba
salidas sin tren como «solo horario». Ese emparejamiento queda solo de reserva, para cuando ningún viaje del feed esté
en el horario (por ejemplo, un horario recién cambiado).

\paragraph{Comprobaciones del Directo} Antes de mostrar nada se comprueba, con los datos de cada consulta:
\begin{itemize}
  \item \emph{Un feed que copia el horario.} Si todas las horas del feed son las del horario al segundo, no es
        directo y sale el horario. Es el caso de los trenes de Euskotren (2\,504 horas, ninguna distinta); los tranvías,
        en el mismo feed, sí traen horas distintas.
  \item \emph{Unas horas que contradicen al GPS.} Si hay al menos 3~vehículos con posición y, en la mayoría de ellos,
        la hora del feed en la parada por la que van se aparta más de 20~minutos de la hora actual
        (\texttt{VEHICLE\_DISAGREEMENT\_SECONDS}), las horas del feed no valen. Es el caso de Álavabus, que las da con una
        hora de más.
  \item \emph{Un viaje que aún no ha salido.} Si no ha salido de su origen y el operador solo repite su hora
        publicada (Renfe, hasta una hora antes, con retraso cero), sale con el horario.
\end{itemize}

\paragraph{Vehículos con Posición y sin Predicción} Su retraso es el que llevan en la parada por la que van (nunca un
adelanto) y se aplica a las paradas que les quedan. No cuentan los que están en la última parada de su viaje ni los que
darían más de 45~minutos (\texttt{MAX\_GPS\_DELAY\_SECONDS}), que suelen ser vehículos parados o que aún hacen el viaje
anterior.

\paragraph{Predicción Vencida} Si un vehículo aún no ha llegado a la parada a la que va y la hora que le daba el
operador ya pasó hace más de 60~s, va al menos ese tiempo más tarde: se suma a esa parada y a las siguientes, y crece
mientras no avance. Con más de una hora de desfase no se extrapola (\texttt{MAX\_STALE\_PREDICTION\_SECONDS}).

\paragraph{Salidas sin Datos y Salidas Pasadas} En los trenes, una salida que ya debería estar en la estación y no está en el
feed va al final como «sin datos del tren». Se miran también las salidas de la última hora, para que un tren muy
retrasado siga en la lista, y se retiran las que ya salieron: las que están en directo, 45~s después de su llegada, y las
programadas, 60~s después.

\paragraph{Vigilancia} El servidor publica para cada red un diagnóstico (viajes del feed, cuántos encuentra en el
horario, cuántos deberían ir en marcha según el horario, si copia el horario o contradice al GPS y la mediana y el
percentil 95 de los retrasos que muestra), y un proceso lo revisa cada hora: si un feed deja de responder, si una red
deja de encontrar sus viajes o si los retrasos no son creíbles (mediana de más de 30~minutos o percentil 95 de más de
60), avisa. Un feed vacío, o uno cuyos viajes no están en el horario, solo cuenta como fallo si según el horario van en
marcha al menos tres viajes de la red; el vacío, además, si la red no es de las que lo publican vacío a ratos aunque
haya servicio (Bilbobus, Lurraldebus y Álavabus). Un feed con vehículos pero ninguno con hora de hoy, como el de Bizkaibus
de madrugada, queda como aviso. Así se detectaron las horas desfasadas de Álavabus, vehículos ya llegados que seguían
contando y trenes de madrugada con $\pm1\,440$~min de retraso.

\subsection{Pasos Para Otra Ciudad}
\label{sec:pasos}

\begin{enumerate}
  \item Comprobar qué publica el operador. Con Trip Updates, basta el emparejamiento de la sección anterior y sus
        comprobaciones. Con SIRI-VM
        (o GTFS-Realtime Vehicle Positions), hace falta el motor de este trabajo.
  \item Construir la base del horario con la fusión de variantes y los alias (sección «\nameref{sec:alias}») y los trazados
        con distancias.
  \item Adaptar la lectura del feed y el emparejamiento, y comprobar con capturas reales que todos los vehículos del
        feed encuentran su viaje.
  \item Poner en marcha la captura periódica y la recalibración diaria. Sin calibrar, el modelo tiene poco con qué
        corregir el horario.
  \item Medir con el protocolo del capítulo «\nameref{cap:evaluacion}» sobre capturas de la propia red antes de mostrar
        nada como «en directo», y volver a medir las constantes de la sección «\nameref{sec:adaptar}».
  \item Comprobar que una implementación propia reproduce los vectores del Apéndice~\ref{ap:vectores}.
\end{enumerate}

\chapter{Conclusiones y Líneas Futuras}
\label{cap:conclusiones}

\section{Resumen}

Esta memoria ha descrito un método para predecir la llegada de autobuses a sus paradas en una red cuyo operador
publica la posición de sus vehículos pero ningún retraso, y cuyo horario intermedio es una interpolación a velocidad
constante. El método se resume en cuatro pasos:

\begin{enumerate}
  \item \textbf{Situar} el GPS sobre el trazado oficial, en una ventana alrededor de la próxima parada que declara el
        feed y con una regla de desempate para los recorridos que vuelven sobre sí mismos, y traducir esa posición a
        la hora a la que el horario pone allí al autobús.
  \item \textbf{Sumar} el tiempo real que cuesta recorrer lo que falta, tramo a tramo, con tiempos aprendidos por
        franja, contraídos hacia el tiempo del tramo a cualquier hora, y con una cascada de respaldo por longitud y por
        factor de ritmo.
  \item \textbf{Elegir} en cada predicción la mejor fuente fiable entre el GPS, la última parada registrada y el
        horario corregido, con reglas explícitas para abstenerse y sin afirmar nunca un retraso que no se pueda
        sostener.
  \item \textbf{Aprender y medir} cada día con el propio feed, y cambiar los parámetros solo si el error baja en
        recorridos que no se usaron para aprenderlos; las reglas nuevas, solo si mejoran en días no vistos al aprender.
\end{enumerate}

\section{Resultados}

En 2,21~millones de llegadas a menos de veinte minutos, recogidas en 17 días, la llegada mostrada cae dentro de $\pm2$~minutos
en el 96\,\% de los casos cuando el autobús está a menos de tres minutos y en el 67\,\% entre diez y veinte minutos,
frente al 24\,\% y el 25\,\% del horario publicado, con intervalos de confianza de, como mucho, tres puntos. En una muestra
de 14~minutos con GPS, la llegada no tiene sesgo y acierta el 94\,\% a menos de tres minutos. Sin GPS, con solo la
parada registrada, el acierto es el mismo a menos de tres minutos (94\,\%) y baja seis puntos entre diez y veinte
minutos (del 73 al 67\,\% en esa muestra). Sin ningún dato del
vehículo, el horario corregido duplica o triplica el acierto del horario publicado, y la corrección por viaje baja su
error medio de 210 a 158~s; con el percentil 40 de los desvíos en lugar de la media, la hora sin autobús acierta
entre 2 y 4 puntos más a $\pm2$~minutos en días no vistos al aprender. Cuando todo el camino tiene tiempos
aprendidos, el 93\,\% de las llegadas a menos de diez minutos cae dentro de $\pm2$~minutos, el techo que permite
para ellas la resolución de este feed. Con el protocolo de evaluación por días, un filtro de Kalman sobre el ritmo del propio autobús y
una red neuronal reducen el error de las sumas de tramos un 4\,\% y entre un 6 y un 7\,\%, respectivamente.

\section{Lecciones}

\subsection{Lo Único Vivo Es la Posición} Antes de modelar nada hubo que comprobar qué datos eran fiables. El retraso
declarado no existía y el horario intermedio era una interpolación. Todo el diseño parte de usar la posición y
diferencias de horario, que hacen irrelevante el retraso acumulado (proposición~\ref{prop:diferencias}).

\subsection{Revisar los Datos de Entrada} Durante semanas, 670 de los 6\,130 viajes (el 11\,\%, en 25 líneas) no se
pudieron situar con el GPS porque el horario oficial dibujaba sus trazados al revés, y esos viajes tampoco aprendían.
No era un fallo del método, sino de los datos, y solo se vio al revisar uno a uno los autobuses que no salían en
directo. Corregir los trazados en la construcción y situar al autobús que espera en la cabecera subió los autobuses
situados del 73{,}8 al 97{,}4\,\%, más que cualquier mejora del modelo.

\subsection{Medir Antes de Adoptar} Muchas ideas razonables empeoraron la medida: corregir con la velocidad del
último minuto,
mezclar lo aprendido con el horario, tratar los festivos como fin
de semana en los tiempos por tramo, suponer que el autobús de la cabecera sale a su hora, arrastrar el retraso del
viaje anterior o seguir al autobús que desaparece del feed; calcular la franja de cada tramo a la hora prevista de paso no
cambió nada. En la capa 3 corregida, en cambio, tratar los
festivos como fin de semana sí mejora, mientras que la primera versión de la corrección, que los trataba por su día de
la semana, los dejaba peor que sin corregir: cada caso se midió por separado. La regla de adoptar solo lo que mejora la medida, por días y de punta a punta, ha sido más útil que
cualquier intuición sobre el tráfico.

\subsection{Abstenerse Es Parte de la Precisión} Mostrar «Estimado» o «Solo horario» cuando no hay datos, ocultar un
retraso no creíble y desconfiar de dos fuentes que se contradicen no aumentan los aciertos medidos, pero evitan los
errores que más daño hacen a la confianza del usuario.

\subsection{Lo Sencillo y Auditable Compite} Medianas por tramo, una contracción con una sola constante, una cascada
de cuatro niveles: cada parámetro tiene un significado físico y se puede comprobar a mano. Frente a una red neuronal
entrenada con lo mismo y con lo que pasa cada día, se pierde entre un 6 y un 7\,\% de error medio; a cambio, cualquier
cifra que muestre la aplicación se puede explicar.

\section{Líneas Futuras}

\begin{itemize}
  \item \textbf{El ritmo del propio autobús.} El filtro de Kalman sobre la desviación del propio autobús mejora las
        sumas de tramos un 4\,\% con dos parámetros (sección «\nameref{sec:comparacion}»). El paso siguiente es medirlo
        sobre la llegada completa, con la regla de adopción de la sección «\nameref{sec:eval-contraccion}».
  \item \textbf{Capas a lo largo de los días.} La captura continua guarda desde el 31 de octubre de 2026 los errores de
        cada bloque por capa, horizonte, franja y tipo de día. Con semanas de datos, las Tablas~\ref{tab:capas}
        y~\ref{tab:cobertura} se podrán repetir con intervalos de confianza estrechos.
  \item \textbf{Temporada y tipo de día.} Con meses de muestras se podrá medir si conviene separar el verano y las
        vacaciones escolares, o el sábado del domingo, en la corrección por viaje y en los tiempos por tramo.
  \item \textbf{Fechas especiales.} Los tiempos aprendidos distinguen laborable y fin de semana, pero no los días de
        fiesta, de eventos o de Navidad. Marcar esas fechas para que no cuenten al aprender evitaría que contaminen los días
        normales; aprenderlas aparte exige capturas del año anterior.
  \item \textbf{Viajes sin autobús.} Distinguir en la aplicación el viaje que ya debería haber salido y no tiene
        autobús en el feed (en torno al 4\,\% de los que deberían estar en marcha), en lugar de darlo como estimado.
        Antes de su hora, la ficha ya dice que aún no tiene autobús asignado.
  \item \textbf{De madrugada.} De 23:58 a 05:00 el feed de Bizkaibus no se actualiza y los nocturnos salen sin directo
        (sección «\nameref{sec:mas-de-24h}»). Si el operador lo mantiene activo de noche, el emparejador de ayer los situará sin
        cambios; para medir esas llegadas y aprender de ellas, la captura tendría que separar también los viajes de ayer.
  \item \textbf{La secuencia de cada vehículo.} En una cabecera, el autobús que empieza un viaje suele ser uno que
        acaba de terminar otro, y su retraso pasa al siguiente. La clave del feed no lo dice (solo el 11\,\% de los
        viajes consecutivos de una clave los hace el mismo autobús), pero las capturas guardan qué vehículo hace cada viaje
        cada día. Si esa secuencia se repite de un día a otro, se podría seguir al autobús en su viaje anterior y
        estimar la salida real; hay que medir antes si se repite.
  \item \textbf{Revisión externa.} Que otras personas repliquen el método con la implementación de referencia y lo
        apliquen a otra red.
  \item \textbf{Intervalos de predicción.} Las muestras permiten estimar, además de la mediana de cada tramo, su
        dispersión. Sumando varianzas por tramo se podría mostrar un intervalo («llega en 6--8 minutos») con una
        cobertura medida.
  \item \textbf{Más capacidad técnica.} Un servidor propio y una base de datos con años de historia quitarían los
        topes de la sección «\nameref{sec:infraestructura}»: aprender festivos, fiestas y temporadas con el año
        anterior, cruzar meteorología y eventos, recalibrar cada hora y servir modelos más capaces. Hay que medir
        cuánto mejora cada paso.
  \item \textbf{Más fuentes de datos.} Pasar del 93\,\% en las llegadas a menos de diez minutos con todo el camino
        aprendido exige información que este feed no tiene: posiciones
        más frecuentes, apertura de puertas o conteo de viajeros.
  \item \textbf{Incidencias.} Los cortes y desvíos publicados por el operador no entran hoy en el cálculo hasta que se
        notan en la posición. Incorporarlos exige medir antes cómo afectan.
  \item \textbf{Planificación de viajes.} Las llegadas estimadas pueden alimentar un buscador de rutas con transbordos
        (por ejemplo, con el algoritmo RAPTOR), de modo que el viaje propuesto use las horas reales y no las
        publicadas.
  \item \textbf{Más redes.} Aplicar el método a otras ciudades y comprobar cuánto de lo medido aquí se mantiene.
\end{itemize}

\section{Reproducibilidad}

Para que cualquiera pueda reproducir el método, esta memoria incluye todas las constantes
(Apéndice~\ref{ap:constantes}), el pseudocódigo (Apéndice~\ref{ap:algoritmos}), los formatos de datos
(Apéndice~\ref{ap:formatos}), vectores de prueba generados con el código de producción (Apéndice~\ref{ap:vectores}) y
una implementación de referencia independiente que los reproduce exactamente (Apéndice~\ref{ap:referencia}). Una
implementación nueva, en cualquier lenguaje, puede darse por correcta cuando reproduce todos los vectores al segundo.


% =====================================================================================================================
\clearpage
\phantomsection
\addcontentsline{toc}{chapter}{Referencias}
\printbibliography[title=Referencias,heading=bibliography]

\apendices
\chapter{Constantes}
\label{ap:constantes}

Este apéndice reúne todas las constantes del método, con su nombre en el código, su valor y la sección donde se explica. Las marcadas en la
columna «Med.» se fijaron midiendo su efecto en Bizkaibus y conviene volver a medirlas en otra red
(sección «\nameref{sec:adaptar}»).

% Los nombres del código van en monoespaciada de 10 puntos, como pide APA para el código.
{\setlength{\tabcolsep}{4pt}\setstretch{1.15}
\begin{longtable}{@{}>{\raggedright\arraybackslash\fontsize{10}{12}\selectfont}p{0.30\textwidth}>{\raggedleft\arraybackslash}p{0.11\textwidth}>{\raggedright\arraybackslash}p{0.27\textwidth}>{\raggedright\arraybackslash}p{0.19\textwidth}c@{}}
\caption{Constantes del Método}\label{tab:constantes}\\
\toprule
Nombre en el código & Valor & Significado & Sección & Med.\\
\midrule
\endfirsthead
\toprule
Nombre en el código & Valor & Significado & Sección & Med.\\
\midrule
\endhead
\bottomrule
\endfoot
\multicolumn{5}{@{}l}{\emph{Construcción de la base (trazados)}}\\
\texttt{MAX\_STOP\_OFFSET\_METERS} & 100 m & Parada mal medida si cae a más de esta distancia del punto que marcan sus metros & «\nameref{sec:trazados-gtfs}» & \\
\texttt{ROAD\_VS\_STRAIGHT\_RATIO} & 3 & Búsqueda de la parada siguiente: veces la línea recta & «\nameref{sec:trazados-gtfs}» & \\
\texttt{SEARCH\_SLACK\_METERS} & 500 m & Margen de esa búsqueda & «\nameref{sec:trazados-gtfs}» & \\
\texttt{EARLIER\_PASS\_MARGIN\_METERS} & 30 m & Ventaja exigida a un paso posterior por la misma zona & «\nameref{sec:trazados-gtfs}» & \\
\midrule
\multicolumn{5}{@{}l}{\emph{Emparejamiento}}\\
\texttt{iDepartureCluster\allowbreak GapSeconds} & 90 s & Separación máxima para fusionar variantes & «\nameref{sec:alias}» & \\
\texttt{MATCH\_TOLERANCE\_SECONDS} & 180 s & Tolerancia del emparejamiento de respaldo & «\nameref{sec:emparejamiento}» & \\
\midrule
\multicolumn{5}{@{}l}{\emph{Localización}}\\
\texttt{METERS\_PER\_DEGREE} & 111\,320 m & Metros por grado de latitud & «\nameref{sec:proyeccion}» & \\
\texttt{MAX\_LOCATION\_AGE\_SECONDS} & 600 s & Antigüedad máxima del dato & «\nameref{sec:condiciones-gps}» & \\
\texttt{HOLD\_SECONDS} & 180 s & Tiempo que se guarda un autobús que falta en el feed & «\nameref{sec:memoria-feed}» & \checkmark\\
\texttt{NOT\_STARTED\_GRACE\_SECONDS} & 60 s & El viaje cuenta como empezado este tiempo antes de su salida & «\nameref{sec:condiciones-gps}» & \\
\texttt{MAX\_WAIT\_AT\_ORIGIN\_SECONDS} & 3\,600 s & Antelación máxima para situar al autobús en la cabecera & «\nameref{sec:cabecera}» & \\
\texttt{SEGMENTS\_BEHIND\_NEXT\_STOP} & 3 & Primer tramo de la ventana: $n-3$ & «\nameref{sec:ventana}» & \\
\texttt{SEGMENTS\_AHEAD\_OF\_NEXT\_STOP} & 2 & Último tramo de la ventana: $n+2$ (paradas hasta $p_{n+3}$) & «\nameref{sec:ventana}» & \\
\texttt{MAX\_DISTANCE\_FROM\_SHAPE\_METERS} & 50 m & Pasillo con trazado & «\nameref{sec:trazado}» & \checkmark\\
\texttt{BEYOND\_NEXT\_STOP\_SLACK\_METERS} & 100 m & Límite de los candidatos esperados & «\nameref{sec:trazado}» & \\
\texttt{AMBIGUITY\_MARGIN\_METERS} & 20 m & Ventaja exigida a un candidato «más allá» & «\nameref{sec:trazado}» & \\
\texttt{MAX\_DISTANCE\_FROM\_ROUTE\_METERS} & 150 m & Pasillo mínimo sin trazado; distancia máxima a la primera parada en la cabecera & «\nameref{sec:sin-trazado}» & \checkmark\\
\texttt{MAX\_DISTANCE\_FROM\_ROUTE\_SEGMENT\_RATIO} & 0,15 & Pasillo sin trazado, fracción del tramo & «\nameref{sec:sin-trazado}» & \checkmark\\
\midrule
\multicolumn{5}{@{}l}{\emph{Modelo}}\\
Franjas & 0, 6, 9, 13, 16 y 20 h & Inicio de cada periodo del día & Definición~\ref{def:franja} & \\
\texttt{LENGTH\_CUTS\_METERS} & 300, 600, 1\,000, 1\,500, 3\,000 m & Cortes de los cubos de longitud & «\nameref{sec:respaldo-largo}» & \\
Cotas de la proporción & 0,3 -- 1,5 & Acotación del respaldo por longitud & «\nameref{sec:respaldo-largo}» & \\
\texttt{MIN\_FACTOR}, \texttt{MAX\_FACTOR} & 0,6 -- 1,2 & Acotación de $k$ & «\nameref{sec:k}» & \checkmark\\
\midrule
\multicolumn{5}{@{}l}{\emph{Predicción}}\\
\texttt{REGISTRATION\_SEGMENT\_SHARE} & 0,5 & Fracción del tramo que falta en la capa 2 & «\nameref{sec:capa2}» & \checkmark\\
\texttt{REGISTRATION\_MAX\_SEGMENT\_SECONDS} & 1\,200 s & Tramo máximo para la capa 2 & «\nameref{sec:capa2}» & \checkmark\\
\texttt{REGISTRATION\_CROSSCHECK\_MAX\_SEGMENT\_SECONDS} & 360 s & Tramo máximo para el contraste & «\nameref{sec:desconfianza}» & \checkmark\\
\texttt{MAX\_GPS\_VS\_STOP\_DISAGREEMENT\_SECONDS} & 360 s & Discrepancia máxima GPS / parada & «\nameref{sec:desconfianza}» & \checkmark\\
\texttt{SCHEDULE\_ORIGIN\_OFFSET\_SECONDS} & 60 s & Salida de cabecera tras la hora publicada & «\nameref{sec:capa3}» & \checkmark\\
\texttt{WRONG\_SCHEDULE\_PASSED\_SECONDS} & 600 s & \texttt{likely\_departed}: hace cuánto & «\nameref{sec:capa3}» & \\
\texttt{WRONG\_SCHEDULE\_GAP\_SECONDS} & 600 s & \texttt{likely\_departed}: horario lento & «\nameref{sec:capa3}» & \\
\texttt{WRONG\_SCHEDULE\_LEARNED\_SHARE} & 0,8 & \texttt{likely\_departed}: parte aprendida & «\nameref{sec:capa3}» & \\
\texttt{TripOffsets::PARTS} & 3 & Partes del recorrido en la corrección por viaje & «\nameref{sec:correccion-viaje}» & \checkmark\\
Calendario de solo domingo & máscara 64 & Marca un festivo entre semana & «\nameref{sec:correccion-viaje}» & \\
\texttt{MAX\_PLAUSIBLE\_DELAY\_SECONDS} & 2\,700 s & Retraso máximo que se afirma & «\nameref{sec:decision}» & \checkmark\\
\texttt{STALE\_GRACE\_SECONDS} & 120 s & Llegada de la capa 2 en el pasado que aún se muestra & «\nameref{sec:decision}» & \\
\texttt{PASSED\_DISTANCE\_METERS} & 150 m & Distancia del GPS a la parada para darla por pasada sin que lo diga el feed & «\nameref{sec:paso}» & \\
\texttt{ARRIVING\_MAX\_SECONDS} & 180 s & Tiempo en «0~min» sin confirmar el paso antes de «puede que ya haya pasado» & «\nameref{sec:llegando}» & \\
\texttt{TTL\_SECONDS} (registro de pasos) & 36 h & Caducidad de los pasos, desde el final de su día & «\nameref{sec:registro-pasos}» & \\
\texttt{DAY\_SECONDS}/2 & 43\,200 s & Adelanto de la salida del feed sobre la hora actual por encima del cual la entrada es de un viaje de ayer & «\nameref{sec:mas-de-24h}» & \\
\texttt{LATE\_GRACE\_SECONDS} & 600 s & De madrugada, la ficha es la del viaje de ayer hasta este tiempo después de su última parada & «\nameref{sec:mas-de-24h}» & \\
Retirada de la lista & $-2$ min & Minutos por debajo de los cuales se retira una fila & «\nameref{sec:retraso-mostrado}» & \\
Umbral de «adelantado» & $-2$ min & Adelanto mínimo que se anuncia & «\nameref{sec:retraso-mostrado}» & \\
\texttt{POSITION\_SANITY\_SECONDS} & 900 s & Discrepancia máxima con la llegada plana, en redes con retraso declarado & «\nameref{sec:otras-redes}» & \\
\midrule
\multicolumn{5}{@{}l}{\emph{Aprendizaje}}\\
Captura & continua, bloques de 1 h, cada 20 s & Bloque e intervalo de captura & «\nameref{sec:captura}» & \\
\texttt{MAX\_GAP\_BETWEEN\_SNAPSHOTS\_SECONDS} & 300 s & Hueco que corta una ventana; separación máxima del par de instantáneas del paso real & «\nameref{sec:paso-real}» & \\
\texttt{OBSERVED\_AFTER\_PREDICTION\_SECONDS} & 600 s & Margen de observación tras la llegada & «\nameref{sec:muestras-ritmo}» & \\
\texttt{MIN\_REMAINING\_SECONDS\_FOR\_FACTOR} & 120 s & Horizonte mínimo de una muestra de ritmo & «\nameref{sec:muestras-ritmo}» & \\
\texttt{MAX\_SAMPLES\_PER\_TRACK} & 30 & Muestras de ritmo por pista & «\nameref{sec:muestras-ritmo}» & \\
\texttt{MIN\_TRACKS\_PER\_SEGMENT} & 2 & Pistas mínimas por tramo y franja & «\nameref{sec:contraccion-tramo}» & \checkmark\\
\texttt{SEGMENT\_BAND\_PRIOR\_TRACKS} & 5 & Peso de la mediana del tramo, $n_0$ & «\nameref{sec:contraccion-tramo}» & \checkmark\\
\texttt{MIN\_SAMPLES\_PER\_LENGTH\_RATIO} & 20 & Muestras mínimas por cubo y franja & «\nameref{sec:respaldo-largo}» & \\
\texttt{MIN\_TRACKS\_PER\_BAND} & 8 & Pistas mínimas para $k_b$ & «\nameref{sec:k}» & \\
\texttt{MIN\_TRACKS\_PER\_LINE} & 25 & Pistas mínimas para $k_{\ell,b}$ & «\nameref{sec:k}» & \checkmark\\
\texttt{MIN\_VEHICLES\_PER\_LINE} & 5 & Vehículos mínimos para $k_{\ell,b}$ & «\nameref{sec:k}» & \\
\texttt{SHRINKAGE\_TRACKS} & 10 & Peso de $k=1$ en la contracción & «\nameref{sec:k}» & \\
\texttt{SAMPLES\_RETENTION\_DAYS} & 120 d & Retención de muestras & «\nameref{sec:retencion}» & \\
\texttt{MAX\_SAMPLES\_PER\_SEGMENT} & 300 & Muestras por tramo y tipo de día & «\nameref{sec:retencion}» & \\
\texttt{MAX\_PACE\_SAMPLES} & 500\,000 & Muestras de ritmo en total & «\nameref{sec:retencion}» & \\
Pliegues & 2 & $\mathrm{crc32}(\texttt{track})\bmod2$ & «\nameref{sec:recalibracion}» & \\
\texttt{TRIP\_OFFSET\_RANKS} & 0,3 a 0,5 & Percentiles de los desvíos entre los que la recalibración elige cada noche: el que más acierta a $\pm2$~min y, con empate, el de menos error medio (0,4 en los vectores) & «\nameref{sec:correccion-percentil}» & \checkmark\\
Pasos fuera de la corrección por viaje & 2\,700 s & Se descartan los pasos que se apartan más de este tiempo de su hora publicada & «\nameref{sec:recalibracion}» & \\
\texttt{TRIP\_OFFSET\_CLIP\_SECONDS} & 900 s & Recorte de cada desvío de la corrección por viaje & «\nameref{sec:recalibracion}» & \checkmark\\
\texttt{TRIP\_OFFSET\_LINE\_PRIOR\_PASSES} & 20 & Peso de la franja en la corrección de la línea & «\nameref{sec:recalibracion}» & \\
\texttt{TRIP\_OFFSET\_TRIP\_PRIOR\_DAYS} & 1 & Peso de la línea en la corrección del viaje, en días & «\nameref{sec:recalibracion}» & \checkmark\\
\texttt{TRIP\_OFFSET\_TEST\_DAYS} & 7 d & Días de prueba de la corrección por viaje & «\nameref{sec:recalibracion}» & \\
\midrule
\multicolumn{5}{@{}l}{\emph{Evaluación}}\\
Paradas evaluadas & $n$ a $n+8$ & Paradas de cada autobús con las que se mide & «\nameref{sec:verdad}» & \\
Horizontes & 0--20 min & Tiempo hasta el paso real que se evalúa & «\nameref{sec:verdad}» & \\
Histogramas & 15 s & Anchura de las clases del error en la captura continua & «\nameref{sec:varios-dias}» & \\
Datos del día para los rivales & 60 s & Antigüedad mínima de los tramos del día que pueden usar & «\nameref{sec:comparacion}» & \\
\midrule
\multicolumn{5}{@{}l}{\emph{Trip Updates (Metro, Euskotren, tranvías, Renfe, Tuvisa, Dbus, Álavabus)}}\\
\texttt{match\_tolerance\_seconds} & 600 s & Tolerancia del emparejamiento de reserva, por paradas y hora & «\nameref{sec:otras-redes}» & \\
\texttt{VEHICLE\_DISAGREEMENT\_SECONDS} & 1\,200 s & Hora del feed que contradice al GPS & «\nameref{sec:otras-redes}» & \\
\texttt{MIN\_VEHICLES\_TO\_DISTRUST} & 3 & Vehículos con posición para desconfiar del feed & «\nameref{sec:otras-redes}» & \\
\texttt{MAX\_GPS\_DELAY\_SECONDS} & 2\,700 s & Retraso máximo creíble por GPS & «\nameref{sec:otras-redes}» & \\
\texttt{STALE\_PREDICTION\_SECONDS} & 60 s & Predicción vencida & «\nameref{sec:otras-redes}» & \\
\texttt{MAX\_STALE\_PREDICTION\_SECONDS} & 3\,600 s & Desfase máximo que se extrapola & «\nameref{sec:otras-redes}» & \\
\texttt{DUE\_SECONDS} & 60 s & Salida que ya debería estar en la estación & «\nameref{sec:otras-redes}» & \\
\texttt{LATE\_TRAINS\_LOOKBACK\_SECONDS} & 3\,600 s & Salidas pasadas que se miran por si van tarde & «\nameref{sec:otras-redes}» & \\
\texttt{ON\_TIME\_SECONDS} & 60 s & Retraso que se redondea a cero & «\nameref{sec:otras-redes}» & \\
\texttt{LIVE\_DEPARTED\_AFTER\_SECONDS} & 45 s & Retirada de una fila en directo & «\nameref{sec:otras-redes}» & \\
\texttt{SCHEDULED\_DEPARTED\_AFTER\_SECONDS} & 60 s & Retirada de una fila programada & «\nameref{sec:otras-redes}» & \\
\end{longtable}
}

\chapter{Pseudocódigo}
\label{ap:algoritmos}

Este apéndice reúne el pseudocódigo de las funciones del método, en el orden en que se usan (el de la
decisión de cada fila está en el capítulo «\nameref{cap:prediccion}»). Los índices de
paradas empiezan en 1, como en el texto: el viaje tiene paradas $p_1,\dots,p_N$ y el tramo $i$ va de $p_i$ a
$p_{i+1}$. $\redondeo$ es el redondeo con las mitades alejándose del cero, $\mediana$ la mediana baja y $\clamp$ la
acotación (sección «\nameref{sec:convenciones}»). «Nulo» significa que el valor no existe.

\begin{figure}[ht]
\caption{Proyección de un Punto Sobre un Segmento}\label{alg:proyeccion}
\begin{algorithmic}[1]
\Function{Proyectar}{$\varphi,\lambda,A,B$}
  \State $m_\lambda\gets\cos(\varphi\cdot\pi/180)\cdot111\,320$;\quad $m_\varphi\gets111\,320$
  \State $x_A\gets(\lambda_A-\lambda)\,m_\lambda$;\ $y_A\gets(\varphi_A-\varphi)\,m_\varphi$;\
         $x_B\gets(\lambda_B-\lambda)\,m_\lambda$;\ $y_B\gets(\varphi_B-\varphi)\,m_\varphi$
  \State $\Delta x\gets x_B-x_A$;\ $\Delta y\gets y_B-y_A$;\ $\ell^2\gets\Delta x^2+\Delta y^2$
  \State $f\gets0$
  \If{$\ell^2>0$} $f\gets\clamp\big(-(x_A\Delta x+y_A\Delta y)/\ell^2,\ 0,\ 1\big)$ \EndIf
  \State \Return $\big(f,\ \sqrt{(x_A+f\Delta x)^2+(y_A+f\Delta y)^2},\ \sqrt{\ell^2}\big)$ \Comment{fracción, distancia, longitud}
\EndFunction
\end{algorithmic}
\nota{En un plano local centrado en el punto.}
\end{figure}

\begin{figure}[ht]
\caption{Posición del Autobús en el Horario (Capa 1, sin el Contraste)}\label{alg:posicion}
\begin{algorithmic}[1]
\Function{Posición}{$J$, entrada $e$, $\tnow$}
  \If{$e$ no tiene $\varphi$, $\lambda$, $\tgps$ o $n$} \Return nulo \EndIf
  \If{$\tnow-\tgps>600$ \textbf{o} $N<2$} \Return nulo \EndIf
  \If{$\tgps<d_1-60$} \Return la posición en la cabecera~\eqref{eq:cabecera} si $d_1-\tgps\le3\,600$ y el GPS está a 150~m o menos de $p_1$; si no, nulo \EndIf
  \State $\nu\gets$ el índice con orden de parada igual a $n$; \textbf{si} no existe, \Return nulo
  \State $i_0\gets\max(1,\nu-3)$;\quad $i_1\gets\min(N-1,\nu+2)$
  \If{$J$ tiene trazado $\Sigma$ \textbf{y} $s_i$ existe para $i=i_0,\dots,i_1+1$}
    \State $\sigma_0\gets s_{i_0}$;\ $\sigma_1\gets s_{i_1+1}$;\ $\sigma_{\lim}\gets s_\nu+100$;\ $E\gets$ nulo;\ $M\gets$ nulo
    \For{$j=1,\dots,M_\Sigma-1$} \Comment{segmentos del trazado, en orden}
      \If{$\sigma_{j+1}<\sigma_0$ \textbf{o} $\sigma_j>\sigma_1$} \textbf{continuar} \EndIf
      \State $(f,d,\cdot)\gets$ \Call{Proyectar}{$\varphi,\lambda,\Sigma_j,\Sigma_{j+1}$};\quad $s^\ast\gets\sigma_j+f(\sigma_{j+1}-\sigma_j)$
      \If{$s^\ast\le\sigma_{\lim}$}
        \If{$E$ es nulo \textbf{o} $d<E.d$} $E\gets(d,s^\ast)$ \EndIf
      \ElsIf{$M$ es nulo \textbf{o} $d<M.d$} $M\gets(d,s^\ast)$
      \EndIf
    \EndFor
    \State $B\gets E$
    \If{$M$ no es nulo \textbf{y} ($B$ es nulo \textbf{o} $M.d+20<B.d$)} $B\gets M$ \EndIf
    \If{$B$ es nulo \textbf{o} $B.d>50$} \Return nulo \EndIf
    \State $s^\ast\gets\clamp(B.s^\ast,\sigma_0,\sigma_1)$
    \For{$i=i_0,\dots,i_1$}
      \If{$s^\ast>s_{i+1}$ \textbf{y} $i<i_1$} \textbf{continuar} \EndIf
      \State $f\gets0$; \textbf{si} $s_{i+1}>s_i$ \textbf{entonces} $f\gets\clamp\big((s^\ast-s_i)/(s_{i+1}-s_i),0,1\big)$
      \State \Return $\big(\Tpos=\redondeo(d_i+f(a_{i+1}-d_i)),\ \tgps,\ B.d\big)$
    \EndFor
  \Else \Comment{sin trazado: rectas entre paradas}
    \State $B\gets$ nulo
    \For{$i=i_0,\dots,i_1$}
      \State $(f,d,L)\gets$ \Call{Proyectar}{$\varphi,\lambda,p_i,p_{i+1}$}
      \If{$B$ es nulo \textbf{o} $d<B.d$} $B\gets\big(d,\ \max(150,0{,}15L),\ d_i+f(a_{i+1}-d_i)\big)$ \EndIf
    \EndFor
    \If{$B.d>B.c$} \Return nulo \EndIf
    \State \Return $\big(\Tpos=\redondeo(B.t),\ \tgps,\ B.d\big)$
  \EndIf
\EndFunction
\end{algorithmic}
\end{figure}

\begin{figure}[ht]
\caption{Tiempo de un Tramo y Tiempo Esperado Entre Dos Horas del Horario}\label{alg:esperado}
\begin{algorithmic}[1]
\Function{TiempoTramo}{$i,b,k$}
  \State $h\gets a_{i+1}-a_i$
  \State $\tau\gets$ valor de $\kappa_i=u_i{>}u_{i+1}$ en la franja $b$; si no existe, el valor \texttt{all}; si no, nulo
  \If{$\tau$ no es nulo} \Return $\tau$ \EndIf
  \If{$s_i$ y $s_{i+1}$ existen \textbf{y} $L=s_{i+1}-s_i>0$ \textbf{y} hay proporciones}
    \State $r\gets$ proporción del cubo $c(L)$ en la franja $b$; si no existe, la \texttt{all}; si no, nulo
    \If{$r$ no es nulo} \Return $\clamp(r,0{,}3,1{,}5)\cdot h$ \EndIf
  \EndIf
  \State \Return $k\cdot h$
\EndFunction
\Statex
\Function{Esperado}{$J,t_0,t_1,b,k$}
  \If{la tabla de tramos está vacía} \Return $k\,(t_1-t_0)$ \EndIf
  \State $S\gets0$
  \For{$i=1,\dots,N-1$}
    \If{$a_{i+1}\le t_0$} \textbf{continuar} \EndIf
    \If{$a_i\ge t_1$} \textbf{terminar} \EndIf
    \State $h\gets a_{i+1}-a_i$;\quad $w\gets1$
    \If{$h>0$} $w\gets\big(\min(a_{i+1},t_1)-\max(a_i,t_0)\big)/h$ \EndIf
    \State $S\gets S+\max(0,w)\cdot$\Call{TiempoTramo}{$i,b,k$}
  \EndFor
  \State \Return $S$
\EndFunction
\end{algorithmic}
\end{figure}

\begin{figure}[ht]
\caption{Registro (Capa 2), Llegada Desde el Registro y Hora de Paso}\label{alg:registro}
\begin{algorithmic}[1]
\Function{Registro}{$J,e,\tnow$}
  \If{falta $n$ \textbf{o} el par $(J,v)$ no es fiable \textbf{o} $n<2$} \Return nulo \EndIf
  \State $\tref\gets\tgps$ si existe; si no, $t_{\mathrm{rec}}$
  \If{$\tref$ no existe \textbf{o} $\tnow-\tref>600$ \textbf{o} $\tref<d_1-60$} \Return nulo \EndIf
  \If{no hay paradas de orden $n-1$ y $n$ \textbf{o} $a_n-a_{n-1}>1\,200$} \Return nulo \EndIf
  \State $b\gets\franja(D,\tref)$;\ $k\gets k(\ell,D,\tref)$
  \State $\theta\gets\redondeo\big(0{,}5\cdot$\Call{Esperado}{$J,a_{n-1},a_n,b,k$}$\big)$
  \State \Return $\big(n,\ \tref,\ a_{n-1},\ a_n,\ \theta,\ \delta_{\mathrm{reg}}=\tref+\theta-a_n,\ g=a_n-a_{n-1}\big)$
\EndFunction
\Statex
\Function{LlegadaRegistro}{$J,e,q,a_\dest$}
  \State $G\gets$\Call{Registro}{$J,e,\tnow$}; \textbf{si} $G$ es nulo \textbf{o} $q<G.n$, \Return nulo
  \State $X\gets0$
  \If{$a_\dest>G.a_n$}
    \State $b\gets\franja(D,G.\tref)$;\ $k\gets k(\ell,D,G.\tref)$
    \State $X\gets$\Call{Esperado}{$J,G.a_n,a_\dest,b,k$}
  \EndIf
  \State \Return $\redondeo(G.\tref+G.\theta+X)$
\EndFunction
\Statex
\Function{HoraDePaso}{$J,e,q,a_\dest$}
  \State $P\gets$ posición con contraste
  \If{$P$ no es nulo}
    \If{$P$ es de cabecera} \Return nulo \EndIf
    \If{$a_\dest<P.\Tpos$ \textbf{y} ($q<n$ \textbf{o} $\mathrm{dist}(e,p_q)>150$)}
      \State \Return $P.\tgps-(P.\Tpos-a_\dest)$
    \EndIf
    \If{$q<n$} \Return $P.\tgps$ \EndIf
    \State \Return nulo
  \EndIf
  \If{$q<n$ \textbf{y} $G\gets$\Call{Registro}{$J,e,\tnow$} no es nulo} \Return $G.\tref-\max(0,G.a_{n-1}-a_\dest)$ \EndIf
  \State \Return nulo
\EndFunction
\end{algorithmic}
\end{figure}

\begin{figure}[ht]
\caption{Posición con Contraste (Se Evalúa Antes que Cualquier Uso del Registro)}\label{alg:contraste}
\begin{algorithmic}[1]
\Function{PosiciónConContraste}{$J,e,\tnow$}
  \State $P\gets$\Call{Posición}{$J,e,\tnow$}
  \If{$P$ no es nulo}
    \State $G\gets$\Call{Registro}{$J,e,\tnow$}
    \If{$G$ no es nulo \textbf{y} $G.g\le360$ \textbf{y} $|(P.\tgps-P.\Tpos)-G.\delta_{\mathrm{reg}}|>360$}
      \State marcar el par $(J,v)$ como no fiable durante esta petición;\quad $P\gets$ nulo
    \EndIf
  \EndIf
  \State \Return $P$
\EndFunction
\end{algorithmic}
\end{figure}

\begin{figure}[ht]
\caption{Capa 3: Horario Corregido, Corrección por Viaje y Parte Aprendida}\label{alg:capa3}
\begin{algorithmic}[1]
\Function{SinDatos}{$J,q,a_\dest,\tnow$}
  \If{$J$ no tiene entrada en el feed \textbf{y} el registro de pasos de $J$ tiene $(n,t)$ con $q<n$} \Return $(\texttt{departed},\min(t,\tnow),0)$ \EndIf
  \State $R\gets(\texttt{scheduled},a_\dest,0)$
  \If{la tabla de tramos está vacía \textbf{o} $J$ no tiene paradas \textbf{o} $a_\dest\le d_1$} \Return $R$ \EndIf
  \State $D'\gets D$, o el domingo de esa semana si $D$ es festivo entre semana
  \State $t_{\mathrm{mid}}\gets\lfloor(d_1+a_\dest)/2\rfloor$;\ $b\gets\franja(D',t_{\mathrm{mid}})$;\ $k\gets k(\ell,D',t_{\mathrm{mid}})$
  \State $C\gets d_1+60+\redondeo\big($\Call{Esperado}{$J,d_1-1,a_\dest,b,k$}$\big)+\redondeo\big($\Call{Corrección}{$J,q,D'$}$\big)$
  \If{$C\ge\tnow$} \Return $(\texttt{estimated},C,0)$ \EndIf
  \If{$C<\tnow-600$ \textbf{y} $a_\dest-C\ge600$ \textbf{y} \Call{ParteAprendida}{$J,d_1-1,a_\dest,b$}$\ \ge0{,}8$}
    \State \Return $(\texttt{likely\_departed},a_\dest,0)$
  \EndIf
  \State \Return $R$
\EndFunction
\Statex
\Function{Corrección}{$J,q,D'$}
  \If{la tabla de correcciones está vacía \textbf{o} ninguna parada de $J$ tiene orden $q$} \Return $0$ \EndIf
  \State $i\gets$ posición de la parada de orden $q$ (0 la cabecera);\ $p\gets\min(2,\lfloor3i/N\rfloor)$
  \State $v\gets(\ell,p_1,p_N,d_1,N)$;\ $\theta\gets$ tipo de día de $D'$;\ $b_1\gets\franja(D',d_1)$
  \If{la tabla tiene $\Delta(v,\theta,\cdot)$} \Return $\Delta(v,\theta,p)$ \EndIf
  \If{la tabla tiene $l(\ell,b_1)$} \Return $l(\ell,b_1)$ \EndIf
  \If{la tabla tiene $g(b_1)$} \Return $g(b_1)$ \EndIf
  \State \Return $0$
\EndFunction
\Statex
\Function{ParteAprendida}{$J,t_0,t_1,b$}
  \State $T\gets0$;\ $A\gets0$
  \For{$i=1,\dots,N-1$}
    \If{$a_{i+1}\le t_0$} \textbf{continuar} \EndIf
    \If{$a_i\ge t_1$} \textbf{terminar} \EndIf
    \State $T\gets T+(a_{i+1}-a_i)$
    \If{$\kappa_i$ tiene valor en $b$ o \texttt{all}} $A\gets A+(a_{i+1}-a_i)$ \EndIf
  \EndFor
  \State \Return $0$ si $T\le0$; si no, $A/T$
\EndFunction
\end{algorithmic}
\end{figure}

\begin{figure}[ht]
\caption{Paso Real por una Parada y Aclarado de Muestras}\label{alg:paso-real}
\begin{algorithmic}[1]
\Function{PasoReal}{pista $((t_k,\theta_k))_{k=1}^{K}$, $a$}
  \For{$k=2,\dots,K$}
    \If{$\theta_k\le\theta_{k-1}$ \textbf{o} $t_k-t_{k-1}>300$} \textbf{continuar} \EndIf
    \If{$\theta_{k-1}\le a\le\theta_k$}
      \State \Return $t_{k-1}+(a-\theta_{k-1})/(\theta_k-\theta_{k-1})\cdot(t_k-t_{k-1})$
    \EndIf
  \EndFor
  \State \Return nulo
\EndFunction
\Statex
\Function{Aclarar}{lista $x_0,\dots,x_{c-1}$}
  \If{$c\le30$} \Return la lista \EndIf
  \State \Return $\big(x_{\lfloor i\,c/30\rfloor}\big)_{i=0}^{29}$ \Comment{división entera}
\EndFunction
\end{algorithmic}
\end{figure}

\begin{figure}[ht]
\caption{Ajuste de los Tiempos por Tramo y de las Proporciones por Longitud}\label{alg:ajuste-tramos}
\begin{algorithmic}[1]
\Function{AjustarTramos}{muestras de tramo, en orden}
  \For{cada muestra $(\kappa,j,b,y)$}
    \State $V[\kappa][\texttt{all}][j]\gets y$;\quad $V[\kappa][b][j]\gets y$ \Comment{la última de cada pista gana}
  \EndFor
  \State $T\gets$ diccionario vacío
  \For{cada tramo $\kappa$ de $V$, con sus grupos en orden de aparición (\texttt{all} primero)}
    \If{$|V[\kappa][\texttt{all}]|<2$} \textbf{continuar} \EndIf
    \State $t\gets\{\texttt{n}:|V[\kappa][\texttt{all}]|\}$;\quad $m\gets\mediana(V[\kappa][\texttt{all}])$
    \State $\mathit{varias}\gets$ (hay al menos dos franjas en $V[\kappa]$)
    \For{cada grupo $g$ de $V[\kappa]$}
      \If{$g=\texttt{all}$ \textbf{y no} $\mathit{varias}$} \textbf{continuar} \EndIf
      \State $n_g\gets|V[\kappa][g]|$; \textbf{si} $n_g<2$, \textbf{continuar}
      \State $x\gets\mediana(V[\kappa][g])$
      \If{$g\ne\texttt{all}$} $x\gets(n_g\,x+5\,m)/(n_g+5)$ \EndIf
      \State $t[g]\gets\redondeo(x)$
    \EndFor
    \If{$t$ tiene algún grupo además de \texttt{n}} $T[\kappa]\gets t$ \EndIf
  \EndFor
  \State \Return $T$ ordenado por clave
\EndFunction
\Statex
\Function{AjustarProporciones}{muestras de tramo, longitudes $L[\kappa]$}
  \For{cada muestra con $L[\kappa]>0$ y $h>0$}
    \State $c\gets c(L[\kappa])$;\ añadir $y/h$ a $Q[c][\texttt{all}]$ y a $Q[c][b]$
  \EndFor
  \State \Return $\{c:\{g:\redondeo_3(\mediana Q[c][g])\ \text{para cada } g \text{ con } |Q[c][g]|\ge20\}\}$
\EndFunction
\end{algorithmic}
\end{figure}

\begin{figure}[ht]
\caption{Ajuste del Factor de Ritmo}\label{alg:ajuste-k}
\begin{algorithmic}[1]
\Function{FactorContraído}{pistas $\mathcal J$, cada una con sus muestras $(H,Y)$}
  \For{cada pista $j$} $\rho_j\gets\sum Y/\sum H$ \EndFor
  \State $m\gets|\mathcal J|$;\quad $\tilde\rho\gets\mediana_j\rho_j$
  \State \Return $\redondeo_3\big(\clamp(1+(\tilde\rho-1)\,m/(m+10),\ 0{,}6,\ 1{,}2)\big)$
\EndFunction
\Statex
\Function{AjustarFactores}{muestras de ritmo}
  \For{cada franja $b$ con $\ge8$ pistas} $k_b\gets$\Call{FactorContraído}{pistas de $b$} \EndFor
  \For{cada línea $\ell\ne$ vacía y franja $b$ con $\ge25$ pistas y $\ge5$ vehículos}
    \State $k_{\ell,b}\gets$\Call{FactorContraído}{pistas de $\ell$ en $b$}
  \EndFor
  \State \Return $\{\texttt{global}:k_b,\ \texttt{lines}:k_{\ell,b}\}$
\EndFunction
\end{algorithmic}
\end{figure}

\begin{figure}[ht]
\caption{Validación Cruzada y Escritura de los Tiempos por Tramo}\label{alg:validacion}
\begin{algorithmic}[1]
\Function{RecalibrarTramos}{muestras $X$, factores vigentes}
  \State $S_a\gets0$;\ $S_d\gets0$;\ $c\gets0$
  \For{$f\in\{0,1\}$}
    \State $T_f\gets$\Call{AjustarTramos}{$\{x\in X:\mathrm{crc32}(x.\texttt{track})\bmod2\ne f\}$}
    \For{cada $x\in X$ con $\mathrm{crc32}(x.\texttt{track})\bmod2=f$ y $\tau\gets T_f(x.\kappa,x.b)$ definido}
      \State $k\gets k(x.\ell,x.D,x.\texttt{at})$;\ $S_a\gets S_a+|k\,x.h-x.y|/60$;\ $S_d\gets S_d+|\tau-x.y|/60$;\ $c\gets c+1$
    \EndFor
  \EndFor
  \State validar igual las proporciones; incluirlas solo si su suma de errores baja
  \If{$c=0$ \textbf{o} $S_d\ge S_a$} \Return sin escribir \EndIf
  \State $T\gets$\Call{AjustarTramos}{$X$}
  \If{$T$ y las proporciones son iguales a las del fichero actual} \Return sin escribir \EndIf
  \State escribir $\{\texttt{generated},\ \texttt{samples}:|X|,\ \texttt{segments}:T,\ \texttt{lengthRatios}\}$
\EndFunction
\end{algorithmic}
\end{figure}

\chapter{Formatos de Datos}
\label{ap:formatos}

\section{Base de Datos del Horario}

Es una base SQLite por red, de solo lectura, que se reconstruye cada día desde el GTFS. Las tablas y columnas que usa el
motor son estas (los tipos son los de SQLite):

\begin{lstlisting}[language=SQL,numbers=none]
CREATE TABLE stops (id TEXT PRIMARY KEY, name TEXT NOT NULL, lat REAL NOT NULL, lon REAL NOT NULL, ...);
CREATE TABLE lines (id TEXT PRIMARY KEY, code TEXT NOT NULL, name TEXT NOT NULL, ...);
CREATE TABLE journey_patterns (id TEXT PRIMARY KEY, line_id TEXT NOT NULL, headsign TEXT);
CREATE TABLE journey_pattern_stops (journey_pattern_id TEXT NOT NULL, seq_order INTEGER NOT NULL,
                                    stop_id TEXT NOT NULL);
CREATE TABLE service_journeys (id TEXT PRIMARY KEY, line_id TEXT NOT NULL,
                               journey_pattern_id TEXT NOT NULL, trip_number TEXT,
                               calendar_id TEXT NOT NULL, first_departure_seconds INTEGER,
                               shape_id TEXT);
CREATE TABLE passing_times (service_journey_id TEXT NOT NULL, seq_order INTEGER NOT NULL,
                            stop_id TEXT NOT NULL, arrival_seconds INTEGER,
                            departure_seconds INTEGER, dist_m REAL);
CREATE TABLE shapes (id TEXT PRIMARY KEY, points TEXT NOT NULL);   -- JSON [[lat, lon, dist], ...]
CREATE TABLE trip_aliases (trip_id TEXT PRIMARY KEY, service_journey_id TEXT NOT NULL);
\end{lstlisting}

\noindent Para Bizkaibus, en octubre de 2026: 100 líneas, 2\,346 paradas, 483 recorridos, 6\,097 viajes fusionados
(23\,332 \texttt{trip\_id} originales en \texttt{trip\_aliases}), 138\,513 horas de paso y 440 trazados. Una fila de
\texttt{passing\_times}:

\begin{lstlisting}[numbers=none]
('trp_A0651_737_OP43FIN_30900_O43FI0651_065151_3', 2, '1529', 32430, 32432, 13223.0)
\end{lstlisting}

\noindent La consulta que usa el motor para un viaje devuelve sus paradas en orden de \texttt{seq\_order}, con
llegada, salida, \texttt{dist\_m} y las coordenadas de la parada; la de un trazado devuelve la lista de puntos.

\section{Tiempos por Tramo: \texttt{segment-times.json}}

\begin{lstlisting}[numbers=none]
{
  "generated": "2026-10-31 11:31",
  "samples": 195786,
  "segments": {
    "1003>706": { "n": 43, "all": 118, "laborable.mediodia": 118, "laborable.punta_tarde": 126, ... },
    ...
  },
  "lengthRatios": [
    { "all": 0.987, "laborable.punta_manana": 1.016, ... },
    { "all": 0.907, ... },
    ...
  ]
}
\end{lstlisting}

\begin{itemize}
  \item \texttt{segments}: un objeto por tramo, con clave «\texttt{desde>hasta}» (identificadores de parada) y valores
        enteros en segundos por franja y \texttt{all}, más \texttt{n}, el número de pistas (informativo).
  \item \texttt{lengthRatios}: proporción por cubo de longitud ($0$ a $5$) y franja. \textbf{Atención}: el
        serializador JSON de PHP escribe un objeto con claves enteras consecutivas desde 0 como una \emph{lista}. Si
        están todos los cubos, \texttt{lengthRatios} es una lista de seis objetos (el índice es el cubo); si falta
        alguno, es un objeto con claves \texttt{"0"}, \texttt{"1"}, \dots{} Un lector debe aceptar las dos formas.
  \item Una proporción entera, como $1$, puede aparecer escrita sin decimales.
\end{itemize}

\section{Factores de Ritmo: \texttt{pace-factors.json}}

\begin{lstlisting}[numbers=none]
{
  "generated": "2026-10-31 11:10",
  "captures": 0,
  "samples": 356001,
  "global": { "laborable.manana": 0.885, "laborable.punta_tarde": 0.928, ... },
  "lines":  { "3411": { "laborable.manana": 1.03, ... }, ... },
  "meta":   { "global": { "laborable.manana": 2260, ... }, "lines": { ... } }
}
\end{lstlisting}

\noindent \texttt{global} y \texttt{lines} contienen los factores ya acotados y redondeados a tres decimales;
\texttt{meta} guarda el número de pistas con que se calculó cada uno. Las claves de línea son el código numérico de la
línea como texto. \texttt{captures} vale 0 cuando se calibra desde el fichero acumulado de muestras, que es lo
habitual, y el número de capturas cuando se calibra directamente desde una carpeta de capturas.

\section{Corrección por Viaje: \texttt{trip-offsets.json}}

\begin{lstlisting}[numbers=none]
{
  "generated": "2026-11-07 11:53",
  "samples": 531799,
  "rank": 0.35,
  "trips": {
    "2151|1|3382|26102|18": { "laborable": [-39, -87, -32] },
    ...
  },
  "lines":  { "651": { "laborable.manana": -6, "fin_semana.punta_manana": -36, ... }, ... },
  "global": { "laborable.manana": 8, "fin_semana.noche": -44, ... }
}
\end{lstlisting}

\begin{itemize}
  \item \texttt{trips}: un objeto por viaje con la clave «línea\texttt{|}primera parada\texttt{|}última
        parada\texttt{|}salida\texttt{|}número de paradas», con la salida en segundos desde medianoche. Por tipo de día
        (\texttt{laborable}, \texttt{fin\_semana}) guarda tres enteros en segundos: la corrección del primer, del
        segundo y del tercer tercio del recorrido.
  \item \texttt{lines} y \texttt{global}: la corrección de cada línea en cada franja y la de cada franja, en segundos
        enteros. La franja es la de la salida de cabecera, y las claves de línea son su código como texto.
  \item \texttt{samples} es el número de pasos con que se calculó y \texttt{rank}, el percentil de los desvíos que
        eligió la recalibración esa noche (sección «\nameref{sec:correccion-percentil}»). Un festivo entre semana se
        consulta como fin de semana.
\end{itemize}

\section{Muestras Acumuladas (JSON Lines)}

Cada línea es un objeto JSON. Una muestra de tramo:

\begin{lstlisting}[numbers=none]
{"track":"2026-10-06|1112|trp_A3411_805_OP9FIN_45000_O9FI3411_341101_7","line":"3411",
 "date":"2026-10-06","at":47066,"band":"laborable.mediodia","from":"1479","to":"1480",
 "scheduled":101,"actual":115.2}
\end{lstlisting}

\noindent Una muestra de ritmo:

\begin{lstlisting}[numbers=none]
{"track":"2026-10-02|1101|trp_A3411_804_OP9FIN_43200_O9FI3411_341104_6","vehicle":"1101",
 "line":"3411","date":"2026-10-02","at":45895,"band":"laborable.manana",
 "scheduledRemaining":217,"actualRemaining":409.8,"actual":46304.8}
\end{lstlisting}

\noindent Al leer, se ignoran las líneas que no son JSON válido o a las que les falta algún campo obligatorio:
\texttt{track}, \texttt{at}, \texttt{band}, \texttt{from}, \texttt{to}, \texttt{scheduled} y \texttt{actual} en las de
tramo; \texttt{track}, \texttt{band}, \texttt{scheduledRemaining} y \texttt{actualRemaining} en las de ritmo. La
recalibración usa además \texttt{date} (retención y tope por tipo de día), \texttt{line} (factores por línea),
\texttt{vehicle} (mínimo de vehículos por línea) y, en las de ritmo, \texttt{at} y \texttt{actual} (validación); las
muestras que escribe la captura los traen siempre.

\section{Registro de Pasos}

Es un almacén clave-valor con dos tablas hash por red y día:

\begin{center}
\small
\begin{tabularx}{\textwidth}{@{}llX@{}}
\toprule
Clave & Campo & Valor\\
\midrule
\texttt{bide:<red>:pass:<fecha>} & referencia $\rho$ & \texttt{"<orden>:<segundos>"}; caduca 36~h después de ese día\\
\texttt{bide:<red>:stops:<fecha>} & \texttt{"<}$\rho$\texttt{>\#<orden>"} & \texttt{"<segundos>"}; caduca 3 días después de ese día\\
\bottomrule
\end{tabularx}
\end{center}

\noindent La primera guarda el último orden visto de cada referencia de viaje y es la que usa la predicción
(sección «\nameref{sec:registro-pasos}»); la segunda guarda la hora en que se vio cada orden y sirve para el historial de una
línea. Se escriben cada minuto. La \texttt{<fecha>} es la del día del viaje: pasada la medianoche, un autobús de un
viaje de ayer se apunta en las tablas de ayer, con los segundos en las horas de su horario (más de $86\,400$; sección
«\nameref{sec:mas-de-24h}»).

\section{El Feed SIRI-VM}

Este es un fragmento mínimo de una actividad con los elementos que se leen (Tabla~\ref{tab:feed}):

\begin{lstlisting}[language=XML,numbers=none]
<VehicleActivity>
  <RecordedAtTime>2026-10-28T15:12:30+01:00</RecordedAtTime>
  <MonitoredVehicleJourney>
    <VehicleJourneyRef>trp_A3414_808_OP9LSEPT_54000_...</VehicleJourneyRef>
    <VehicleRef>1112</VehicleRef>
    <VehicleLocation><Longitude>-2.9214</Longitude><Latitude>43.2601</Latitude></VehicleLocation>
    <LocationRecordedAtTime>2026-10-28T15:12:30+01:00</LocationRecordedAtTime>
    <MonitoredCall><StopPointRef>1480</StopPointRef><Order>3</Order></MonitoredCall>
  </MonitoredVehicleJourney>
</VehicleActivity>
\end{lstlisting}

\noindent Los retrasos (\texttt{Delay}) siguen el patrón \texttt{\^{}(-?)PT(?:(\textbackslash d+)H)?(?:(\textbackslash
d+)M)?(?:(\textbackslash d+)S)?\$}; cualquier otro valor se lee como cero.

\chapter{Vectores de Prueba}
\label{ap:vectores}

Los vectores de prueba son entradas concretas con la salida \emph{exacta} del código de producción. Los genera
\texttt{vectores.php}, un programa que va con el código privado: carga el motor tal como está (las clases del
emparejador, los factores y los tiempos por tramo, y las funciones de calibración del script de aprendizaje, sin ejecutar su programa principal) y
escribe \texttt{referencia/vectores.json}. Una implementación en cualquier lenguaje puede darse por correcta cuando
reproduce todos los vectores; \texttt{referencia/comprobar.py} lo hace para la implementación de referencia del
Apéndice~\ref{ap:referencia}.

Los valores reales se dan con todos sus decimales para comprobar también los cálculos en coma flotante; los
resultados que el método redondea deben coincidir exactamente.

\section{Funciones Básicas}

\begin{table}[ht]
\centering
\caption{Mediana Baja y Percentiles}\label{tab:v-percentil}
\begin{tabular}{@{}llr@{}}
\toprule
Valores & $p$ & Resultado\\
\midrule
5, 1, 4, 2, 3 & 0,5 & 3\\
1, 2, 3, 4 & 0,5 & 2 \quad(no 2,5: mediana baja)\\
10, 20, \dots, 100 & 0,9 & 90\\
\bottomrule
\end{tabular}
\end{table}

\begin{table}[ht]
\centering
\caption{Franja de una Fecha y un Instante}\label{tab:v-franja}
\begin{tabular}{@{}llrl@{}}
\toprule
Fecha & Día & Segundos & Franja\\
\midrule
2026-10-28 & miércoles & 28\,950 (8:02:30) & \texttt{laborable.punta\_manana}\\
2026-10-28 & miércoles & 21\,599 (5:59:59) & \texttt{laborable.madrugada}\\
2026-10-28 & miércoles & 21\,600 (6:00:00) & \texttt{laborable.punta\_manana}\\
2026-10-31 & sábado & 75\,600 (21:00:00) & \texttt{fin\_semana.noche}\\
2026-10-28 & miércoles & 90\,000 (25:00:00) & \texttt{laborable.madrugada}\\
2026-11-01 & domingo & 46\,800 (13:00:00) & \texttt{fin\_semana.mediodia}\\
\bottomrule
\end{tabular}
\end{table}

\begin{table}[ht]
\centering
\caption{Factor de Ritmo}\label{tab:v-k}
\begin{tabular}{@{}lrrl@{}}
\toprule
Línea & Segundos & $k$ & Por qué\\
\midrule
3414 & 28\,950 & 1,2 & Factor de línea 1,35, acotado\\
9999 & 28\,950 & 0,9 & Sin factor de línea: el global de la franja\\
9999 & 79\,200 & 0,6 & Global 0,4 de la noche, acotado\\
9999 & 50\,400 & 1,0 & Sin factor para el mediodía\\
03414 & 28\,950 & 1,2 & Se lee como la 3414\\
\bottomrule
\end{tabular}
\nota{Con $k_{\texttt{laborable.punta\_manana}}=0{,}9$, $k_{\texttt{laborable.noche}}=0{,}4$ y $k_{3414,\texttt{laborable.punta\_manana}}=1{,}35$, el miércoles 2026-10-28.}
\end{table}

\begin{table}[ht]
\centering
\caption{Proyección Sobre un Segmento de Prueba}\label{tab:v-proyeccion}
\begin{tabular}{@{}>{\raggedright\arraybackslash}p{0.22\textwidth}rrr@{}}
\toprule
Punto $(\varphi,\lambda)$ & $f$ & $d_\perp$ (m) & $\|B-A\|$ (m)\\
\midrule
$(43{,}2601;\,-2{,}925)$ & 0,5000000000000222 & 11,132000000369544 & 810,6880955946056\\
$(43{,}2605;\,-2{,}935)$ & 0 & 409,14505231598264 & 810,6827696075283\\
Segmento degenerado $A=B$, punto en $A$ & 0 & 0 & 0\\
\bottomrule
\end{tabular}
\nota{Segmento de $A=(43{,}26;\,-2{,}93)$ a $B=(43{,}26;\,-2{,}92)$.}
\end{table}

\noindent La longitud del segmento cambia ligeramente entre las dos primeras filas porque el factor
$\cos\varphi$ se toma en la latitud del autobús, no en la del segmento.

\begin{table}[ht]
\centering
\caption{Cubos de Longitud}\label{tab:v-cubo}
\begin{tabular}{@{}rr@{\qquad}rr@{\qquad}rr@{}}
\toprule
Metros & Cubo & Metros & Cubo & Metros & Cubo\\
\midrule
0 & 0 & 999,9 & 2 & 2\,999,9 & 4\\
299,9 & 0 & 1\,000 & 3 & 3\,000 & 5\\
300 & 1 & 1\,499,9 & 3 & 12\,000 & 5\\
& & 1\,500 & 4 & &\\
\bottomrule
\end{tabular}
\end{table}

\begin{table}[ht]
\centering
\caption{Respaldo por Longitud}\label{tab:v-respaldo}
\begin{tabular}{@{}rrlr>{\raggedright\arraybackslash}p{0.19\textwidth}@{}}
\toprule
$h$ (s) & $L$ (m) & Franja & Resultado (s) & Por qué\\
\midrule
120 & 800 & \texttt{laborable.punta\_manana} & 97,32000000000001 & $0{,}811\cdot120$\\
120 & 800 & \texttt{fin\_semana.noche} & 92,28 & $0{,}769\cdot120$ (\texttt{all})\\
300 & 2\,000 & \texttt{laborable.manana} & 450 & $2{,}0$ acotado a $1{,}5$\\
300 & 5\,000 & \texttt{laborable.manana} & 90 & $0{,}1$ acotado a $0{,}3$\\
300 & 100 & \texttt{laborable.manana} & nulo & Cubo 0 sin proporción\\
300 & desconocida & \texttt{laborable.manana} & nulo & Sin longitud\\
\bottomrule
\end{tabular}
\nota{Con las proporciones $r_{2,\texttt{all}}=0{,}769$, $r_{2,\texttt{laborable.punta\_manana}}=0{,}811$, $r_{4,\texttt{all}}=2{,}0$ y $r_{5,\texttt{all}}=0{,}1$.}
\end{table}

\section{Calibración}

\subsection{Tiempos por Tramo} Con las quince muestras siguientes (pista, franja, tramo, tiempo real), la calibración
da exactamente
\texttt{A>B}: $\{\texttt{n}:8,\ \texttt{all}:61,\ \texttt{laborable.punta\_manana}:69,\ \texttt{laborable.mediodia}:60\}$,
\texttt{C>D}: $\{\texttt{n}:3,\ \texttt{laborable.noche}:200\}$ y
\texttt{D>E}: $\{\texttt{n}:2,\ \texttt{laborable.manana}:60\}$; el tramo \texttt{B>C}, con una sola pista, no se
guarda.

\begin{center}
\small
\begin{tabular}{@{}llll@{\qquad}llll@{}}
\toprule
Pista & Franja & Tramo & $y$ & Pista & Franja & Tramo & $y$\\
\midrule
b1 & punta\_manana & A>B & 90 & m5 & mediodia & A>B & 61\\
b2 & punta\_manana & A>B & 100 & b1 & punta\_manana & B>C & 75\\
m0 & mediodia & A>B & 60 & c0 & noche & C>D & 200\\
m1 & mediodia & A>B & 58 & c1 & noche & C>D & 210\\
m2 & mediodia & A>B & 62 & c2 & noche & C>D & 190\\
m3 & mediodia & A>B & 64 & d1 & manana & D>E & 50\\
m4 & mediodia & A>B & 59 & d1 & manana & D>E & 70\\
 & & & & d2 & manana & D>E & 60\\
\bottomrule
\end{tabular}
\end{center}

\noindent El cálculo de \texttt{A>B} está en el ejemplo~\ref{ej:contraccion}. En \texttt{D>E}, la pista d1 tiene dos
muestras y cuenta la última (70); la mediana baja de $\{70,60\}$ es 60.

\subsection{Factor de Ritmo} Con una muestra por pista, de $H=600$~s y $Y=600\,\rho_j$:

\begin{center}
\small
\begin{tabular}{@{}lr@{}}
\toprule
Proporciones por pista $\rho_j$ & $k$\\
\midrule
0,80, 0,82, 0,84, 0,86, 0,88, 0,90, 0,92, 0,94 & 0,938\\
cuarenta veces 0,85 & 0,88\\
1,6, 1,7, \dots, 2,5 & 1,2\\
\bottomrule
\end{tabular}
\end{center}

\subsection{Paso Real} Con la pista $(28\,900,\,28\,880)$, $(28\,960,\,28\,940)$, $(29\,020,\,29\,030)$: para
$a=28\,950$, $P=28\,966{,}666666666668$; para $a=28\,800$, nulo.

\subsection{Aclarado} De 100 muestras de una pista se conservan las de índices 0, 3, 6, 10, 13, 16, 20, 23, 26, 30,
33, 36, 40, 43, 46, 50, 53, 56, 60, 63, 66, 70, 73, 76, 80, 83, 86, 90, 93 y 96.

\section{Corrección por Viaje}

\subsection{Consulta} La tabla de los vectores tiene el viaje \texttt{9999|S0|S4|28800|5} con $[20,40,60]$ en
\texttt{laborable} y $[-30,0,30]$ en \texttt{fin\_semana}; la línea \texttt{L2}, con 15~s en
\texttt{laborable.punta\_manana}; y por franja, 25~s en \texttt{laborable.punta\_manana} y $-10$~s en
\texttt{fin\_semana.noche}. Los tercios~\eqref{eq:tercio} de las posiciones 0 a 4 de un recorrido de cinco paradas son
0, 0, 1, 1 y 2; los de las posiciones 1 y 2 de tres, 1 y 2; la posición 9 de cinco da 2, y un recorrido sin paradas,
0. En un festivo, la capa 3 calcula con el domingo de esa semana: el lunes 12 de octubre y el sábado 17 dan el 18; el
domingo 1 de noviembre, el mismo día; y el martes 8 de diciembre, el 13.

\begin{center}
\small
\begin{tabular}{@{}lllrl@{}}
\toprule
Viaje & Fecha & Salida, posición & Resultado (s) & Por qué\\
\midrule
\texttt{9999|S0|S4|28800|5} & 2026-10-28 & 28\,800, 4 de 5 & 60 & Tercer tercio, laborable\\
\texttt{9999|S0|S4|28800|5} & 2026-10-28 & 28\,800, 2 de 5 & 40 & Segundo tercio\\
\texttt{9999|S0|S4|28800|5} & 2026-10-31 & 28\,800, 1 de 5 & $-30$ & Sábado: fin de semana\\
\texttt{L2|A|B|28800|3} & 2026-10-28 & 28\,800, 1 de 3 & 15 & Línea en su franja\\
\texttt{L2|A|B|72000|3} & 2026-10-31 & 72\,000, 1 de 3 & $-10$ & La línea no tiene esa franja\\
\texttt{L3|A|B|28800|3} & 2026-10-28 & 28\,800, 1 de 3 & 25 & Sin línea: la franja\\
\texttt{L3|A|B|50400|3} & 2026-10-28 & 50\,400, 1 de 3 & 0 & Nada para esa franja\\
\texttt{09999|S0|S4|28800|5} & 2026-10-28 & 28\,800, 4 de 5 & 60 & Se lee como la 9999\\
\bottomrule
\end{tabular}
\end{center}

\subsection{Calibración} Con los nueve desvíos siguientes (fecha, $D'$, viaje, posición y desvío en segundos), la
calibración~\eqref{eq:correccion-tablas} da para el viaje \texttt{L1|A|D|28800|4} en laborable $[127,150,420]$; para
\texttt{L1|A|D|32400|4}, $[-60,-60,-60]$ en laborable y $[30,30,40]$ en fin de semana; y para
\texttt{L2|X|Y|61200|6}, $[300,300,300]$. Las líneas quedan en \texttt{L1}: $\{$\texttt{laborable.punta\_manana}: 140,
\texttt{laborable.manana}: $-60$, \texttt{fin\_semana.manana}: 30$\}$ y \texttt{L2}:
$\{$\texttt{laborable.punta\_tarde}: 300$\}$, con los mismos valores por franja.

\begin{center}
\small
\begin{tabular}{@{}llllr@{}}
\toprule
Fecha & $D'$ & Viaje & Posición & Desvío\\
\midrule
2026-10-26 & 2026-10-26 & \texttt{L1|A|D|28800|4} & 1, 2 y 3 de 4 & 100, 160 y 220\\
2026-10-27 & 2026-10-27 & \texttt{L1|A|D|28800|4} & 1 y 3 de 4 & 140 y 2\,000\\
2026-10-26 & 2026-10-26 & \texttt{L1|A|D|32400|4} & 2 de 4 & $-60$\\
2026-10-31 & 2026-10-31 & \texttt{L1|A|D|32400|4} & 1 de 4 & 30\\
2026-10-12 & 2026-10-18 & \texttt{L1|A|D|32400|4} & 3 de 4 & 50\\
2026-10-26 & 2026-10-26 & \texttt{L2|X|Y|61200|6} & 5 de 6 & 300\\
\bottomrule
\end{tabular}
\end{center}

\noindent Para \texttt{L1} en \texttt{laborable.punta\_manana}, el desvío de 2\,000~s se recorta a 900. Ordenados,
$100,140,160,220,900$, el percentil 40 es el de la posición $\lfloor0{,}4\cdot4\rfloor=1$: $g=140$, y la línea, con el
peso 20 de la franja, también vale 140. El primer tercio del viaje (posición 1 de 4) tiene un desvío cada día, 100 y
140: $\redondeo\big((100+140+140)/3\big)=127$. El tercer tercio, con 220 y 900, da $\redondeo(1\,260/3)=420$. El
festivo del 12 de octubre cuenta como fin de semana: con el 31 de octubre, la franja \texttt{fin\_semana.manana} tiene
30 y 50, y su percentil 40 es 30; el primer tercio da $(30+30)/2=30$, el segundo toma el valor de la línea y el tercero
da $(50+30)/2=40$.

\section{El Emparejador Completo}

El viaje sintético tiene cinco paradas, \texttt{S0} a \texttt{S4}, en línea recta hacia el este sobre el paralelo
$43{,}26^\circ$ desde la longitud $-2{,}93^\circ$, a 0, 400, 1\,000, 2\,500 y 3\,200~m (longitudes
$-2{,}93$, $-2{,}925065928003549$, $-2{,}917664820008872$, $-2{,}8991620500221806$ y $-2{,}890527424028391$), con
llegada igual a salida en cada parada: $28\,800$, $28\,890$, $29\,010$, $29\,310$ y $29\,430$~s. El trazado, cuando se
usa, son las cinco paradas con sus distancias. La fecha es el miércoles 2026-10-28, la línea \texttt{9999} y el
vehículo \texttt{V1}. Los parámetros son:

\begin{itemize}
  \item tramos aprendidos: \texttt{S0>S1} $\{\texttt{n}:10,\ \texttt{all}:100,\ \texttt{laborable.punta\_manana}:110\}$
        y \texttt{S1>S2} $\{\texttt{n}:6,\ \texttt{all}:150,\ \texttt{laborable.punta\_manana}:140\}$;
  \item proporciones: cubo 2, $\{\texttt{all}:0{,}769,\ \texttt{laborable.punta\_manana}:0{,}811\}$; cubo 4,
        $\{\texttt{all}:0{,}64,\ \texttt{laborable.punta\_manana}:0{,}668\}$;
  \item factor global $k_{\texttt{laborable.punta\_manana}}=0{,}9$, sin factores de línea.
\end{itemize}

\noindent En el escenario de «horario muy lento», las horas de \texttt{S3} y \texttt{S4} pasan a $30\,810$ y
$30\,930$, y la tabla de tramos se sustituye por \texttt{S0>S1}, \texttt{S1>S2}, \texttt{S2>S3} y \texttt{S3>S4} con 110,
140, 240 y 100~s en \texttt{laborable.punta\_manana} (y \texttt{n}$=9$).

Las entradas del feed de cada escenario son:

\begin{center}
\small
\begin{tabularx}{\textwidth}{@{}l>{\raggedright\arraybackslash}X@{}}
\toprule
Escenario & Entrada\\
\midrule
1 y 2 & $\varphi=43{,}2601$, $\lambda=-2{,}9213653740062107$ (700~m), $\tgps=t_{\mathrm{rec}}=28\,950$, $n=3$\\
3 & $\varphi=43{,}260538986704994$ (60~m al norte), mismo $\lambda$ y horas, $n=3$\\
4 & sin posición, $t_{\mathrm{rec}}=28\,950$, $n=3$\\
5 & $\varphi=43{,}2601$, $\lambda=-2{,}895461496024842$ (2\,800~m), $\tgps=t_{\mathrm{rec}}=29\,400$, $n=3$\\
6, 7, 8, 10 y 11 & ninguna\\
9 & como el 1, con $\tgps=t_{\mathrm{rec}}=28\,300$\\
12 & $\varphi=43{,}2601$ (11~m al norte de \texttt{S0}), $\lambda=-2{,}93$, $\tgps=t_{\mathrm{rec}}=28\,200$, $n=2$, $\tnow=28\,210$\\
13 & $\varphi=43{,}26538986704995$ (600~m al norte), mismo $\lambda$ y horas, $n=2$\\
14 y 16 & como el 1, con $\tnow=29\,100$ y $29\,320$\\
15 & $\varphi=43{,}2601$, $\lambda=-2{,}917048061009316$ (1\,050~m, 50~m pasada \texttt{S2}), $\tgps=t_{\mathrm{rec}}=29\,020$, $n=3$, $\tnow=29\,040$\\
\bottomrule
\end{tabularx}
\end{center}

\noindent Los escenarios 10 y 11 usan la tabla de corrección de la sección anterior; el 11, además, con la fecha
marcada como festivo (un calendario de solo domingo que circula ese día por excepción). Los resultados están en la
Tabla~\ref{tab:escenarios}. Las posiciones calculadas son $\Tpos=28\,950$ con
$d_\perp=11{,}132000000369544$~m en los escenarios 1 (método «trazado») y 2 (método «tramo recto»); en el 12,
$\Tpos=\tgps=28\,800$ con la misma distancia y el método «cabecera»; en el 14 y el 16, la del 1; en el 15,
$\Tpos=29\,020$ con el método «trazado»; no hay posición en los demás.

\section{Del Feed a la Pantalla}

Las reglas de esta sección rodean al emparejador: cómo se lee el feed, a qué viaje del horario se asigna cada bus, qué
pasa con un bus que falta en una lectura o con un viaje que pasa de medianoche, qué se guarda de cada paso, qué retraso
se muestra y cómo se elige cada noche el percentil de la corrección por viaje. Sus vectores salen del mismo motor y
\texttt{comprobar.py} los comprueba igual que los anteriores.

\subsection{Lectura del Feed} Un fragmento SIRI con cuatro actividades, leído con fecha 2026-10-28. La referencia
\texttt{trp\_A0651\_737\_ABC1\_28500} da la línea \texttt{651} (sin el cero), el viaje \texttt{737} y la salida
$28\,500$; su \texttt{Delay} \texttt{PT2M30S} son 150~s; la hora de la posición, \texttt{07:02:30Z}, es
$\tgps=28\,950$ en Madrid, y la de la lectura, $t_{\mathrm{rec}}=28\,960$. La segunda actividad no trae posición, su
orden sale de \texttt{VisitNumber} y su \texttt{Delay} \texttt{-PT45S} son $-45$~s. La tercera se registró el día
anterior: entra en el feed, pero sin horas ni posición. La cuarta tiene una referencia de otro formato y no entra.

\subsection{Emparejado} Con la referencia exacta \texttt{trp\_A9999\_101\_X\_28800} del viaje \texttt{J} en la tabla de
alias:

\begin{center}
\small
\begin{tabularx}{\textwidth}{@{}>{\raggedright\arraybackslash}X l l@{}}
\toprule
Buses de la línea \texttt{9999}, viaje \texttt{101} & Viaje pedido & Bus elegido\\
\midrule
Uno con la referencia exacta y otro con otra referencia y la misma salida & \texttt{J} & el exacto\\
Sin referencia exacta, salida a 120~s & \texttt{J} & ese\\
Sin referencia exacta, salida a 180~s justos & \texttt{J} & ese\\
Sin referencia exacta, salida a 200~s & \texttt{J} & ninguno\\
Dos sin referencia exacta, a 150 y a 60~s & \texttt{J} & el de 60~s\\
Solo el de la referencia exacta de \texttt{J} & \texttt{K} & ninguno (no se reutiliza)\\
El de \texttt{J} y otro sin referencia exacta & \texttt{K} & el otro\\
\bottomrule
\end{tabularx}
\end{center}

\subsection{Buses que Faltan} Con $\tnow=29\,010$, de la lectura anterior se conservan, marcados como retenidos, el
bus visto hace 110~s y el que solo tiene hora de lectura de hace 180~s justos. No se conservan: el de hace 190~s ni el
de hace 210~s; el de un viaje que en la lectura nueva lleva otro bus; el que aparece en la lectura nueva en otro viaje;
el que no tiene hora; el de hora posterior a ahora; ni el que no tiene identificador de vehículo.

\subsection{Viajes que Pasan de Medianoche} A las 00:30 ($\tnow=1\,800$), el viaje que salió a las 23:30
($84\,600$) y el de las 24:10 del horario de ayer ($87\,000$) son de ayer; uno que sale a las 12:30 ($45\,000$, justo
12~h por delante) es de hoy, y uno a $45\,001$, de ayer. Para los de ayer, sus horas pasan a contarse desde la
medianoche de ayer: un dato de las 00:28:20 ($1\,700$) pasa a $88\,100$. Un bus sin hora de hoy no cuenta en ningún día.

\subsection{Registro de Pasos} De cada referencia se guarda «orden:hora», con la hora de la lectura o, sin ella, la de
la posición: \texttt{3:29000} y \texttt{2:28900}. No se guardan las entradas sin orden, sin referencia o sin hora. Al
leer, por viaje del horario queda el orden mayor de sus referencias (\texttt{3:29000} y \texttt{5:29100} dan orden 5 a
las $29\,100$); los valores nulos o mal formados se ignoran. Con ese paso, un viaje sin bus en el feed tiene pasadas las
paradas de orden menor:

\begin{center}
\small
\begin{tabular}{@{}lrlll@{}}
\toprule
Paso registrado & $\tnow$ & \texttt{S0}, \texttt{S1}, \texttt{S2} & \texttt{S3} & \texttt{S4}\\
\midrule
orden 4 a las $29\,300$ & $29\,400$ & pasadas, a las $29\,300$ & horario & capa 3: $29\,408$\\
orden 2 a las $29\,000$ & $28\,990$ & \texttt{S0} pasada a las $28\,990$ (ahora) & capa 3: $29\,310$ & capa 3: $29\,408$\\
\bottomrule
\end{tabular}
\end{center}

\noindent En el segundo caso, \texttt{S1} queda en el horario ($28\,890$) y \texttt{S2} con la capa 3 ($29\,110$).

\subsection{El Retraso Mostrado} Los minutos son los segundos entre 60 redondeados con mitades hacia fuera del cero:
29~s dan 0; 30, 1; 89, 1; 90, 2; $-30$, $-1$; $2\,700$, 45; y un retraso de más de $2\,700$~s no se muestra. Del bus
del escenario 1, con GPS, se muestra el retraso que llevará en su próxima parada, 10~s, aunque su retraso en la
posición sea 0; sin GPS (escenario 4), el de la parada registrada, 10~s; y con un dato de más de 10~min y sin hora de
lectura, 0.

\subsection{Elección del Percentil} Con un viaje de la línea \texttt{L1} y cinco desvíos en cada uno de cuatro días,
se aprende con los tres primeros y se prueba en el último:

\begin{center}
\small
\begin{tabular}{@{}lrrrrr@{}}
\toprule
Percentil & 0,30 & 0,35 & 0,40 & 0,45 & 0,50\\
\midrule
Error medio (s) & 144,4 & 144,4 & 147,4 & 149,2 & 151,2\\
A $\pm2$~min (\%) & 20 & 20 & 40 & 60 & 60\\
\bottomrule
\end{tabular}
\end{center}

\noindent Se elige 0,45: el que más acierta a $\pm2$~min y, en el empate con 0,50, el de menos error medio, aunque 0,30
tenga el menor. Sin corregir, el error medio es 126~s y el acierto, 60\,\%: la tabla nueva no bajaría el error medio y
no se escribe.

\chapter{Implementación de Referencia}
\label{ap:referencia}

La implementación de referencia, en Python 3 y sin dependencias externas, se escribió a partir de la especificación de
esta memoria y no del código PHP. Reproduce todos los vectores del Apéndice~\ref{ap:vectores}. Su propósito es servir
de punto de partida para portar el método a otro lenguaje: es corta, sigue el orden de los capítulos y usa los mismos
nombres que el texto.

\section{Cómo Comprobarla}

Los ficheros van adjuntos a este PDF: \texttt{referencia.py}, \texttt{comprobar.py}, \texttt{vectores.json} y
\texttt{desvio-horario.py}. La mayoría de los lectores de PDF los muestran en un panel de adjuntos, desde el que se
pueden guardar; también se pueden extraer con \texttt{pdfdetach -saveall} (de Poppler). Con los cuatro en una carpeta:

\begin{lstlisting}[language=bash,numbers=none]
python3 comprobar.py
# 407 de 407 comprobaciones coinciden con el código de producción
\end{lstlisting}

\noindent La carpeta \texttt{referencia} del repositorio público trae además \texttt{evaluar.py}, el evaluador
público: repite la evaluación de la sección «\nameref{sec:varios-dias}» sobre capturas del feed y escribe los mismos
histogramas de 15~s que la de producción, idénticos clase a clase en las 92\,686 predicciones de 40~minutos de feed
del 7 de noviembre de 2026. Los vectores los generó el autor con el código de producción, que es privado, con el
programa \texttt{vectores.php}, que necesita ese código y por eso no se publica; \texttt{vectores.json} trae además un caso que
\texttt{comprobar.py} no usa (\texttt{k\_un\_bus}, el factor sin contraer de una sola pista). No hace falta ese código para
usarlos: una implementación nueva, en cualquier lenguaje, está bien hecha cuando reproduce todos los vectores. El
programa \texttt{desvio-horario.py} recalcula las cifras de la sección «\nameref{sec:desvio-horario}»; necesita la base del
horario y las muestras de tramo, que se generan con el propio método (capítulos «\nameref{cap:datos}»
y «\nameref{cap:aprendizaje}»).

\embedfile[desc={Implementación de referencia},filespec=referencia.py,ucfilespec=referencia.py]{../referencia/referencia.py}
\embedfile[desc={Comprobación contra los vectores de prueba},filespec=comprobar.py,ucfilespec=comprobar.py]{../referencia/comprobar.py}
\embedfile[desc={Vectores de prueba generados con el código de producción},filespec=vectores.json,ucfilespec=vectores.json]{../referencia/vectores.json}
\embedfile[desc={Cifras de la sección 1.1},filespec=desvio-horario.py,ucfilespec=desvio-horario.py]{../referencia/desvio-horario.py}

\section{Lo que Hay que Respetar al Portar}

\begin{enumerate}
  \item \textbf{Redondeo.} $\redondeo$ redondea las mitades alejándose del cero. En Python, \texttt{round(2.5)} da 2;
        hay que usar la función \texttt{php\_round} de la referencia. En JavaScript, \texttt{Math.round(-2.5)} da $-2$
        y tampoco sirve para negativos. El redondeo a tres decimales es $\redondeo(1000x)/1000$.
  \item \textbf{Mediana baja.} $x_{(\lfloor p(n-1)\rfloor)}$ sobre los valores ordenados, sin promediar los dos
        centrales. Las funciones de mediana de las bibliotecas estadísticas suelen interpolar.
  \item \textbf{División entera.} $t_{\mathrm{mid}}=\lfloor(d_1+a_\dest)/2\rfloor$ y la hora
        $\lfloor t/3600\rfloor\bmod24$ se calculan con división entera.
  \item \textbf{Truncado.} El campo \texttt{at} de las muestras de tramo es la parte entera del paso real, no su
        redondeo.
  \item \textbf{CRC32 sin signo.} El pliegue de validación usa el CRC32 de IEEE de los bytes UTF-8 de la pista, como
        entero sin signo; en lenguajes con enteros de 32 bits con signo hay que convertirlo antes del módulo.
  \item \textbf{Lista u objeto.} \texttt{lengthRatios} puede llegar como lista o como objeto
        (Apéndice~\ref{ap:formatos}).
  \item \textbf{Orden de evaluación.} La posición, con su contraste, se evalúa antes que cualquier uso del registro
        (sección «\nameref{sec:desconfianza}»).
  \item \textbf{Bucles por horas.} El tiempo esperado recorre los tramos comparando horas, no índices
        (sección «\nameref{sec:tiempo-esperado}»).
  \item \textbf{Festivos.} La capa 3 de un festivo entre semana se calcula con el domingo de esa semana, también la
        franja de la salida con la que se busca la corrección de la línea.
  \item \textbf{Dos redondeos.} En la capa 3, el tiempo esperado y la corrección por viaje se redondean por separado.
  \item \textbf{El día del viaje.} La referencia recibe la entrada, $\tnow$ y la fecha ya en las horas del día del viaje.
        Pasada la medianoche, una entrada de un viaje de ayer~\eqref{eq:dia-viaje} se calcula con $\tgps$,
        $t_{\mathrm{rec}}$ y $\tnow$ más $86\,400$~s y con la fecha de ayer; una entrada sin ninguna hora de hoy no se usa
        (sección «\nameref{sec:mas-de-24h}»). El registro de pasos y la memoria del feed, que dependen de un almacén
        externo, no están en la referencia.
\end{enumerate}

\section{Código}

\lstinputlisting[language=Python,caption={\texttt{referencia.py}},label=lst:referencia]{../referencia/referencia.py}

\lstinputlisting[language=Python,caption={\texttt{comprobar.py}},label=lst:comprobar]{../referencia/comprobar.py}

\chapter{Glosario}
\label{ap:glosario}

\begin{description}
  \item[Alias de viaje.] Correspondencia entre un \texttt{trip\_id} original del GTFS y el viaje fusionado que lo
        contiene (sección «\nameref{sec:alias}»).
  \item[Cabecera.] Primera parada de un viaje; su hora de salida es $d_1$.
  \item[Capa.] Cada una de las tres fuentes de la predicción: GPS, parada registrada y horario corregido
        (capítulo «\nameref{cap:prediccion}»).
  \item[Captura.] Conjunto de versiones del feed descargadas durante un bloque de captura (hasta una hora)
        (sección «\nameref{sec:captura}»).
  \item[Contracción.] Media ponderada entre una estimación propia y un valor común, con un peso que crece con los datos
        propios (secciones «\nameref{sec:marco-contraccion}» y «\nameref{sec:contraccion-tramo}»).
  \item[Corrección por viaje.] Lo que suele apartarse cada viaje del horario corregido, por tipo de día y tercio del
        recorrido (sección «\nameref{sec:correccion-viaje}»).
  \item[Cubo de longitud.] Intervalo de longitudes de tramo para el respaldo por longitud
        (ecuación~\eqref{eq:cubo}).
  \item[Día del viaje.] Día de servicio al que pertenece un viaje del horario. Sus horas se cuentan desde su medianoche
        y pueden pasar de las 24:00 (sección «\nameref{sec:mas-de-24h}»).
  \item[Desfase.] Diferencia $\Delta(s)=R(s)-T(s)$ entre la hora real y la programada en un punto del recorrido.
  \item[Feed.] Fuente de datos en tiempo real publicada por el operador; aquí, SIRI Vehicle Monitoring.
  \item[Festivo entre semana.] Día en que el horario hace circular por excepción los calendarios de solo domingo; en la
        capa 3 cuenta como fin de semana.
  \item[Franja.] Combinación de tipo de día y periodo del día (definición~\ref{def:franja}).
  \item[GTFS.] \emph{General Transit Feed Specification}, formato abierto de horarios de transporte público.
  \item[Horario corregido.] Horario recalculado desde la cabecera con los tiempos aprendidos (capa 3).
  \item[Horizonte.] Tiempo que falta para la llegada en el momento de la predicción.
  \item[Interpolación del horario.] Cálculo de las horas de las paradas intermedias repartiendo el tiempo entre puntos
        de control a velocidad constante (\texttt{timepoint=0}).
  \item[Mediana baja.] Elemento central de una lista ordenada; con un número par de elementos, el menor de los dos
        centrales.
  \item[Muestra de ritmo.] Par (tiempo programado, tiempo real) que faltaba para llegar a una parada desde una
        posición (sección «\nameref{sec:muestras-ritmo}»).
  \item[Muestra de tramo.] Tiempo real que tardó un autobús en ir de una parada a la siguiente
        (sección «\nameref{sec:muestras-tramo}»).
  \item[Pasillo.] Distancia máxima al recorrido para aceptar una posición GPS.
  \item[Paso real.] Instante, reconstruido, en que un autobús pasó por una parada (definición~\ref{def:paso-real}).
  \item[Pista.] Recorrido de un vehículo en un viaje en una fecha; unidad de independencia de las muestras.
  \item[Pliegue.] Mitad de las pistas en la validación cruzada.
  \item[Proporción por longitud.] Mediana de tiempo real entre tiempo programado de los tramos de un cubo de longitud.
  \item[Registro.] Último dato del feed sobre la próxima parada de un autobús y su hora, usado por la capa 2.
  \item[Registro de pasos.] Almacén que guarda el último orden de parada visto de cada viaje, por día del viaje
        (sección «\nameref{sec:registro-pasos}»).
  \item[SIRI-VM.] \emph{Service Interface for Real Time Information, Vehicle Monitoring}: norma europea para publicar la
        actividad de los vehículos.
  \item[Tramo.] Par de paradas consecutivas de un viaje.
  \item[Trazado.] Polilínea geográfica que sigue el recorrido de un viaje (\texttt{shapes.txt}).
  \item[$\Tpos$.] Hora a la que el horario pone al autobús en el punto del recorrido donde está.
  \item[Ventana.] Conjunto de tramos alrededor de la próxima parada donde se busca la proyección del GPS.
\end{description}


\end{document}
